
%% Template by Michal Forisek


\documentclass[a4paper]{report}
\usepackage{ulem}
\usepackage{slovak}
\usepackage[utf8]{inputenc}
\usepackage{a4wide}
\usepackage{tabularx}
\usepackage{amsfonts}
\usepackage{amssymb}
\usepackage{amsmath}
\usepackage{epsfig}
\usepackage[usenames,dvipsnames]{color}
\usepackage{mathrsfs}
\usepackage{verbatim}
\usepackage{hyperref}
\usepackage{ifthen}

% vim: set fdm=marker:
%% Original by Michal Forisek

\ifx \RestyleAlgo \undefined
    \def\RestyleAlgo#1{\restylealgo{#1}}
\fi

\RestyleAlgo{boxed} %% nastav styl algorithm2e

%% zakladne definicie
\newcommand{\quoteme}[1]{\clqq#1\crqq}
\def\todo#1{[{\color{red} TODO:} {\bf  #1}]}
\def\fixme#1{[{\color{red} FIXME:} {\bf  #1}]}
\def\verify#1{\todo{verify: #1}}

\def\xor{\oplus}
\def\concat{\|}
%\def\inr{\in_{R}}
\def\toa #1 {\overset{#1}{\rightarrow}}
\def\inr{\overset{\$}{\leftarrow}}
\def\assign{\leftarrow}
\def\send{\rightarrow}
\def\isomorph{\cong}
\def\nsd{NSD}
\def\union{\cup}
\newcommand{\unit}[1]{\ensuremath{\, \mathrm{#1}}}
\DeclareMathOperator{\dlog}{dlog}

\def\compactlist{
  \setlength{\itemsep}{1pt}
  \setlength{\parskip}{0pt}
  \setlength{\parsep}{0pt}
}
\def\mod{\,{\rm mod}\,}

%%% original od Misofa:
%% {{{

\catcode`\@=11

\def\R{{\cal R}}
\def\cent{{c\kern-0.3em|\kern0.1em}}
\def\N{{N}} % FIXME FIXME 

\let\eps=\varepsilon

\def\relupdown#1#2#3{\mathrel{\mathop{#1}\limits^{#2}_{#3}} }

\let\then=\Rightarrow
\let\neht=\Leftarrow

\def\krok#1{\relupdown{\Longrightarrow}{}{#1}}
\def\thenrm{\relupdown{\Longrightarrow}{}{rm}}

\def\bicik{\upharpoonright}
\def\B{{\mathbf B}}
\def\kodTS#1{{\tt <}#1{\tt >}}

\newtheorem{definicia}{Definícia}[section]
\newtheorem{HLPpoznamka}{Poznámka}[section]
\newtheorem{HLPpriklad}{Príklad}[section]
\newtheorem{HLPcvicenie}[HLPpriklad]{Cvičenie}
\newtheorem{zadanie}{Úloha}[section]
\newenvironment{poznamka}{\begin{HLPpoznamka}\rm}{\end{HLPpoznamka}}
\newenvironment{priklad}{\begin{HLPpriklad}\rm}{\end{HLPpriklad}}
\newenvironment{cvicenie}{\begin{HLPcvicenie}\rm}{\end{HLPcvicenie}}
\newtheorem{veta}{Veta}[section]
\newtheorem{lema}[veta]{Lema}
\newtheorem{dosledok}[veta]{Dôsledok}
\newtheorem{teza}[veta]{Téza}
% \newtheorem{dokaz}{Dôkaz}[section]

\long\def\odsadene#1{
\leftskip=\parindent
\parindent=0pt
\vskip-5pt

\parskip=5pt
#1
\parskip=0pt

\parindent=\leftskip
\leftskip=0pt

} % end \odsadene




%%%%%%%%%%% PROSTREDIE PRE PISANIE KOMENTAROV

%\newenvironment{komentar}{%
%\vskip\baselineskip
%\tabularx{0.95\textwidth}{|X|}
%\sl
%}
%{\endtabularx
%\vskip\baselineskip
%}

\newenvironment{komentar}{%
\vskip\baselineskip\noindent
\tabularx{\textwidth}{>{\hsize=.2\hsize}X>{\hsize=1.8\hsize}X}
\sl ~ & \sl
}
{\endtabularx
\vskip\baselineskip
}

%\newenvironment{komentar}{%
%\vskip\baselineskip
%\trivlist\vspace{-4pt}\raggedleft\item\relax\tabularx{0.9\textwidth}{X}\sl}
%{\endtabularx\vspace{-4pt}\endtrivlist
%\vskip\baselineskip
%}

\newenvironment{dokaz}{\trivlist
  \item[\hskip \labelsep{\bfseries Dôkaz.}]}{\endtrivlist}
  
%\newenvironment{dokaz}{%
%\vskip\baselineskip\noindent
%\tabularx{\textwidth}{||X||}
%\sl
%}
%{\endtabularx
%\vskip\baselineskip
%}

%%%%%%%%%%% PROSTREDIE PRE MOJE ITEMIZE 

\newenvironment{myitemize}{%
\begin{itemize}
\itemsep-3pt
}
{\end{itemize}
}

%%%%%%%%%%% MATICKE MAKRA

\font\tenrm=csr10

\def\eps{\varepsilon}
% \def\R{{\mathbb R}}
\def\lvec#1{\overrightarrow{#1}}
\def\uhol{{\measuredangle}}
\def\then{\Rightarrow}
% \def\lg{{\rm lg}}
\def\lg{\log_2}
%\def\div{\mathbin{\rm div}}
\def\div{{\rm div}}

%%%%%%%%%%% PDF

\newif\ifpdf
\ifx\pdfoutput\undefined
  \pdffalse
\else
  \pdfoutput=1 \pdftrue
\fi

%%%%%%%%%%% OBRAZKY 

\newcommand{\myincludegraphics}[2][]{\includegraphics[#1]{images/#2}}

%%%%%%%%%%% SLOVNICEK

\openout2=\jobname.slo

\newcommand{\definuj}[3][]{%
\def\tmpvoid{}\def\tmpfirst{#1}%
\ifx\tmpvoid\tmpfirst%
  {\sl #2}\label{definicia:#2}\write2{#2 & #3 & \pageref{definicia:#2} \cr}%
\else%
  {\sl #2}\label{definicia:#2}\write2{#1 & #3 & \pageref{definicia:#2} \cr}%
\fi}

\newcommand{\definujsilent}[2]{%
\label{definicia:#1}\write2{#1 & #2 & \pageref{definicia:#1} \cr}%
}

\newcommand\myglossary{
  \immediate\closeout2 
  %\if@twocolumn\@restonecoltrue\onecolumn\else\@restonecolfalse\fi
  \chapter{Slovníček pojmov}
  \begin{tabular}{|l|l|r|}
  \hline
  {\bfseries slovenský pojem} & {\bfseries anglický preklad} & {\bfseries str.} \\ 
  \hline
  \InputIfFileExists{\jobname.srs}{}{~ & ~ & ~ \\}
  \hline
  \end{tabular}
  %\if@restonecol\twocolumn\fi
}

%%%%%%%%%%% UVODZOVKY

\catcode`\"=13
\def "{\begingroup\clqq\def "{\endgroup\crqq}}
\def\dospecials{\do\ \do\\\do\{\do\}\do\$\do\&%
  \do\#\do\^\do\^^K\do\_\do\^^A\do\%\do\~\do\"}

%%%%%%%%%%% DANGER BENDS 

\font\manual=manfnt % font used for the METAFONT logo, etc.
\def\dbend{{\manual\char127}} % dangerous bend sign

\newlength{\bendwidth}   \settowidth{\bendwidth}{\dbend}    \newlength{\hangwidth}

\def\hangone{%
  \hangwidth=\bendwidth%
  \advance\hangwidth 5pt%
  \hangindent\hangwidth%
}
\def\hangtwo{%
  \hangwidth=\bendwidth%
  \multiply\hangwidth 2%
  \advance\hangwidth 6pt% 
  \hangindent\hangwidth%
}

\def\medbreak{\par\ifdim\lastskip<\medskipamount \removelastskip\penalty-100\medskip\fi}
\let\endgraf=\par

\def\d@nger{\medbreak\begingroup\clubpenalty=10000
%\def\d@nger{\begingroup\clubpenalty=10000
%  \def\par{\endgraf\endgroup\medbreak} \noindent\hangone\hangafter=-2
  \def\par{\endgraf\endgroup} \noindent\hangone\hangafter=-2
  \hbox to0pt{\hskip-\hangindent\dbend\hfill}}
\outer\def\danger{\d@nger}

\def\dd@nger{\medbreak\begingroup\clubpenalty=10000
%  \def\par{\endgraf\endgroup\medbreak} \noindent\hangtwo\hangafter=-2
  \def\par{\endgraf\endgroup} \noindent\hangtwo\hangafter=-2
  \hbox to0pt{\hskip-\hangindent\dbend\kern1pt\dbend\hfill}}
\outer\def\ddanger{\dd@nger}

\def\enddanger{\endgraf\endgroup} % omits the \medbreak
\def\enddangerhop{\endgraf\endgroup\medbreak}




\def\@nakedcite#1#2{{#1\if@tempswa , #2\fi}}
\DeclareRobustCommand\nakedcite{%
  \@ifnextchar [{\@tempswatrue\@nakedcitex}{\@tempswafalse\@nakedcitex[]}}
\def\@nakedcitex[#1]#2{%
  \let\@citea\@empty
  \@nakedcite{\@for\@citeb:=#2\do
    {\@citea\def\@citea{,\penalty\@m\ }%
     \edef\@citeb{\expandafter\@firstofone\@citeb\@empty}%
     \if@filesw\immediate\write\@auxout{\string\citation{\@citeb}}\fi
     \@ifundefined{b@\@citeb}{\mbox{\reset@font\bfseries ?}%
       \G@refundefinedtrue
       \@latex@warning
         {Citation `\@citeb' on page \thepage \space undefined}}%
       {\hbox{\csname b@\@citeb\endcsname}} }}{#1}}

\long\def\FIXME#1{
  \begin{center}
  \begin{minipage}{0.8\textwidth}
  {\bf FIXME:~}\sl #1
  \end{minipage}
  \end{center}
}


\catcode`\@=12
%% }}}


\begin{document}

\thispagestyle{empty}
\begin{minipage}{0.25\textwidth}
\includegraphics[width=0.9\textwidth]{img/komlogo-new}
\end{minipage}
\begin{minipage}{0.69\textwidth}
\begin{center}
\sc Katedra Informatiky \\
Fakulta Matematiky, Fyziky a Informatiky \\
Univerzita Komenského, Bratislava
\end{center}
\end{minipage}

\vfill
\begin{center}
\begin{minipage}{0.8\textwidth}
\hrule
\bigskip\bigskip
\centerline{\LARGE\sc Matematická Logika}
\smallskip
\centerline{(spísané poznámky, draft)}
\bigskip
\bigskip
\centerline{\large\sc Milan Plžík, Pavol Struhár, Ivan Kováč, Peter Perešíni}
\bigskip\bigskip
\hrule
\end{minipage}
\end{center}
\vfill
{~}
\hfill verzia zo dňa {\bf\today} 
\eject % EOP i

\section*{Úvod}

Tieto poznámky obsahujú študijné materiály
k predmetu \emph{Matematická logika}
na Fakulte matematiky, fyziky a informatiky UK.

Základ poznámok bol spísaný podľa prednášky doc. Eduarda Tomana v roku 2009.
Poznámky ale nie sú oficiálny študijný materiál, preto autori neručia
žiadnym spôsobom za ich aktuálnosť či vhodnosť. Navyše, obsah prednášky sa
počas rôznych rokov môže meniť a preto je silne odporúčané dopísať si
prípadné rozdiely medzi poznámkami a prednáškou.

Aby sme umožnoli jednoduchšie spravovanie a udržali poznámky dlhšie
aktuálne, rozhodli sme sa verejne publikovať zdrojové kódy na stránke
\url{http://code.google.com/p/matlogika}. Ak máte akékoľvek pripomienky,
návrhy, opravy, môžete nám ich prostredníctvom tejto stránky oznámiť.

Za autorov, PPershing.

\tableofcontents

\chapter{Dokončenie úvodu do matematickej logiky}
\section{Prerekvizity a označenia}

V tejto časti si zhrnieme najdôležitejšie označenia, pojmy, vety z
prednášky Úvod do matematickej logiky.

\noindent Označenia:
\begin{itemize}
    \item $\provable A$ -- formula $A$ je dokázateľná
    \item $M \models A$ -- $M$ je modelom $A$
    \item $\implies, \leftrightarrow$ -- implikácia, ekvivalencia vo
    výrokovej/predikátovej logike
    \item $\Rightarrow, \Leftrightarrow$ -- implikácia a ekvivalencia v
    našom jazyku
\end{itemize}

\noindent Axiómy:
\begin{itemize}
    \item[A1:] $(A \implies (B \implies A))$
    \item[A2:] $((A \implies (B \implies C)) \implies 
                ((A \implies B) \implies (A \implies C)))$
    \item[A3:] $((\neg B \implies \neg A) \implies (A \implies B))$
    \vskip 0.5cm
    \item[A4:] $(\forall x) A \implies A_x[t]$
    \item[A5:] $(\forall x) (A \implies B) \implies (A \implies
    (\forall x) B)$.
\end{itemize}

\noindent Pravidlá:
\begin{itemize}
    \item Modus ponens: $\displaystyle \frac{A,A\implies B}{B}$
    \item Pravidlo zovšeobecnenia:
            $\displaystyle \frac{A,x}{(\forall x)A}$
    \item Jednoduchý sylogizmus:
            $\displaystyle \frac{A, A\implies B, B\implies C}{C}$, resp.
            $\provable (A \implies B) \implies ((B \implies C) 
                \implies (A \implies C))$.
\end{itemize}

\begin{veta}[O dedukcii]
    \begin{equation*}
        T \provable A \implies B \Leftrightarrow T,A \provable B
    \end{equation*}
\end{veta}

\begin{veta}[Postove vety a iné užitočné tvrdenia:]
    \noindent
    \begin{itemize}
        \item $\provable A \implies A$
        \item $\provable A \implies \neg \neg A$
        \item $\provable \neg \neg A \implies A$
        \item $\provable (A \implies B) \implies (\neg B \implies \neg A)$
    \end{itemize}
\end{veta}

\begin{lema}[O neutrálnej formule]
   Nech $T,A \provable B$ a tiež $T, \neg A \provable B$. Potom
   $T \provable B$.
\end{lema}

\begin{lema}[O distribúcii kvantifikátorov]
    Ak je $\provable A \implies B$, potom
    $\provable (\forall x)A \implies (\forall x) B$ a
    $\provable (\exists x)A \implies (\exists x)B$.
\end{lema}

\begin{veta}[O ekvivalencii]
    Nech formula $A'$ vznikne z formuly A nahradením všetkých výskytov 
    podformúl $B_1,\dots ,B_n$ v uvedenom poradí formulami 
    $B'_1, \dots, B'_n$. Ak platí $\provable B_i \leftrightarrow B'_i$ pre
    $i \in \{1,\dots,n\}$, potom $\provable A \leftrightarrow A'$.
\end{veta}

\todo{vety o variantoch, zavedeni kvantifikatorov, ...}
\begin{lema}[Duálny tvar axiómy špecifikácie]
    $\provable A_x[t] \implies (\exists x) A$
\end{lema}

\begin{lema}[Pravidlo zavedenia existenčného kvantifikátora]
    Ak $\provable A \implies B$ a $x$ nie je voľná v $B$, potom
    $\provable (\exists x) A \implies B$.
\end{lema}

% vim: set fdm=marker: sw=2: sts=2: ts=2:

\section{Prenexné tvary formúl}

Ako sme mali vo výrokovej logike isté normálne tvary - konjunktívnu a
disjunktívnu normálnu formu, budeme mať aj v predikátovej logike isté
špeciálne tvary. Zaujímavé sú najmä prexenxná forma a ešte Skolemov
normálny tvar, čo je špeciálny prípad prexenej formy.
V prípade prenexného tvaru ide o preskupenie kvantifikátorov na
začiatok formuly - tvar vzniká aplikovaním kvantifikátorov na otvorenú
formulu.


\begin{definicia}
    Formula $A$ je v prenexnom tvare, ak $A$ je v nasledujúcom tvare:
    \begin{equation*}
     (Q_1 x_1) (Q_2 x_2) (Q_3 x_3) \dots (Q_n x_n) B
    \end{equation*}
    kde $x_1, \dots , x_n$ sú navzájom rôzne premenné,
    $Q_i \in\{\forall, \exists\}$ sú
    kvantifikátory a  formula $B$ je bez kvantifikátorov.
    Formulu $B$ nazveme otvoreným jadrom formuly $A$, 
    formulu $(Q_1 x_1) (Q_2 x_2) (Q_3 x_3) \dots (Q_n x_n)$ nazveme
    prefixom formuly $A$.
\end{definicia}

\begin{poznamka}
    \noindent
    \begin{itemize}
     \item $x_1, \dots, x_n$ sú navzájom rôzne pre vylúčenie viacnásobných
        kvantifikácií.
     \item ak $n=0$, tak $A$ je otvorená a nemá prefix
     \item $B$ je najväčšia otvorená podformula formuly $A$.
    \end{itemize}
\end{poznamka}

\begin{priklad}
    Formula elementárnej aritmetiky:
    \begin{equation*}
        (\forall x) (\forall y) (\exists z) (x+y=z)
    \end{equation*}
\end{priklad}

\begin{veta}
 Nech $A$ je ľubovoľná formula predikátovej logiky. Potom existuje
 formula $A'$ v prenexnom tvare taká, že
 $\provable A \leftrightarrow A'$.
 \label{veta:prenex}
\end{veta}

Pri prevádzaní formuly na prenexný tvar budeme využívať nasledujúce
\emph{prenexné} operácie, každá z nich nahrádza podformulu
jej ekvivalentom.
\begin{itemize}
    \item[a)] $B$ nahraď variantom formuly $B$ (premenovanie viazaných
    premenných)
    \item[b)] $\neg(Q x) B$ nahraď $(\overline{Q} x) \neg B$
    \item[c)] ak $x$ nie je voľná v $B$, tak podformulu $B\implies (Qx)C$
            nahraď podformulou $(Qx) (B\implies C)$
    \item[d)] ak $x$ nie je voľná v $C$, tak $((Qx)) B \implies C$
        nahraď $(\overline{Q} x) (B \implies C)$
    \item[e)] ak $x$ nie je voľná v $B$, $\square \in \{\land,\lor\}$.
     Potom $B \squareop ((Qx)C)$ resp. $((Qx)B)\squareop C$ nahraď
     $(Qx)(B \squareop C)$
\end{itemize}
\begin{poznamka}
    Asi stojí za zmienku upozorniť, že v časti c) $x$ nie je voľná vo
    formule $C$ narozdiel od častí b), d), kde $x$ nie je voľná v $B$.
\end{poznamka}

\begin{lema}
    Prenexnými operáciami dostaneme ekvivalentné formuly
\end{lema}
\begin{dokaz}
  %%% {{{ 
  \noindent
  \begin{itemize}
    \item[a)] Veta o variantoch
    \item[b)] Platí
        \begin{eqnarray*}
            \provable \neg(\forall x) B \leftrightarrow
                \neg (\forall x) \highlightb{\neg\neg} B \\
            \provable \highlighta{\neg (\forall x) \neg} \neg B
                \leftrightarrow \highlightb{(\exists x)} \neg B
        \end{eqnarray*}
        a taktiež
        \begin{eqnarray*}
            \provable \neg(\exists x) B \leftrightarrow
                \neg (\exists x) \highlightb{\neg\neg} B \\
            \provable \highlighta{\neg (\exists x) \neg} \neg B
                \leftrightarrow \highlightb{(\forall x)} \neg B                
        \end{eqnarray*}
        z čoho vieme priamo odvodiť dané tvrdenie.
    \item[c)] Nech $Q=\forall$. Chceme ukázať
        $\provable (\forall x) (B \implies C) \leftrightarrow
            (B \implies (\forall x) C)$ kde $x$ nie je voľná v $B$.
        \begin{itemize}
        \item[$\Rightarrow$]
            5. Axióma predikátovej logiky


        \item[$\Leftarrow$]
            \begin{itemize}
            \item[1] $\provable (\forall x) C \implies C$ --
                axióma špecifikácie

            \item[2] $\provable \highlighta{
                 \underbrace{(B \implies (\forall x) C)}_X
                \implies 
                 (\underbrace{((\forall x) C \implies C)}_Y
                    \implies \underbrace{(B \implies C)}_Z)}$ --
                    Jednoduchý sylogizmus

            \item[*] $\provable X \implies (Y \implies Z)$.

            \item[*] $\provable (X \implies (Y \implies Z) \implies
                (Y \implies (X \implies Z))$ -- pravidlo zámeny
                predpokladov

            \item[3] $\provable \highlighta{
                [(B \implies (\forall x) C)
                \implies 
                 (((\forall x) C \implies C)
                    \implies (B \implies C))]
                }
                \implies \highlightb{
                   [((\forall x) C \implies C) \implies
                    ((B \implies (\forall x) C) \implies (B \implies
                    C))]}$.

            \item[4] $\provable
                  \highlightb{
                   ((\forall x) C \implies C) \implies
                    \highlighto{
                    [(B \implies (\forall x) C) \implies (B \implies
                    C)]}}$ - MP 2,3

            \item[5] $\provable \highlighto{
                (B \implies (\forall x)) \implies (B
                    \implies C)}$ - MP 1,4

            \item[6] $\provable (B \implies (\forall x) C) \implies
            \highlighta{(\forall x)} (B \implies C)$ -- 
                pravidlo zavedenia veľkého kvantifikátora
            \end{itemize}
        \end{itemize}
     Druhou možnosťou je $Q=\exists$. Našim cieľom
     $\provable (\exists x) (B \implies C) \leftrightarrow 
        (B \implies (\exists x) C)$ za predpokladu že $x$ nie je voľná v $B$.
        \begin{itemize}
        \item[$\Rightarrow$]
            \begin{itemize}
            \item[1] $\provable C \implies (\exists x) C$ --
                duálna verzia axiómy špecifikácie

            \item[2] $\provable (B \implies C) \implies
                 ((C \implies (\exists c) C) \implies (B \implies
                 (\exists x) C ))$ -- jednoduchý sylogizmus (JS)

            \item[3] $\provable
                \highlighta{
                [(B \implies C) \implies (( C \implies (\exists x) C)
                \implies (B \implies (\exists x) C))]}
                \implies
                \highlightb{
                  [(C \implies (\exists x) C) \implies (( B \implies
                  C) \implies (B \implies (\exists x) C))]}$ -- 
                  pravidlo zámeny predpokladov

            \item[4] $\provable \highlightb{
                  ((C \implies (\exists x) C) \implies 
                  \highlighto{[( B \implies
                  C) \implies (B \implies (\exists x) C))]}}$ -- MP 2,3

            \item[5] $\provable
                  \highlighto{(( B \implies
                  C) \implies (B \implies (\exists x) C)))}$ -- MP 1,4

            \item[6] $\provable \highlighta{(\exists x)}
                ( B \implies C) \implies (B \implies (\exists x) C)$
                -- pravidlo zavedenie existenčného kvantifikátora
            \end{itemize}


        \vskip 5mm
        \item[$\Leftarrow$]
            \begin{itemize}
                \item[1] $\provable C \implies (B \implies C)$ -- A1

                \item[2] $\provable (\exists x) C \implies (B \implies
                C)$ -- pravidlo zavedenia existenčného
                kvantifikátora

                \item[3] $\provable \neg B \implies (B \implies C)$ -- postova
                teoréma

                \item[4] $\provable (B \implies C) \implies 
                    (\exists x)(B \implies C)$ -- 
                    duálny tvar axiómy špecifikácie

                \item[5] $\provable \highlighta{\neg B \implies 
                    (\exists x) (B \implies C)}$ -- JS 3,4

                \item[*] $\provable
                    \highlighta{
                    [\neg \underbrace{B}_X \implies 
                        \underbrace{(\exists x) ( B \implies C)}_Z]
                    }
                    \implies
                    \highlightb{
                    [(\underbrace{(\exists x) C}_Y 
                        \implies 
                      \underbrace{(\exists x) (B \implies C)}_Z)
                    \implies
                    ((\underbrace{B}_X \implies 
                        \underbrace{(\exists x) C}_Y)
                      \implies 
                        \underbrace{(\exists x) (B
                    \implies C)}_Z)]}$ -- dokážeme neskôr

                \item[6] $\provable
                    \highlightb{
                    [(\exists x) C \implies (\exists x) (B \implies C)]
                    \implies
                    \highlighto{
                    [(B \implies (\exists x) C) \implies (\exists x) (B
                    \implies C ))]}}$ -- MP 5,*

                \item[7] $\provable
                    \highlighto{
                    (B \implies (\exists x) C) \implies (\exists x) (B
                    \implies C ))}$ -- MP 2,6
            \end{itemize}
            Ešte treba dokázať formulu (*)
            \begin{itemize}
            \item[a] $\neg X \implies Z, Y \implies Z, X \implies Y,
                \highlighta{X}
                \provable Z$
            \item[b] $\neg X \implies Z, Y \implies Z, X \implies Y,
                \highlightb{\neg X}
                \provable Z$
            \item[c] $\neg X \implies Z, Y \implies Z, X \implies Y
                \provable Z$ -- veta o neutrálnej formule ($X,\neg
                X$).
            \item[d] $\provable (\neg X \implies Z) \implies (
                    (Y \implies Z) \implies ((X \implies Y) \implies
                    Z))$ -- veta o dedukcii
            \end{itemize}
        \end{itemize}
    \item[d)]
        \begin{itemize}
        \item $Q=\forall$: Ukazujeme
            $\provable (\exists x) (B \implies C) \leftrightarrow
             ((\forall x) B \implies C)$ ak $x$ nie je voľná v $C$.
            \begin{itemize}
            \item[1]
                $\provable ((\forall x) B \implies C) \leftrightarrow
                    (\neg C \implies \neg (\forall x) B)$
            \item[2]
                $\provable ((\forall x) B \implies C) \leftrightarrow
                    (\neg C \implies \neg (\forall x) \highlighta{\neg
                    \neg} B)$
            \item[3]
                $\provable ((\forall x) B \implies C) \leftrightarrow
                    (\neg C \implies \highlightb{(\exists x)} \neg B)$
            \item[4]
                $\provable 
                    (\exists x) (\highlighta{\neg C \implies \neg B})
                \leftrightarrow
                    (\highlightb{\neg C \implies (\exists x) \neg B})$
                 -- časť c) tohoto dôkazu
            \item[5]
                $\provable (\neg C \implies \neg B) \leftrightarrow (B \implies
                C)$ -- axióma A3.
            \item[6]
                $\provable 
                    (\exists x) (\highlighta{B \implies C}) \leftrightarrow
                    (\highlightb{
                        \highlighto{\neg (\exists x) \neg} B \implies 
                        \highlightp{\neg \neg}
                    C)}$ --
                    použijeme úplné delo - vetu o ekvivalentných zámenách
                    na 4,5 a to rovno dvakrát
            \item[7]
                $\provable (\exists x) (B \implies C) \leftrightarrow
                    (\highlighto{(\forall x)} B \implies C)$
            \end{itemize}

        \item $Q=\exists$: Chceli by sme
            $\provable (\exists x) (B \implies C) \leftrightarrow
             ((\forall x) B \implies C)$ ak $x$ nie je voľná v $B$.
            \begin{itemize}
            \item[1]
                $\provable ((\exists x) B \implies C) \leftrightarrow
                    (\neg C \implies \neg (\exists x) B)$
            \item[2]
                $\provable ((\exists x) B \implies C) \leftrightarrow
                    (\neg C \implies \neg (\exists x)
                                \highlighta{\neg \neg} B)$
            \item[3]
                $\provable ((\exists x) B \implies C) \leftrightarrow
                (\neg C \implies \highlightb{(\forall x)} \neg B)$
            \item[4]
                $\provable (\forall x) (\highlighta{\neg C \implies \neg B})
                \leftrightarrow
                    (\highlightb{\neg C \implies (\forall x) \neg B})$ 
                -- platí časť c) tohoto dôkazu
            \item[5]
                $\provable (\neg C \implies \neg B) \leftrightarrow (B \implies
                C)$ -- štandardná formula dokázaná vo výrokovej logike
            \item[6]
                $\provable (\forall x) (\highlighta{B \implies C})
                        \leftrightarrow
                    (\highlightb{\highlightp{\neg (\forall x) \neg} B 
                                        \implies \highlighto{\neg \neg} C})$
            \item[7]
                $\provable (\forall x) (B \implies C) \leftrightarrow
                    (\highlightp{(\exists x)} B \implies C)$
            \end{itemize}
        \end{itemize}
        \item[e)] Ukazujeme
            $\provable (Qx) (B\squareop C) \leftrightarrow (B\squareop (Qx)C)$,
            kde $x$ nie je voľná v $B$.
            Na základe operácii c), d) toto vieme dokázať, pretože platí
            \begin{align*}
                \provable (A \lor B) \leftrightarrow (\neg A \implies B) \\
                \provable (A \land B) \leftrightarrow (A \implies \neg B)
            \end{align*}
  \end{itemize}
  %%% }}}
\end{dokaz}

\begin{dokaz}[Dôkaz vety \ref{veta:prenex} o prexexných tvaroch]
%%% {{{
Budeme postupovať matematickou indukciou vzhľadom na zložitosť formuly $A$.
\begin{itemize}
    \item $A$ je atomická formula. $A$ je potom v prenexnom tvare.

    \item $A=\neg B$. Na $B$ sa vzťahuje IP, teda vieme zostrojiť
        $B'$ takú, že platí
        $\provable B \leftrightarrow B'$.
        Položíme $A'=\neg B'$ a aplikovaním prenexnej operácie b) dostaneme 
        $A''$ v správnom tvare.

    \item $A=B \implies C$. Na $B,C$ platí IP a teda existujú formuly
        $B',C'$ v prenexnom tvare, pre ktoré platí
        $\provable B \leftrightarrow B'$, $\provable C \leftrightarrow C'$.
        Nech $A' = B' \implies C'$. Na základe vety o ekvivalencii platí
        $\provable A \leftrightarrow A'$. Teraz potrebujeme dostať
        $A'$ do prenexného tvaru.
        Vezmime variant $C''$ formuly $C'$ taký, že $B',C''$ nemajú
        žiadnu spoločnú premennú.
        \begin{equation*}
            \provable A \leftrightarrow (B' \implies C'')
        \end{equation*}
        Teraz použijeme prenexné operácie c), d) a formulu
        $B' \implies C''$ prevedieme do prenexného tvaru.

    \item $A=(\forall x)B$. Z indukčného predpokladu vyplýva
        existencia $B'$, $\provable B \leftrightarrow B'$.
        Môžu nastať 2 prípady
        \begin{itemize}
        \item $x$ nie je viazaná v $B'$. Položme $A' = (\forall x) B'$
        \item $x$ je viazaná v $B'$. Potom máme $A' = B'$.
        \end{itemize}
\end{itemize}
%%% }}}
\end{dokaz}
\begin{poznamka}
    Ak $A$ obsahuje spojky 
    $\land,\lor$, môžeme použiť prenexnú operáciu e) alebo formulu nahradiť
    ekvivaletnou formulou obsahujúcou $\neg,\implies$.
    Ak sa vo vormule vyskytuje $\leftrightarrow$, nemôžeme priamo
    použiť operácie e), d) ale $A\leftrightarrow B$ prepíšeme na
    $(A\implies B) \land (B \implies A)$.
\end{poznamka}

\begin{priklad}
    \noindent
    Formula $A: B \leftrightarrow (\forall x) C$ kde $x$ nie je voľná
    v $B$ a $y$ sa nevyskytuje v $B,C$.
    \begin{align*}
       (B \implies (\forall x) C) \land ((\forall x) C \implies B) &\\
       (B \implies (\forall x) C) \land ((\forall y) C_x[y] \implies B) 
       & \mbox{ -- podľa a)} \\
       (\forall x)(B \implies C) \land (\exists y) (C_x[y] \implies B) 
       & \mbox{ -- podľa c), d)} \\
       (\forall x)(\exists y)((B \implies C) \land (C_x[y] \implies B)
    \end{align*}
\end{priklad}

\begin{priklad}
    Formula elementárnej aritmetiky:
    \begin{align*}
        (\exists x) (x=y) \implies (\exists x)((x=0) \lor
                    \neg (\exists y)(y<0))& \\
        (\exists x) (x=y) \implies (\exists u)((u=0) \lor
                    \neg (\exists v)(v<0)) &\mbox{ -- podľa a)}\\
        (\exists x) (x=y) \implies (\exists u)((u=0) \lor
                    (\forall  v) \neg(v<0)) &\mbox{ -- podľa b)}\\
        (\exists x) (x=y) \implies (\exists u)(\forall v)
                ((u=0) \lor \neg(v<0)) &\mbox{ -- podľa e)}\\
        (\forall x)(\exists u)(\forall v) (x=y) \implies 
                ((u=0) \lor \neg(v<0)) &\mbox{ -- podľa c), d)}
    \end{align*}
\end{priklad}


\section{Skolemov tvar formuly}

\begin{definicia}[Skolemov normálny tvar]$\!\!\!$\footnote{Pozn.:
    Táto definícia sa mierne líši od štandardnej, v tej sa nemôžu
    vyskytovať žiadne existenčné kvantifikátory
    }
  Uvažujme uzevretú formulu $A$.
  Ak $A$ má prefix $(\exists x_1)(\exists x_2) \dots (\exists x_k)
    (\forall x_{k+1}) (\forall x_{k+2}) \dots (\forall x_n)$, potom 
    hovoríme, že formula $A$ je vyjadrená v Skolemovom normálnom
    tvare, pričom $0 \le k \le n$.
\end{definicia}

\begin{veta}
    Nech $A$ je formula predikátovej logiky. Potom k nej môžeme
    zostrojiť formulu $A'$ v Skolemovom normálnom tvare, pričom platí
    $\provable A \iff \provable A'$.
    \label{veta:skolem}
\end{veta}

\begin{poznamka}
    Všimnime si, že predchádzajúca veta nehovorí nič o existencii
    formuly $A'$ takej, že $\provable A \lequiv A'$ ako to
    bolo u prenexného tvaru. Taká formula totiž v prípade Skolemovho
    normálneho tvaru nemusí existovať.
\end{poznamka}

\begin{definicia}[Hodnosť formuly]
    Uvažujme formulu $A$ vyjadrenú v prenexnom tvare. Potom pod
    hodnosťou formuly $A$ označíme počet veľkých kvantifikátorov,
    ktoré v prefixe predchádzajú \emph{posledný} existenčný
    kvantifikátor (počítame zľava doprava).
\end{definicia}

\begin{priklad}
    Uvažujme formulu
    \begin{equation*}
        (\forall x) (\forall y) (\exists z) 
        (\forall u) (\exists w) (\exists v)
        (\forall s) (\forall t) A
    \end{equation*}
    Jej hodnosť je 3.
\end{priklad}

\def\existsxton{(\exists x_1) (\exists x_2) \dots (\exists x_n)}
\def\forallxton{(\forall x_1) (\forall x_2) \dots (\forall x_n)}

\begin{dokaz}[vety \ref{veta:skolem}]
%%% {{{
    Budeme predpokladať, že $A$ je uzavretá (Podľa vety o uzávere
    platí $\provable A$ práve vtedy, keď je dokázateľný jej uzáver.
    Dôkaz bude prebiehať matematickou indukciou vzhľadom na hodnosť
    $m$ formuly $A$.
    \begin{itemize}
        \item $m=0$ - $A$ je v Skolemovom normálnom tvare
        \item IP: tvrdenie vety platí pre každú
        formulu ktorej hodnosť je $m-1$.
        Nech je teda $A$ tvaru
        \begin{equation}
            A: \existsxton (\forall y) B(x_1,x_2,\dots,x_n,y)
        \end{equation}
        kde $B$ \emph{nie je} nutne bezkvantifikátorová.
        Predpokladajme, že v $B$ sú voľné iba $x_1,x_2,\dots,x_n,y$.
        Keďže hodnosť $A$ je $m$, tak vo formule $B$ poslednému
        existenčnému kvantifikátoru predchádza práve $m-1$ všeobecných
        kvantifikátorov.
        Nech $P^{n+1}$ je $(n+1)$-árny predikát, ktorý sa nevyskytuje v
        $A$ (a teda ani $B$). Uvažujme nasledujúcu formulu
        \begin{equation}
            A^*: \existsxton \left[
                (\forall y) [B(x_1,\dots,x_n,y) \implies 
                    P^{(n+1)}(x_1,\dots,x_n,y)] \implies 
                    (\forall y) P^{(n+1)} (x_1,\dots,x_n,y)
            \right]
        \end{equation}
        Postupne ukážeme, že $\provable A \Leftrightarrow \; \provable
        A^*$.
        %%% {{{ \provable A <==> \provable A*
        \begin{itemize}
        \item[$\Rightarrow$]
        %%% {{{ ==>
            \begin{itemize}
            \item[0] $\provable A$

            \item[1] $\provable \highlightb{(B \implies P^{(n+1)}) \implies
                            (B \implies P^{(n+1)})}$ -- inštancia
                            $X\implies X$.

            \item[2] $\provable (X\implies (Y\implies Z)) \implies
                            (Y\implies (X\implies Z))$ -- veta o zámene
                            predpokladov

            \item[3] $\provable
                \highlightb{\underbrace{(B \implies P^{(n+1)})}_{X} \implies 
                        (\underbrace{B}_{Y} \implies
                        \underbrace{P^{(n+1)}}_{Z}))} \implies
                \highlighta{[ \underbrace{B}_{Y} \implies
                 (\underbrace{(B \implies P^{(n+1)})}_{X} \implies
                        \underbrace{P^{(n+1)}}_{Z})]}$ -- inštancia
                        kroku 2

            \item[4] $\provable
                 \highlighta{B \implies
                 ((B \implies P^{(n+1)}) \implies
                        P^{(n+1)})}$ -- MP 1,3.\footnote{
                        Toman: Ideme to obliecť, zatiaľ je
                        obnažená.}

            \item[x] $\provable \highlightc{\highlighto{(\forall y)}
                B} \implies \highlightp{\highlighto{(\forall y)}(
                (B \implies P^{(n+1)}) \implies P^{(n+1)})}$ --
                veta o distribúcii kvantifikátorov aplikovaná na
                5

            \item[y] $\highlightp{(\forall y)(
                (B \implies P^{(n+1)}) \implies P^{(n+1)})}
                \implies
                \highlightb{(\highlighto{(\forall y)}(B \implies P^{(n+1)}) 
                    \implies
                     \highlighto{(\forall y)}P^{(n+1)})}$ --
                Ukážeme nasledovne:
                %%% {{{
                 \begin{itemize}
                    \item[a] $ \highlightdim{(\forall x) (X \implies
                        Y)} \provable 
                            (\forall x) (X \implies Y)$ -- predpoklad
                            je vždy dokázateľný
                    \item[b] $ \provable \highlightdim{
                        (\forall x) (X \implies Y)}
                                \implies (X \implies Y)$ -- axióma
                                špecifikácie
                    \item[c] $ \highlightdim{(\forall x) (X \implies
                        Y)} \provable
                            X \implies Y$ -- MP a,b
                    \item[d] $ \highlightdim{(\forall x) (X \implies
                        Y)} \provable
                            (\forall x) X \implies (\forall x) Y$ --
                            pravidlo zavedenia všeobecného
                            kvantifikátora aplikované na c
                    \item[e] $ \provable \highlightdim{(\forall x) (X
                                \implies Y)}
                        \implies
                            ((\forall x) X \implies (\forall x) Y)$ --
                            veta o dedukcii v predikátovej logike.%
                            \footnote{Pozor, treba si riadne
                                premyslieť, že ju môžeme použiť. Totiž, v
                                dôkaze d sme nikne nepoužili pravidlo
                                zovšeobecnenia na premennú, ktorá by bola
                                voľná v $(\forall x) (X \implies Y)$ --
                                použili sme to iba na premennú $x$ a tá je
                                viazaná.
                            }
                 \end{itemize}
                 %%% }}}

            \item[5] $\provable \highlightc{(\forall y) B} \implies 
                \highlightb{((\forall y)
                    (B \implies P^{(n+1)}) \implies
                    (\forall y) P^{(n+1)})}$ -- pravidlo jednoduchého sylogizmu
                    aplikované na x,y

            \item[z] $\provable 
                \forallxton \left[
                \highlightc{(\forall y) B} \implies 
                \highlightb{((\forall y)
                    (B \implies P^{(n+1)}) \implies
                    (\forall y) P^{(n+1)})}
                \right] $ -- $n$-krát aplikované pravidlo zovšeobecnenia
                 na krok 5

            \item[6] $\provable \highlighta{\existsxton} 
                    \highlightc{(\forall y) B} \implies
                    \highlighta{\existsxton} \left[
                        \highlightb{
                        (\forall y) (B \implies P^{(n+1)})
                        \implies (\forall y) P^{(n+1)}) } \right]$ --
                $n$-krát použijeme
                    $\provable (\forall x)(A \implies B) \implies
                     ((\exists x) A \implies (\exists x)B)$, čo sa dá
                     dokázať nasledovne:
                     %%% {{{
                   \begin{itemize}
                    \item[c] $ \highlightdim{
                        (\forall x) (X \implies Y)}
                                \provable X \implies Y$ -- už sme
                                mali

                    \item[j] $ \provable Y \implies (\exists x) Y$ --
                        duálny tvar axiómy špecifikácie

                    \item[k] $ \highlightdim{
                        (\forall x) (X \implies Y)} \provable
                            X \implies (\exists x) Y$ -- JS c,j

                    \item[l] $ \highlightdim{
                        (\forall x) (X \implies Y)} \provable
                            (\exists x) X \implies (\exists x) Y$ --
                            pravidlo zavedenia existenčného
                            kvantifikátora, $x$ nie je voľné v
                            $(\exists x) Y$.

                    \item[m] $ \provable \highlightdim{
                        (\forall x) (X \implies Y)} \implies
                            (\exists x) X \implies (\exists x) Y$ --
                            Veta o dedukcii, čitateľ si môže
                            premyslieť, že ju môžeme použiť
                   
                   \end{itemize}
                   %%% }}}

            \item[7] $\provable
                 \highlighta{\existsxton} \left[
                 \highlightb{(\forall y) (B \implies P^{(n+1)})
                 \implies (\forall y) P^{(n+1)}} \right]$ -- MP
                 indukčného predpokladu a 6
            \end{itemize}
        %%% }}}
        \vskip 0.5cm
        \item[$\Leftarrow$]
        Predpokladáme, že je dokázateľné $\provable A^* : 
                \existsxton [(\forall y) (B \implies P^{(n+1)}) 
                \implies (\forall y) P^{(n+1)}]$.
        Treba si uvedomiť,
        že na predikát $P^{(n+1)}$ nekladieme žiadne nároky a teda je
        to dokázateľné pre ľubovoľný taký predikát. No ale
        formula $B$ sa dá chápať ako špeciálny prípad $(n+1)$-árneho
        predikátu. Preto bude dokázateľná aj inštancia
        formuly $A^*$, ktorá vyzerá nasledovne:
            $\provable  \existsxton [(\forall y) (B \implies B) 
                \implies (\forall y) B]$

        %%% {{{ <==
            \begin{itemize}
            \item[1] $\provable \highlightc{
                \existsxton [(\forall y) (B \implies B) 
                \implies (\forall y) B]}$

            \item[2] $\provable (\exists x) (A \implies B) \implies
                ((\forall x) A \implies (\exists x) B)$ -- Dá sa
                dokázať napríklad takto:
                %%% {{{
                \begin{itemize}
                \item $\provable A \implies (\neg B \implies \neg (A
                        \implies B))$

                \item $\provable (\forall x) A \implies (\forall x) (\neg B
                   \implies \neg (A \implies B))$ -- veta o zavedení
                   všeobecných kvantifikátorov

                \item $(\forall x) A \provable
                    (\forall x) (\neg B \implies \neg (A \implies B))$
                    -- veta o dedukcii

                \item $(\forall x) A \provable
                    \highlighta{(\forall x) \neg B} \implies
                    \highlightb{(\forall x) \neg (A \implies B)}$ --
                    veta o distribúcii kvantifikátorov\footnote{
                        formálne by sme mali ešte spraviť medzikrok
                        $(\forall x) A \provable \neg B \implies
                                \neg (A \implies B)$
                    }

                \item $(\forall x) A \provable
                    \highlightb{\highlightc{\neg} (\forall x) 
                                 \neg (A \implies B)} \implies
                    \highlighta{\highlightc{\neg} (\forall x) 
                                 \neg B}$ -- obmena implikácie

                \item $\highlighta{(\forall x)} A \provable
                    \highlightb{\highlightp{(\exists x)} (A \implies
                    B)} \implies
                    \highlightp{(\exists x)} B$ -- nahradenie
                    kvantifikátorov

                \item $\highlightb{(\exists x) (A \implies B)} \provable
                     \highlighta{(\forall x) A } \implies
                    (\exists x) B$ -- 2 krát veta o dedukcii
                \end{itemize}
                %%% }}}
            \item[3] $\provable \highlightc{
                \existsxton [(\forall y) (B \implies B) 
                \implies  (\forall y) B]} \implies $ \\
                \phantom{x}$\quad\quad
                \highlighta{
                [ \forallxton (\forall y_n) (B \implies B)
                \implies \existsxton (\forall y) B]}$ -- $n$-krát
                využijeme 2 nasledujúcim štýlom:
                Použijeme na dané pravidlo vetu o zavedení
                existenčného kvantifikátora. Následne, pravú stranu
                hlavnej implikácie vidíme ako kandidáta na krok 2.
                Preto ju na novom riadku rozpíšeme podľa tohoto pravidla
                a následne použítím jednoduchého sylogizmu tieto 2
                riadky zložíme. Tím dostaneme variantu 2 s pridanou
                ďalšou premennou a toto opakujeme príslušný počet
                krát.

            \item[4] $\provable \highlighta{
                \highlighto{\forallxton (\forall y_n) (B \implies B)}
                \implies \highlightp{\existsxton (\forall y) B}}$.
                -- MP 1,3
            \item[5]
                $\provable \highlighto{\forallxton (\forall y) (B
                \implies B)}$
                -- $(n+1)$-krát pravidlo zovšeobecnenia použité na 
                dokázateľnú formulu $B \implies B$
            \item[6]
                $\provable \highlightp{\existsxton (\forall y) B}$
                -- MP 4,5
            \end{itemize}
        %%% }}}
        \end{itemize}
        %%% }}}

        Ešte potrebujeme transformovať formulu $A^*$ na formulu s
        hodnosťou $m-1$. Nech
        $B:(Q_1 z_1)(Q_2 z_2) \dots (Q_l z_l) C$, kde $C$ je formula
        bez kvantifikátorov. Vieme, že hodnosť $B$, resp.
        hodnosť $(Q_1 z_1)(Q_2 z_2) \dots (Q_l z_l)$ je $m-1$
        nakoľko hodnosť $A$ je $m$ a v $A$ sa pred $B$ nachádza
        jeden všeobecný kvantifikátor.
        %%% {{{
        \begin{itemize}
        %%% pomocne makra
        \def\lefta{\scalebox{3}[1]{\highlighta{[}} }
        \def\leftb{\scalebox{3}[1]{\highlightb{[}} }
        \def\righta{\scalebox{3}[1]{\highlighta{]}} }
        \def\rightb{\scalebox{3}[1]{\highlightb{]}} }

        \item[$A^*$] 
        $=\highlightdim{\existsxton} 
            \lefta (\forall y)
            \leftb
              (Q_1 z_1) (Q_2 z_2) \dots (Q_l z_l) C \implies
              P^{(n+1)} \rightb \implies (\forall y)
              P^{(n+1)} \righta$

        \item[] $\Leftrightarrow$
        $ \highlightdim{\existsxton} \lefta
                    (\forall y) (\overline{Q_1} z_1)
                    (\overline{Q_2} z_2) \dots
                    (\overline{Q_l} z_l) \leftb
                        C \implies P^{(n+1)} \rightb
                            \implies (\forall u)
                                P^{(n+1)} \righta$ -- veta o variantoch
                                + $l$ prenexných operácií d)
        \item[] $\Leftrightarrow$
        $ \highlightdim{\existsxton} \highlightc{(\exists y)}
                \lefta
                    (\overline{Q_1} z_1)
                    (\overline{Q_2} z_2) \dots
                    (\overline{Q_l} z_l) \leftb
                        C \implies P^{(n+1)} \rightb
                            \implies (\forall u)
                                P^{(n+1)} \righta$ -- prenexná
                                operácia d) aplikovaná na $y$
        \item[] $\Leftrightarrow$
        $ \highlightdim{\existsxton} (\exists y)
                    (Q_1 z_1)
                    (Q_2 z_2) \dots
                    (Q_l z_l) \lefta \leftb
                        C \implies P^{(n+1)} \rightb
                            \implies (\forall u)
                                P^{(n+1)} \righta$ -- $l$-krát prenexná
                                operácia d)

        \item[] $\Leftrightarrow$
        $ \highlightdim{\existsxton} (\exists y)
                    (Q_1 z_1)
                    (Q_2 z_2) \dots
                    (Q_l z_l) (\forall u) \lefta \leftb
                        C \implies P^{(n+1)} \rightb
                            \implies
                                P^{(n+1)} \righta$ -- prenexná
                                operácia c) aplikovaná na $u$

            %%% zmaz pomocne makra
            \let\lefta\undefined
            \let\leftb\undefined
            \let\righta\undefined
            \let\rightb\undefined
        \end{itemize}
        %%% }}}
        Posledná formula má hodnosť $m-1$ a tak ju vieme podľa
        indukčného predpokladu previesť na Skolemov normálny tvar.
    \end{itemize}
%%% }}}    
\end{dokaz}

\section{Predikátová logika s rovnosťou}

Predikátovú logiku môžeme rozšíriť o nové axiómy, ktoré budú
hovoriť o predikáte ``=''.

\par{Axiómy rovnosti:}
\noindent
\begin{itemize}
%%% {{{
    \item[R1:] ak $x$ je premenná, potom formula $x=x$ je axióma
    \item[R2:] ak $x_1,\dots,x_k, y_1, \dots, y_k$ sú premenné a 
        $f$ je $k$-árny funkčný symbol, potom je axiómou formula
        \begin{equation*}
            (x_1 = y_1) \implies ( (x_2 = y_2) \implies ( \dots
                \implies ((x_k = y_k) \implies
                    [f(x_1,\dots,x_k) = f(y_1,\dots,y_k)] \dots )))
        \end{equation*}
    \item[R3:] ak $x_1,\dots,x_k, y_1, \dots y_k$ sú premenné a 
        $P$ je $k$-árny predikátový symbol, potom je axiómou formula
        \begin{equation*}
            (x_1 = y_1) \implies ( (x_2 = y_2) \implies ( \dots
                \implies ((x_k = y_k) \implies
                    [P(x_1,\dots,x_k) \implies P(y_1,\dots,y_k)] \dots )))
        \end{equation*}
%%% }}}        
\end{itemize}

\begin{priklad}[Teória neostrého čiastočného usporiadania $\le$]
\noindent
%%% {{{
    \begin{itemize}
        \item[1.] $(\forall x) \langle x,x\rangle \in \varphi$ -- Identita

        \item[2.] $(\forall x) (\forall y) [
            (\langle x,y \rangle \in \varphi \land
             \langle y,x \rangle \in \varphi) \implies (x=y)]$ -- 
             Antisymetrickosť

        \item[3.] $(\forall x) (\forall y) (\forall z)
            [( \langle x,y \rangle \in \varphi \land 
               \langle y,z \rangle \in \varphi) \implies
               \langle x,z \rangle \in \varphi]$. Tranzitívnosť
    \end{itemize}
    Ak pridáme trichotomickosť, dostaneme teóriu neostrého
    usporiadania:
    \begin{itemize}
        \item[4.] $(\forall x) (\forall y) [x \not=y \implies
            (\langle x,y \rangle \in \varphi \lor 
             \langle y,x \rangle \in \varphi)]$

        \item[4'] $(\forall x) (\forall y) [x = y \lor
             \langle x,y \rangle \in \varphi \lor 
             \langle y,x \rangle \in \varphi]$
    \end{itemize}
%%% }}}    
\end{priklad}

\begin{lema}
Rovnosť je symetrická a tranzitívna.
    \begin{itemize}
        \item[1] $\provable (x=y) \implies (y=x)$ -- symetria
        \item[2] $\provable (x=y) \implies ((y=z) \implies (x=z))$ --
            tranzitívnosť
    \end{itemize}
\end{lema}
\begin{dokaz}
%%% {{{
\noindent
\begin{itemize}
    \item Symetria:
        \begin{itemize}
        \item[1] $\provable (x=y) \implies ((x=x) \implies ((x=x)
                \implies (y=x)))$
            pretože \\
                $\provable (x_1=y_1) \implies ((x_2=y_2) \implies
                    ((x_1=x_2) \implies (y_1=y_2)))$ je inštancia R3
        \item[2] $\provable A \implies (B \implies (B \implies C))$.
        \item[3] $\provable B \implies (B \implies (A \implies C))$ --
            2x použité pravidlo zámeny predpokladov + veta o dedukcii
        \item[4] $\provable (x=x) \implies (x=x) \implies (x=y)
                    \implies (y=x)$
        \item[5] $\provable x=x$ -- R1
        \item[6] $\provable (x=y) \implies (y=x)$ -- 2x MP na 5,4
        \end{itemize}

    \item Tranzitívnosť:
        \begin{itemize}
            \item[1] $\provable (y=x) \implies ((z=z) \implies
                ((y=z) \implies (x=z)))$, pretože \\
                $\provable (x_1 = y_1) \implies ((x_2=y_2) \implies
                 ((x_1=x_2) \implies (y_1=y_2)))$ je inštancia R3.
            \item[2] $\provable A \implies (B \implies (C \implies D))$
            \item[3] $\provable B \implies (A \implies (C \implies D))$
                -- pravidlo zámeny predpokladov
            \item[4] $\provable \highlightc{(z=z)} \implies [
                \highlightp{(y=x) \implies ((y=z) \implies
                (x=z))}]$
            \item[5] $\provable \highlightc{z=z}$ -- R1
            \item[6] $\provable \highlightp{
                    (y=x) \implies ((y=z) \implies (x=z))}$
                -- MP 4,5
            \item[7] $\provable (A \implies B) \implies ((B \implies
                C)) \implies (A \implies C))$ -- JS
            \item[8] $\provable \highlighta{[(x=y) \implies (y=x)]}
                \implies \Big[\highlightb{
                [\highlightp{(y=x) \implies ((y=z) \implies (x=z))}] \implies
                [\highlighto{(x=y) \implies ((y=z) \implies (x=z)))}]} \Big]$ 
                -- inštancia 7
            \item[9] $\provable \highlighta{(x=y) \implies (y=x)}$ -- symetria
            \item[10] $\provable \highlightb{ \highlightp{
                [(y=x) \implies ((y=z) \implies (x=z))]} \implies
                [\highlighto{
                (x=y) \implies ((y=z) \implies (x=z))}]}$ -- MP 9,8
            \item[11] $\provable \highlighto{
                (x=y) \implies ((y=z) \implies (x=z))}$ -- MP 6,10
        \end{itemize}
\end{itemize}
%%% }}}
\end{dokaz}

\begin{veta}
    Nech $t_1,\ldots,t_n,s_1,\ldots,s_n$ sú termy, pričom platí
        $\forall i \in \{1,\dots,n\}:\; \provable t_i = s_i$.
    Potom
    \begin{itemize}
    \item[i)] Ak $t$ je term, ktorý vznikne z termu $s$ nahradením
        niektorých výskytov termov $s_i$ za $t_i$, potom 
        $\provable t=s$.
    \item[ii)] Ak $A'$ je formula, ktorá vznikne z formuly $A$
    dosadením $t_i$ za niektoré termy $s_i$, okrem prípadov, keď
    $t_i$ je premenná $x$ v kvantifikácii $(\forall x)$ 
    resp. $(\exists x)$. Potom
    $\provable A \leftrightarrow A'$.
    \end{itemize}
\end{veta}

\begin{dokaz}
%%% {{{
\noindent
\begin{itemize}
    \item[i)] Dôkaz matematickou indukciou vzhľadom na zložitosť termu $t$.
        %%% {{{
        \begin{itemize}
        \item Nech $t$ je premenná alebo $t$ je $s_i$ pre nejaké $i$.
            Potom zjavne $\provable t=s_i$.

        \item Nech term $t$ je $f(r_1,\dots,r_k)$, 
            term $s$ je $f(r_1',\dots,r_k')$.
            Pre $r_1,\dots,r_k$ platí IP, čize $\provable r_i = r_i'$.
            Potom $\provable (r_1=r_1') \implies  \dots \implies
                (r_k = r_k') \implies (f(r_1,\dots,r_k) = f(r_1',
                    \dots, r_k'))$, čo $k$-násobným použitím MP na 
                    indukčný predpoklad vedie k
                    $\provable f(r_1,\dots,r_k) = f(r_1', \dots,
                    r_k')$.
                
        \end{itemize}
        %%% }}}
    \item[ii)] Zámena termov prebieha len v atomických podformulách
        formuly $A$ (každý term je časťou nejakej atomickej podformuly).
        Máme 2 možnosti, ako vyzerá atomická podformula:
        %%% {{{
        \begin{itemize}
        \item Nech $P$ je atomická podformula tvaru
            $P(r_1,\dots,r_l)$. Potom po zámene dostaneme z 
            $P(r_1, \dots, r_k)$ formulu $P':P(r_1',\dots,r_k')$.
            Chceme ukázať $\provable P(r_1,\dots,r_k) \leftrightarrow
                                     P(r_1',\dots,r_k')$, ak vieme, že
            podľa IP platí $\forall i: \provable r_i = r_i'$
            %%% {{{
            \begin{itemize}
            \item [$\Rightarrow:$] 
                \begin{itemize}
                \item $\provable (r_1=r_1') \implies
                                 (r_2=r_2') \implies \dots \implies
                                 (r_k=r_k') \implies
                                 (P(r_1,\dots,r_k) \implies
                                 P(r_1',\dots,r_k'))$ -- inštancia R3
                \item $\provable r_i = r_i'$ -- IP
                \item $\provable P(r_1,\dots,r_k) \implies
                    P(r_1',\dots, r_k')$ -- $k$-krát MP
                \end{itemize}
            \item [$\Leftarrow:$]
                \begin{itemize}
                \item $\provable (r_1'=r_1) \implies
                                 (r_2'=r_2) \implies \dots \implies
                                 (r_k'=r_k) \implies
                                 (P(r_1',\dots,r_k') \implies
                                 P(r_1,\dots,r_k))$ -- inštancia R3
                \item $\provable r_i' = r_i$ -- IP + symetria
                \item $\provable P(r_1',\dots,r_k') \implies
                    P(r_1,\dots, r_k)$ -- $k$-krát MP
                
                \end{itemize}
            \end{itemize}
            %%% }}}
        \item Nech atomická podformula je tvaru $r_1 = r_2$. 
        Tento prípad sa dá ukázať podobne ako predchádzajúci,
        na prednáške bol ale iný dokaz:
        Chceme ukázať $\provable r_1=r_2 \leftrightarrow r_1'=r_2'$.
            %%% {{{
            \begin{itemize}
            \item[$\Rightarrow:$]
                \begin{itemize}
                \item $\provable r_1=r_1'$ -- IP
                \item $\provable r_2=r_2'$ -- IP
                \item $\provable r_1 = r_2 \implies (
                            r_2 = r_2' \implies r_1 = r_2')$ --
                            tranzitívnosť
                \item $r_1=r_2 \provable r_1=r_2'$ -- VD + MP
                \item $r_1=r_2 \provable r_2'=r_1$ -- symetria + MP
                \item $\provable r_2' = r_1 \implies (
                            r_1 = r_1' \implies r_2' = r_1')$ --
                            tranzitívnosť
                \item $r_1=r_2 \provable r_2' = r_1'$ -- 2-krát MP
                \item $r_1=r_2 \provable r_1' = r_2'$ -- symetria + MP
                \item $\provable r_1=r_2 \implies r_1'=r_2'$ -- VD
                \end{itemize}
            \item[$\Leftarrow:$]
                \begin{itemize}
                \item $\provable r_1=r_1'$ -- IP
                \item $\provable r_2=r_2'$ -- IP
                \item $\provable r_1' = r_2' \implies (
                            r_2' = r_2 \implies r_1' = r_2)$ --
                            tranzitívnosť
                \item $r_1'=r_2' \provable r_1'=r_2$ -- VD + MP
                \item $\provable r_1 = r_1' \implies (
                            r_1' = r_2 \implies r_1 = r_2)$ --
                            tranzitívnosť
                \item $r_1'=r_2' \provable r_1 = r_2$ -- 2xMP
                \item $\provable r_1'=r_2' \implies r_1=r_2$ -- VD
                \end{itemize}
            \end{itemize}
            %%% }}}
        \end{itemize}
        %%% }}}
        Dokázali sme, že atomické podformuly sú ekvivalentné. Spolu s
        vetou o ekvivalencii to ale znamená, že aj pôvodné formuly sú
        ekvivalentné.
\end{itemize}
%%% }}}
\end{dokaz}

\begin{veta}
    Majme term $t$, termy $t_1,\dots,t_n, s_1,\dots, s_n$ a formulu
    $A$.
    Potom platí
    \begin{itemize}
        \item[i)] $\provable t_1=s_2 \implies t_2=s_2 \implies
            t_n=s_n \implies (t[t_1,\dots,t_n] = t[s_1,\dots,s_n])$.
        \item[ii)] $\provable t_1=s_2 \implies t_2=s_2 \implies
            t_n=s_n \implies (A[t_1,\dots,t_n] \leftrightarrow 
                              A[s_1,\dots,s_n])$.
    \end{itemize}
    Ak navyše $x$ je premenná, ktorá nie je obsiahnutá v terme $t$,
    potom platí
    \begin{itemize}
        \item[iii)] $\provable A_x[t] \leftrightarrow 
            (\forall x)((x=t) \implies A)$
        \item[iv)] $\provable A_x[t] \leftrightarrow 
            (\exists x)((x=t) \land A)$
    \end{itemize}
\end{veta}
\begin{dokaz}
    \noindent
    \begin{itemize}
        \item[i),ii)] Ak $t_1,\dots,t_n,s_1,\dots,s_n$ neobsahujú
            premenné, tak to vyplýva priamo z predchádzajúcej vety a vety
            o dedukcii. V prípade, že tieto termy obsahujú premenné,
            tieto premenné nahradíme rôznymi konštantami, použijeme
            vetu o konštantách,\footnote{Veta o konštantách hovorí, že
            $T \provable A \iff T \provable A[c_1,\dots,c_m]$ kde
            $c_1,\dots,c_m$ sú nové konštanty. Inak povedané,
            to, že niečo vieme dokázať s premennými je ekvivalentné
            tomu, že to vieme dokázať ak premenné nahradíme novými
            konštantami} predchádzajúcu vetu a vetu o dedukcii.
        \item[iii)]
            %%% {{{
            \begin{itemize}
            \item[$\Rightarrow:$]
                \begin{itemize}
                \item $\provable x=t \implies (A \leftrightarrow
                    A_x[t])$ -- podľa ii).
                \item $\provable \underbrace{x=t}_X \implies (
                    \underbrace{A_x[t]}_Y \implies
                    \underbrace{A}_Z)$ -- platí totiž
                    $\provable (B \implies (C \leftrightarrow D))
                    \implies (B \implies (D \implies C))$
                \item $(X \implies (Y \implies Z)) \implies
                       (Y \implies (X \implies Z))$ -- pravidlo zámeny
                       predpokladov
                \item $\provable A_x[t] \implies (x=t \implies A)$ --
                    MP na pravidlo zámeny predpokladov
                \item $\provable A_x[y] \implies (\forall x)(x=t
                    \implies A)$ -- pravidlo zavedenia všeobecného
                    kvantifikátora
                \end{itemize}

            \item[$\Leftarrow:$]
                \begin{itemize}
                \item $\provable \underbrace{(\forall x)(x=t \implies
                    A)}_X \implies
                        (\underbrace{(t=t)}_Y \implies
                        \underbrace{A_x[t]}_Z)$ -- Axióma špecifikácie

                \item $(X \implies (Y \implies Z)) \implies
                       (Y \implies (X \implies Z))$ -- pravidlo zámeny
                       predpokladov

                \item $\provable (t=t) \implies ((\forall x) ((x=t)
                    \implies A) \implies A_x[t])$ -- MP

                \item $\provable t=t$ - axióma R1

                \item $\provable (\forall x) ((x=t)
                    \implies A) \implies A_x[t]$ -- MP
                \end{itemize}
            \end{itemize}
            %%% }}}
        \item[iv)]
            \begin{itemize}
            \item[1] $\provable (\forall x)((x=t) \implies A)
                \leftrightarrow A_x[t]$ -- iii)
            \item[2] $\provable (\forall x)((x=t) \implies \neg A)
                \leftrightarrow \neg A_x[t]$ -- inštancia 1
            \item[3] $\provable \neg \neg A_x[t] \leftrightarrow
                \neg (\forall x) ((x=t) \implies \neg A)$ -- obmena
                ekvivalencie
            \item[4] $\provable A_x[t] \leftrightarrow (\exists x) \neg
                ((x=t) \implies \neg A)$ -- odstránenie $\neg \neg$ a
                zámena kvantifikátora
            \item[5] $\provable \neg(X \implies \neg Y)
                \leftrightarrow  (X \land Y)$ -- rozpísanie
                $\land$
            \item[6] $\provable A_x[t] \leftrightarrow (\exists x)
                ((x=t) \land A)$ -- veta o ekvivalencii aplikovaná na
                4,5

            \end{itemize}
    \end{itemize}
\end{dokaz}



\chapter{Matematická logika}
% vim: set fdm=marker:
\section{Pravdivosť a dokázateľnosť}

\begin{definicia}[Logická platnosť formuly]
%%% {{{
    Nech $L$ je jazyk prvého rádu a $A$ je formula jazyka
    $L$. Hovoríme, že formula $A$ je \emph{logicky platná},
    označujeme $\models A$,
    ak je splnená v ľubovoľnej realizácii $m$ jazyka $L$.

    \begin{align*}
        \mathcal{M} & \models A[e] \quad \forall \mathcal{M} \\
        \mathcal{M} & \models A \\
        & \models A
    \end{align*}
%%% }}}    
\end{definicia}

\begin{poznamka}
    Formula $A$ je logicky platná, práve vtedy, keď je pravdivá
    bez ohľadu na realizáciu symbolov jazyka $L$.
\end{poznamka}

\begin{definicia}[Teória]
%%% {{{
    Nech $L$ je jazyk prvého rádu a $T$ je množina formúl
    jazyka $L$. Hovoríme, že $T$ je teória 1. rádu predikátovej logiky
    s jazykom $L$ (t.j. množina formúl $T$ je množina axióm teórie).
%%% }}}
\end{definicia}

\begin{definicia}[Model teórie]
%%% {{{
    Nech $T$ je teória v jazyku $L$, $\mathcal{M}$ je realizácia jazyka $L$.
    Hovoríme, že $\mathcal{M}$ je modelom teórie $T$, ak pre každú formulu $A$
    patriacu $T$ platí $\mathcal{M} \models A$.
%%% }}}
\end{definicia}

\begin{definicia}[Tautologický dôsledok]
%%% {{{
    Hovoríme, že formula $A$ je
    sémantickým/tautologickým dôsledkom (vetou teórie) množiny formúl $T$,
    resp. $A$ je $T$-platná, ak $A$ je splnená v každom modeli teórie $T$.
    Túto skutočnosť označujeme $T \models A$.
%%% }}}
\end{definicia}

\begin{priklad}{Teória ostrého usporiadania}
    Majme predikát < na množine $N$, tak, že platí
    \begin{itemize}
    \item[1] $(\forall x)(\forall y) ((x < y) \implies \neg (y < x))$
    - asymetrickosť
    \item[2] $(\forall x)(\forall y)(\forall z) (((x<y) \land (y<z)) \implies
        (x<z))$ - tranzitívnosť
    \item[3] $(\forall x)(\forall y)( (x\not=y) \implies ((x<y) \lor (y<x)))$ -
     trichotomickosť
    \end{itemize}
    Ak sú splnené axiómy 1,2, tak daná množina je modelom \fixme{}. Ak je
    navyše splnená aj axióma 3, množina tvorí teóriu ostrého
    usporiadania.
\end{priklad}

\begin{priklad}{Elementárna aritmetika}
%%% {{{
    Jazyk prvého rádu rozšírme o nasledujúce symboly:
    \begin{itemize}
        \item $0$ - konštanta (nulárny funkčný symbol),
        \item $S$ - $S(x)=x+1$, čiže nasledovník (unárny funkčný symbol),
        \item $+,*$ - binárne funkčné symboly
    \end{itemize}
    Axiómy elementárnej aritmetiky:
    \begin{itemize}
        \item[1.] $\neg S(x) = 0$
        \item[2.] $S(x) = S(y) \implies x=y$
        \item[3.] $x+0 = x$
        \item[4.] $x+s(y) = s(x) + y$
        \item[5.] $x * 0 = 0$
        \item[6.] $x * S(y) = (x*y)+x$
    \end{itemize}
    Zoberme si realizáciu 
        $\mathcal{N}=\langle N^+,0,S,+,* \rangle, N^+ =
        \{1,2,\dots\}$.
        Potom $\mathcal{N}$ je model pre elementárnu aritmetiku,
        zvykne sa označovať aj ako \emph{štandardný model}.
        Takejto aritmetike sa hovorí Robinsonova aritmetika.
        Ak pridáme axiómu indukcie, dostaneme Peanovu aritmetiku.
%%% }}}
\end{priklad}

\begin{poznamka}
    Ak si nebudeme všímať axiómy, ale iba relačnú štruktúru
    $\mathcal{N}=\langle N^+,0,S,+,* \rangle$.
    Zoberme namiesto $S$ konštantu $1$.
    Ciže $\mathcal{N}'=\langle N^+,0,1,+,* \rangle$.
    $\mathcal{N}'$ realizuje jazyk teórie telies.
    Lenže $\mathcal{N}'$ nie je modelom tohoto jazyka. \fixme{preco?}
\end{poznamka}

\begin{priklad}[Teória grúp]
%%% {{{
    Špeciálne symboly sú $+,-$.
    Axiómy sú:
    \begin{enumerate}
            \item $((x+y)+z) = (x+(y+z))$
            \item $(x+0) = (0+x) = x$
            \item $x, -x, x+(-x) = 0 = (-x)+x$
            \fixme{co tam znamena x,-x?}
    \end{enumerate}
%%% }}}
\end{priklad}

Ďalším cieľom je stotožniť dokázateľné formuly s tautológiami.

\begin{veta}[o korektnosti]
    Ak $T$ je teória v jazyku $L$ a ak formula $A$ je taká,
    že $T \provable A$, potom $T \models A$.
\end{veta}

\begin{dokaz}
%%% {{{
    Nech $A_1, A_2, \ldots A_n$ je odvodenie (dôkaz) formuly $A$
    z predpokladov $T$ (v teórii $T$). $A_1, A_2, \ldots A_n(A)$ kde
    $A_n=A$. Nech $\mathcal{M}$ nech je ľubovoľný model teórie $T$ 
    (čiže platí $\mathcal{M} \models A$).
    Ukážeme (indukciou podľa dĺžky dôkazu), že platí $\mathcal{M} \models A_i$
    pri predpoklade, že pre $A_j, j < i$ platí $m \models A_j$.

    Do dôkazu sa $A_i$ môže dostať niekoľkými spôsobmi:
    \begin{enumerate}
    \item $A_i \in T, \mathcal{M}$ je model $T$. Potom 
        $\mathcal{M} \models A_i$ a teda $T \models A_i$

    \item $A_i$ sa dostane do dôkazu ako axióma predikátovej logiky:
    \begin{enumerate}
        \item $A_i$ je axióma výrokovej logiky -- je poskladaná z
            atomických formúl a logických spojok $\neg$ a $\implies$. Potom
            $A_i$ je tautológia výrokovej logiky (ak formula je tautológia,
            jej pravdivostná hodnota nezávisi od ohodnotenia premenných, a
            $\mathcal{m} \models A_i$, teda $\mathcal{M} \models A_i[e]$).

        \item $A_i$ je axioma špecifikácie, teda je tvaru 
            $A_i: (\forall x) B \implies B_x[t]$, 
            $t$ je substitúcia za $x$ do $B$. 
            Chceme ukázať, že formula bude pravdivá
            pri každom ohodnotení $e$.
            Zaujíma nás prípad, kedy $(\forall x) B$ je pravdivý.\footnote{
                v opačnom prípade implikácia triviálne platí}
            To znamená, že pre ľubovoľné indivíduum $m$ platí (Z Tarského
            definície) $B[e(x/m)]$, teda $e(x)=m$.
        
            Tvrdenie zo zimného semestra:
            \begin{equation*}
                \mathcal{m} \models A_{x_1, \ldots x_n}[t_1, \ldots t_n][e] 
                    \Leftrightarrow 
                \mathcal{m} \models A[e(x_1/m_1, \ldots x_n/m_n)]
            \end{equation*}
            a
            \begin{equation*}
              t_i[e] = m_i.
            \end{equation*}

            Ak položím tvrdenie $B_x[t][e]$, zmení sa to na $B[e(x/m)]$ a
            táto formula je pravdivá v $\mathcal{m} \models B[e(x/m)]$.
            \fixme{co tym chcel autor povedat?}

        \item $A_i: (\forall x) (B \implies C) \implies (B \implies
            (\forall x) C)$ a $x$ nie je voľná v $B$.
            Mali by sme dokázať, že platí $\mathcal{M} \models A_i$.
            Zaujímavý prípad je, keď
            $\mathcal{M} \models (B \implies C)[e(x/m)]$ platí, vtedy sa
            pozeráme na platnosť $(B \implies (\forall x)C)$.
            Posledná formula je ale ekvivalentná s $\neg B \lor (\forall x) C$.
            Dôležitý je tiež predpoklad, že $x$ nie je voľná v $B$,
            a teda nezávisí od ohodnotenia viazanej premennej.
            Ak $B$ nie je pravdivá, tak disjunkcia je pravdivá a
            problém je vyriešený.
            Ak by $B$ bola pravdivá, tak by malo byť $(\forall x) C$
            pravdivé. Lenže to musí byť, inak by neplatilo
            $(\forall x) (B \implies C)$.
    \end{enumerate}

    \item $A_i$ je niektorá axióma rovnosti:
    \begin{enumerate}
        \item $A_i: x=x$. Potom $A[e]$ je $m=m$ a teda pri každom
            ohodnotení $\mathcal{M} \provable A$.

        \item $A_i: x_1 = y_1 \implies x_2 = y_2 \implies \ldots 
            \implies f(x_1, \ldots, x_n) = f(y_1, \ldots, y_n).$.
            Zaujíma nás prípad, keď  $e(x_i) = e(y_i)$, teda
            $e(x_i)=e(y_i)=m_i$. Vtedy dostávame
            \begin{equation*}
              m_1=m_1 \implies m_2=m_2 \implies \ldots
                f(m_1, \dots, m_n) =f(m_1, \dots, m_n)
            \end{equation*}

        \item $A_i: x_1 = y_1 \implies x_2 = y_2 \implies \ldots 
            \implies P(x_1, \ldots, x_n) = P(y_1, \ldots, y_n).$.
            Zaujíma nás opäť prípad, keď $x_i$ aj $y_i$ reprezentujeme 
            rovnakým indivíduom: $e(x_i/m_i)$ a $e(y_i/m_i)$.
            Potom
            $P_m(m_1, \ldots, m_n) = P_m(m_1, \ldots, m_n)$, kde $(m_1,
        \ldots m_n) \in P$ buď je alebo nie je v relácii.
    \end{enumerate}

    \item Odvodzovacie pravidlá, $A_i: A_j, A_k$, pričom $x < i$ a
        $k < i$.
        \begin{enumerate}
        \item Modus Ponens:
            $A_i: A_k$, ak vieme
            $T \provable A_j, T \provable A_j \implies A_k$.
            Podľa IP platí
            $T \models A_j, T \models A_j \implies A_k$ a vďaka
            korektnosti MP teda aj $T \models A_k$.

        \item Pravidlo zovšebecnenia: $A_i: (\forall x) A_j$.
            Z IP platí $\models A_j \Rightarrow 
                \mathcal{M} \models A_i[e(x/m)]$
            pre ľubovoľné indivíduum $m$. Tým pádom ale
            $\models (\forall x) A_j$.
            
        \end{enumerate}
    \end{enumerate}
%%% }}}
\end{dokaz}

\begin{priklad} % existencia nedokazatelnej formuly a jej negacie v el. ar.
%%% {{{
    Uvažujme elementárnu aritmetiku, ktorá má svoj štandardný
    model. Uvažujme formulu $x=0$ v $N$.
    \begin{itemize}
        \item Nech $e(x) = m \neq 0$. Potom formula $A: x=0$ nie je
        splnená pre ohodnotenie $e$ a teda $m \notmodels A[e]$.
        \item Nech $e(c) = 0$. Potom formula $A': \neg x=0$ je je
        splnená v ohodnotení $e$, t.j. $m \notmodels A'[e]$.
    \end{itemize}
    To ale znemaná, že formula $A$ ani jej negácia
    nie sú dokázateľné v elementárnej aritmetike.
%%% }}}
\end{priklad}


\begin{poznamka} % neexistencia vety o dedukcii v predikatovej logike
%%% {{{
    Vetu o dedukcii v predikátovej logike nemožno vysloviť pre
    ľubovoľnú formulu.
    \begin{align*}
            A: & \neg x=0 \\
            A': & \neg y=0 \\
    \end{align*}
    $A'$ je inštancia formuly $A$. Potom ale platí:
    \begin{align*}
            A \provable A' & \quad \mbox{čiže}\\
            \neg x=0 \provable \neg y=0
    \end{align*}
    Môžem podľa vety o dedukcii napísať, že 
    $\provable \neg x=0 \implies y=0$? Nie, pretože
    \begin{align*}
        A \implies A' \\
        e: \\
        e(x)= m \neq 0 \\
        e(x)= 0 \\
        m \notmodels (A \implies A')[e]
    \end{align*}
%%% }}}
\end{poznamka}

\begin{dosledok}{Vety o korektnosti}
    Ak teória $T$ v jazyku $L$ má model $m$, potom $T$ je
    bezposporná.
\end{dosledok}
\begin{dokaz}
%%% {{{
    Nech $m$ je model $T$ a $A$ je ľubovoľná uzavretá formula jazyka $L$.
    Potom práve jedna z formúl $A$, $\neg A$ je pravdivá v modeli $m$
    (podľa Tarského definície splniteľnosti).
    Lenže tá, ktorá nie je pravdivá, nie je ani dokázateľná 
    (podľa vety o korektnosti).
%%% }}}
\end{dokaz}

\paragraph{}
\par Tento výsledok hovorí, ze ak máme vyšetriť bezospornosť nejakej teórie,
treba nájsť jej model. Keď si zoberieme predikátovú logiku:

\par Každá realizácia jazyka predikátovej logiky je modelom (neskôr nám bude
stačiť zobrať model, ktorý má jedno jediné indivíduum).

\par
Syntaktický prístup ku dôkazu bezospornosti: Teóriu, ktorej bezospornosť chceme
dokázať, transformuje pomocou istých pravidiel do teórie, o ktorej vieme, že je
bezsosporná  (my budeme predikátovú logiku transformovať do výrokovej logiky).

\begin{priklad}
    Máme jazyk $L$ -- jazyk prvého rádu, ktorý rozšírime o konštantu $c \notin L$.
    Dostaneme tak jazyk $L' = L \union \{c\}$.
    Proces \fixme{} bude prebiehať nasledovne:
    Každý term formuly $A$ nahradíme konštantou $c$,ďalej  z formuly
    vynechávame všetky kvantifikátory a premenné bezprostredne spojené s
    kvantifikátorom.

    Každej formule $A$ na jazyku $L$ priradíme formulu $A^*$ na jazyku $L'$:
    \begin{enumerate}
            \item ak A je tvaru $A: P(t_1, \cdots t_n)$, tak $A^*: P(c, c, \cdots,c)$.
            \item $A: B \square C$, potom $A^*: B^* \square C^*$.
            \item $A: \neg B$, potom $A^*: \neg B^*$
            \item $A: (Qx) B$, potom $A^*: B^*$.
    \end{enumerate}

Ak $L$ je jazyk bez rovnosti a $\provable^L A$, potom $A^*$ je tautológia. Ak $L$
je jazyk s rovnosťou a $\provable^L A$, potom $A^*$ je tautologický dôsledok $c=c$.

\par Tvrdíme, že pre žiadnu formulu $A$ nie je $\provable A$ aj $\vdash \neg A$. Ak
by to platilo, dostali by sme sa do sporu, že vo výrokovej logike je $\provable
A^*$ aj $\provable \neg A^*$.
\fixme{toto je koniec toho prikladu?}
\end{priklad}

\section{Veta o úplnosti}
\startFIXME

G\"odelova veta, ktorú si teraz vyslovíme a dokážeme, má 2 varianty.

\begin{veta}{G\"odel, 1}
    Nech $T$ je teória v jazyku $L$ a ak $A$ je
    ľubovoľná formula jazyka $L$, potom $T \provable A \iff T \models A$,
    čiže je splnená v každom modeli teórie $T$.
\end{veta}

\begin{veta}{G\"odel, 2}
    Teória $T$ je bezosporná práve vtedy, ked $T$ má model.
\end{veta}

\begin{poznamka}
    Varianta 1 G\"odelovej vety vyplýva z varianty 2.

    Veta o dedukcii mala nasledovný dôsledok:

    Ak máme teóriu $T$, $A$ je jej formula a $A'$ je uzáver formuly $A$,
    tak $T \provable A \iff T \union \{ \neg A' \}$ je sporná teória,
    t.j. podľa vety 2 $T \union \{ \neg A' \}$ nemá model.

    Toto znamená, že v každom modeli teórie $T$ je pravdivý uzáver $A'$.
    $T \models A' \Rightarrow T \models A$. Z platnosti vety 2 vyplýva veta 1.
\end{poznamka}

\paragraph{Dôkaz:}

\par Budeme sa snažiť zostrojiť model pre teóriu, ktorá je bezosporná. Majme
bezospornú teóriu v jazyku $L$. Potrebujeme relácie, zobrazenia a $M$ --
univerzum. K dispozícii máme len syntaktické prostriedky teórie. Do úvahy
prichádzajú iba termy bez premenných -- majú jednoznačnú realizáciu (sami sebe
budú realizáciou).

$$A[e] \iff T \provable A $$.

Možné problémy pri konštrukcii modelu:
\begin{enumerate}
    \item Jazyk $L$ neobsahuje žiadne konštanty (a teda žiadne termy bez
            premenných).
    \item Ak jazyk $L$ bude jazyk s rovnosťou, môže sa stať, že v teórii $T$
            bude $T \provable t=s$, ale $t$ a $s$ sú rôzne termy bez
            premenných (rôzne konštanty).
    \item Nech $m$ je ľubovoľná realizácia jazyka $L$ a $A$ je uzavretá
            formula jazyka $L$. Potom práve jedna z formúl $A$, $\neg A$ je
            pravdivá, ale žiadna z nich nemusí byť dokázateľná v $T$.
    \item Môže sa stať, že formula $(\exists x)B$ je dokázateľná v teórii
            $T$, ale pre žiaden term $t$ bez premenných formula $B_x[t]$ nie
            je dokázateľná v $T$. To znamená, že $(\exists x)B$ 
            podľa Tarského definície pravdivosti $(\exists x)B$ je
            nepravdivá, čo je spor s vetou o korektnosti.
\end{enumerate}

Ako odstránime tieto nedostatky:
Odstránenie bodu 2 je jednoduché -- riešime vhodnou faktorizáciou, čiže
zavedieme si množinu $\tau$, čo je množina všetkých termov bez premenných
a na nej zavedieme reláciu ekvivalencie.

Body 1 a 4 riešime rozšírením o konštanty.

Tretí bod sa rieši tzv. konzervatívnym rošírením teórie (Henkinovským).
Budú to tzv.  konzervatívne teórie (na pôvodnom jazyku nezískame žiadne nové
teorémy ani nestratíme žiadne).

\begin{definicia} % uplna teoria
%%% {{{
    Hovoríme, že teória $T$ s jazykom $L$ je \emph{úplná}, ak $T$ je
    bezosporná teória a pre ľubovoľnú uzavretú formulu $A$ na jazyku
    $L$ buď $A$ alebo $\neg A$ je dokázateľná v $T$.
%%% }}}
\end{definicia}

\par Nech $T_h(m)$ je množina všetkých pravdivých uzavreté formuly $T$. 
Potom $T_h$ je úplná (nemusíme brať otvorené formuly 
-- napr. $x=0$ v aritmetike nemusí byť pravdivé, 
lebo závisí od ohodnotenia $x$).

\begin{poznamka}
 V úplnej teórii (a teda špeciálne v $T_h$) nemôže nastať problém 3,
 ktorý sme spomínali.
\end{poznamka}

\begin{definicia}
%%% {{{ Henkinova teoria
    Hovoríme, že teória $T$ s jazykom $L$ je \emph{Henkinova}, ak pre
    ľubovoľnú uzavretú formulu $(\exists x)B$ jazyka $L$ platí
    \begin{equation}
        T \provable (\exists x)B \implies B_x[c]
    \end{equation}
    pre nejaké $c$ -- konštantu.
%%% }}}
\end{definicia}

\begin{poznamka}
    Ak je teória Henkinova, tak sme vyriešili problémy 1,4.
\end{poznamka}

\begin{lema}
    Ak $T$ je úplná a Henkinova teória, tak potom $T$ má model.
\end{lema}
\begin{dokaz}
    Nech $L$ je jazyk teórie $T$, $\tau$ je množina všetkých
    termov jazyka $L$ bez premenných. Na množine $\tau$ definujem reláciu
    ekvivalencie nasledovne:
    \begin{equation}
        \forall t_1, t_2 \in \tau :
            t_1 \equiv t_2 \iff T \provable t_1 = t_2
    \end{equation}
    Rovnosť je reflexívna, symetrická, tranzitívna, teda týmto
    spôsobom definovaná relácia je relácia ekvivalencie a rozdeľuje
    nám množinu $\tau$ na triedy ekvivalencie:
    \begin{equation*}
        [t] = \{ s \in \tau: t \equiv s\}
    \end{equation*}
    Univerzum budú tvoriť termy bez premenných.
    Relačná štruktúra $M$ bude mnomžina všetkých termov bez premenných.
    Nech $f$ je ľubovoľný $n$-árny funkčný symbol a nech
     $[t_1], [t_2], \ldots, [t_n] \in M$. Definujeme funkciu $f$ v relačnej
     štruktúre $M$ nasledovne:
     \begin{equation*}
      f_m([t_1], \ldots, [t_n]) = [f(t_1, \ldots t_n)]
     \end{equation*}
    Ešte treba ukázať, že táto definícia je konzistentná, čize
    $f_m([t_1], \ldots, [t_n])=f_m([s_1], \ldots, [s_n])$ ak
    $s_i \equiv t_i$, čiže nezáleží na výbere reprezentantov.
    Taktiež je dobré si uvedomiť, že
    $[t_{x_1,\dots,x_n}[t_1,\dots,t_n]]=t[e], e(x_i/t_i)$.

    Podobne, nech $P$ je $n$-árny predikátový symbol rôzny od $=$ ($=$ sme
    si už zaviedli). V tom prípade definujeme
    \begin{equation*}
     ([t_1], \ldots [t_n]) \in P_m \iff T \provable P(t_1, \ldots, t_n)
    \end{equation*}
\end{dokaz}

\paragraph{Poznámka:} Doteraz sme pracovali s $[\ldots]$. Výsledok sa nezmení,
ak zoberiem $s_1 \in [t_1], \ldots s_n \in [t_n]$. Nezávisí to od reprezentanta
danej triedy.
\par
$$ s_1, \ldots s_n $$
$$ s_i \in [t_i], i = 1, \ldots n $$
$$ [t_1], \ldots, [t_n]) \in P_m \iff T \provable P(s_1, \dots s_n) $$

\par

\fixme{zarovnanie na ...}
\begin{align*}
    x_1, x_2, &\ldots,& x_n \mbox{-- premenné}.	\\
    e(x_i) &=& t_i, i=1,\ldots,N	\\
    t[e] &=& [t_{x_1,\ldots,x_n}][t_1,\ldots, t_n]	\\
    t[e] &=& [t]
\end{align*}

% \begin{unknown}
$A$ -- ľubovoľná uzavretá formula z $L$, potom $m \models A \iff T \provable A$.
$P(t_1, t_2, \ldots t_n)$ -- atomická formula,  uzavretá, $t_1. \ldots t_n$ sú
termy bez premenných.

$m \models A \iff ([t_1], \ldots, [t_n]) \in P_m \iff T \provable P(t_1,\ldots
t_n)$ pre ľubovoľné ohodnotenie.

\par $A: t_1 = t_2$. $m \models t_1 = t_2 \iff [t_1] = [t_2] \iff T \provable t_1 =
t_2$.

\par $A: \neg B$. Indukčný predpoklad: pre $B$ bolo tvrdenie dokázané ($m\models
B \iff T \provable B)$. Formula $A$ je pravdivá v $m$ práve vtedy, keď formula $B$
nie je pravdivá $B$ nie je pravdivá, $T \unprovable B, T \vdash \neg B$.

\par $A. B \implies C$, pričom $B$ a $Cc$ sú dokázané podľa indukčného
predpokladu. Ak platí $m \models B \implies C$, potom $m \models \neg B \lor m
\models C$. $T \provable \neg B \lor T \vdash C$.

\begin{align*}
    T \provable B \implies C \\
    \provable C \implies (B \implies C)
\end{align*}

Opačná implikácia:
\begin{equation*}
    T\provable B\implies C\Rightarrow m \models B \implies C
\end{equation*}

Teória $T$ je úplná teória a $B \implies C$ je uzavretá. Pre každú uzavretú
formulu platí, jedna z možností $T \provable \neg B$ (potom $m \models \neg B$)
alebo $T \provable B$, a z toho vyplýva, že $T \vdash C$, a teda $m \models C$.
Tvrdenie je teda dokázané.

\begin{equation}
    m \models A \iff T \provable A
\end{equation}

\par Predpokladajme, že formula $A$ je v tvare $A: (\forall x) B$. Pre každú
inštanciu formuly $B$ tvrdenie platí, $\provable (\forall x) B \implies B_x[t]$
(axióma špecifikácie). Potom aj $T \provable B_x[t]$ (lebo teória je Henkinova). $m \models B[e(x/t)]$.

\par Ak je pravdivá negácia, potom $T \provable \neg A$, čo znamená, že $(\exists
x) \neg B$. $T \provable (\exists x) \neg B \implies B_x[c]$, a to znamená, že $T
\provable \neg B_x[c]$.

\par Nech $A$ je ľubovoľná formula z teórie $T$. Pre ňu platí, že $T \provable A$
(predpoklad). Potrebujem dokázať, že $T \provable A \iff m \models A$. Zoberiem si
$A'$ -- uzáver formuly $A$. Na $A'$ sa vzťahuje tvrdenie, ktoré sme už dokázali,
a teda $m \models A' \iff T \provable A'$. Podľa vety o uzávere:
$m \models A \iff  m\models A' T \provable A' \iff T \vdash A$

\par Pre ľubovoľnú teóriu, ktorá je idealizovaná (Henkinova a úplná), viem
zostrojiť model.

\stopFIXME

\section{Rozšírenia teórie}

\startFIXME

Budeme definovať rozšírenia (ktorá ani nepomôžu, ani neuškodia :-). (Henkinova
veta, Lindenbaumova veta)

\begin{definicia}[Rozšítenie jazyka]
%%% {{{
    Jazyk $L'$ je \emph{rozšírením jazyka} $L$, ak každý
    špeciálny symbol jazyka $L$ (prípadne aj symbol $=$)
    je v jazyku $L'$ obsiahnutý s rovnakým významom a s rovnakou árnosťou.
%%% }}}
\end{definicia}

\begin{priklad}
    Nech jazyk $L'$ je jazyk s rovnosťou a špeciálnymi symbolmi $<$ a $0$.
    Jazyk $L$ má symboly $<$ a $=$. Jazyk $L'$ je rozšírením jazyka $L$.
\end{priklad}    

\begin{definicia}[Rozšítenie teórie]
%%% {{{
    Majme teóriu $T'$ s jazykom $L'$. Hovoríme, že $T'$ je
    rozšírením teórie $T$ s jazykom $L$, ak platí,
    že $L'$ je rozšírením jazyka $L$ a pre každý formulu $A$ jazyku $L$
    platí $T \provable A, T' \vdash A$. \fixme{je tam namiesto ciarky
    implikacia ci ekvivanencia}
%%% }}}    
\end{definicia}


\begin{veta}[Henkinova]
%%% {{{
    \begin{align*}
        \provable (\exists x) \neg B \implies \neg B_x[c] \\
        T \provable  \neg B_x[c] \\
        m \models  (\forall x) B \iff T (\forall c) B \\
    \end{align*}
    \fixme{Co vlastne hovori tato veta? Co su predpoklady a co
    vystup?}
    \fixme{je o chvilu zase ta ista veta?}
%%% }}}    
\end{veta}
\fixme{nema tu byt dokaz?}


\par Relačná štruktúra $<M, P_m, t_m>$ -- model pre množinu formúl $T$. Vezmeme
ľubovoľnú formulu $A \in T$ a platí $m \models A$, potom $T$ je modelom.

\par $T \provable A$. $A'$ -- uzáver A. $$ m \models A' \iff T \vdash A'$$. Ak v
modeli $m$ je splnená formula $A'$, potom musí byť splnená aj $A$. Štruktúra,
ktorú sme uvažovali, je modelom pre našu teóriu.

\paragraph{Záver:} Pre idealizovanú teóriu viem vytvoriť model.

\par Ďalšou úlohou je rozšíriť bezospornú teóriu do úplnej a Henkinovskej.

\paragraph{Príklad (čo je to model):} Teória grúp -- je to teória s rovnosťou a
používa špeciálny symbol $+$ a existuje v ňom neutrálny prvok. Axiomy:

$$\forall x,y,z (x+y)+z = x+(y+z)$$
$$(x+0) = (0+x) = x$$
$$(x+(-x)) = ((-x)+x) = 0$$.

\par Ak si vezmeme relačnú štruktúru $m=<N^+,0,+>$, táto realizuje jazyk teórie
grúp, ale nie je modelom (nie všetky axiomy nie sú splnené). Pokiaľ ale vezmeme
$m=<Z,0,+>$, táto množina tvorí model teórie grúp.  

\par Máme daný jazyk $L$, jazyk $L'$, ktorý je rozšírením $L$ (obsahuje každý
symbol z $L$ s rovnakou hodnotou a árnosťou). Vezmime jazyk $L= \{0, 1, +,
\cdot \}$. Ak nad týmto jazykom budeme pridávať axiómy, môžeme získať teórie ako
teória okruhov, teória oborov integrity, teória telies, \ldots . Zistíme, že
teória telies je rozšírenie teórie oborov integrity a tie sú rozšírením teórie
okruhov. 

\par Ak máme teóriu $T$, $T'$ je rozšírenie $T$ a je bezosporná. Potom teória
$T$ je taktiež bezosporná. Ak by sme pripustili, že teória $T$ je sporná, potom
by sa našla taká formula, že $T \provable A$ aj $T \vdash \neg A$, z čoho vyplýva
$T' \provable A$ aj $T' \vdash \neg A$.

\begin{definicia}[Konzervatívne rozšírenie]
%%% {{{
    Majme teóriu $T$ a teóriu $T'$. Hovoríme, že $T'$ je
    \emph{konzervatívne rozšírenie} teórie $T$,
    ak $T'$ je rozšírenie $T$ a pre každú formulu $A$ na jazyku $L$,
    pre ktorú platí $T' \provable A$ implikuje $T \provable A$.
%%% }}}
\end{definicia}    

\begin{lema} 
    Nech $T'$ je konzervatívne rozšírenie $T$, potom $T$ je
    bezosporná vtedy a len vtedy, keď $T'$ je bezosporná.
\end{lema}
\begin{dokaz}
%%% {{{
    Vyplýva to z toho, že ak zoberieme konzervatívne rozšírenie,
    pre ľubovoľnú formulu pôvodného jazyka platí, že:
    \begin{equation*}
        T' \provable_L A \Rightarrow T \vdash A
    \end{equation*}
%%% }}}
\end{dokaz}    

\begin{veta}[Henkinova] 
    K ľubovoľnej teórii $T$ môžeme zostrojiť Henkinovu
    teóriu $T_H$, ktorá je konzervatívnym rozšírením teórie $T$.
\end{veta}    

\begin{dokaz}
%%% {{{
    Teóriu $T$ rozšírime a priradíme k nej teóriu $T_H$ tak, že
    jednak rozšírime jazyk teórie a jednak pridáme axiómy.
    Nech teória $T$ má jazyk $L$.
    Budeme tvoriť jazyk rozšírený jazyk $L_H$ (budeme pridávať konštanty):

    \fixme{<skontroluj, malé písmená>}

    \par Vezmime si formulu $A$ napísanú v jazyku $L$, ktorá má jedinú premennú $x$
    priradíme jej premenné $c_A$ a $c_{\neg A}$ a nasledovné axiomy: 

    $
    \begin{array}{lll}
    \mbox{(*)} & (\exists x) A \implies A_x[c_A] & (H1) \\
    \mbox{(**)} & A_x{c_{\neg A}} \implies (\forall x) A & (H2) \\
    \end{array}
    $

    Stačila by nám aj jedna konštanta:
    $$B: \neg A$$
    $$ (*) (\exists x) \neg A \implies \neg A_x[c_{\neg A}]$$
    $$\neg \neg A_x[c_{\neg A}] \implies (\forall x) \neg \neg A$$
    $$A_x[c_{\neg A}] \implies (\forall x) A$$

    Alebo:
    $$B: \neg A$$
    $$\neg A_x[c_{??} \implies (\forall x) \neg A$$
    $$(\exists x) A \implies A_x[c_A]$$
    \textbf{</skontroluj, malé písmená>}

    \par $C_1$ bude množina všetkých konštánt $c_A$ a $c_{\neg A}$. Týmito konštantami
    rozšírime jazyk $L$. Induktívne pokračujeme ďalej a podobne vytvárame množinu $C_2$
    (zoberieme ľubovoľnú formulu $A$ nad rozšíreným jazykom, ktorá mala jedinú voľnú
    premennú $x$ a obsahovala aspoň jednu konštantu prvého rádu).
    Pridané konštanty budeme nazývať konštanty druhého rádu. Pokračujeme až po
    $C_n$.

    \par Vezmeme formuly, ktoré majú jedinú voľnú premennú $x$ a obsahujú nejakú
    pridávanú konštantu rádu $n$.

    Označme
    \begin{equation*}
        L(C) = L \union C_1 \union C_2 \union \ldots
    \end{equation*}
    Teóriu $T_H$ bude tvoriť rozšírenie $L(C)$ jazyka $L$, axiómy budú axiómy $T$
    ku ktorým pridáme všetky axiómy $H1$.
    Je ihneď zrejmé, že $T_H$ je rozšírenie teórie $T$.
    Potrebujeme ukázať, že $T_H$ je konzervatívne rozšírenie teórie $T$, t.j.
    že pre každú formulu $A$ jazyka $L$ platí 
    $T_H \provable A \Rightarrow T \provable A$

    Nech platí, že v $T_H$ je dokázateľná formula $A$
    a nech $B_1, B_2, \ldots ,B_n$ sú všetky Henkinovské
    axiómy vyskytujúce sa v dôkaze $A$. Keďže dôkaz je konečný, týchto
    axióm je konečný počet. Ďalej môžeme uvažovať (ako sme už ukázali),
    že všetky axiómy $B_1, \ldots, B_n$ sú typu $H1$.
    Máme teda
    \begin{equation*}
        T, B_1, \dots ,B_n \provable A
    \end{equation*}
    Keďže axiómy $B_1,\ldots,B_n$ sú uzavreté formuly, môžeme použiť vetu o dedukcii
    a dostávame
    \begin{equation*}
        T \provable B_1 \implies B_2 \implies \ldots \implies B_n \implies A
    \end{equation*}
    Axiómy navyše môžeme poprehadzovať tak, aby $B_1$ obsahovala konštantu maximálneho rádu.
    Teda $B_1$ je tvaru $(\exists x) D \implies D_x[c_D]$, pričom $c_D$ nie je je použitá
    v $B_2, \ldots, B_n$. Použijeme vetu o konštantách a dostávame
    \begin{equation*}
        T \provable ((\exists x) D \implies D_x[w] ) \implies (B_2 \implies \ldots \implies B_n \implies A)
    \end{equation*}
    kde $w$ je nová premenná.
    Teraz môžeme použiť pravidlo zavedenia existenčného kvantifikátora --
    ak $T \provable A \implies B$ a $w$ nie je voľná v $B$, potom $T\provable (\exists w) A \implies B$.
    Čiže
    \begin{equation*}
        T \provable (\exists w) ((\exists x)D \implies D_x[w]) \implies (B_2 \implies
            \ldots \implies B_n \implies A)
    \end{equation*}
    Ďalej použijeme prenexnú operáciu $(Qx) (A\implies B) \iff (A\implies
    (Qx)B)$ (za predpokladu, že $x$ nie je voľná v $A$), aby sme $w$
    preniesli dovnútra. Výsledkom je
    \begin{equation*}
        T \provable ((\exists x) D \implies (\exists w)D_x[w]) \implies B_2 \implies \ldots \implies B_n \implies A
    \end{equation*}
    Z vety o variantoch ale máme
    \begin{equation*}
        T \provable ( \exists x ) D \implies (\exists w) D_x[w] \\
    \end{equation*}
    z čoho pomocou pravidla modus ponens získame
    \begin{equation*}
        T \provable B_2 \implies B_3 \implies \ldots \implies B_n \implies A
    \end{equation*}
    Opakovaním postupu ďalších $n-1$ krát dostaneme $T \provable A$.
    Ukázali sme teda, že $T_H$ s jazykom $L(C)$ je konzervatívne rozšírenie jazyka $L$.
%%% }}}
\end{dokaz}

\stopFIXME

\begin{veta}[Lindenbaum]
    Ak $T$ je bezosporná teória s jazykom $L$, potom existuje úplné
    rozšírenie $T'$ teórie $T$ s rovnakým jazykom $L$.
\end{veta}

\begin{dokaz}[Lindenbaumovej vety]
%%% {{{
    \fixme{vyfaktorovat tento dokaz na viacero nezavislych tvrdeni}
    Nech $\mathscr{S}$ je množina všetkých uzavretých formúl jazyka $L$.
    Ďalej definujeme
    $\mathcal{B} = \{ S \mid S \subseteq \mathscr{S}, T \union S
        \mbox{ je bezosporná}\}$.
    Množina $\mathcal{B}$ je čiastočne usporiadaná na $\subseteq$ a
    má konečnú vlastnosť -- keď zoberiem ľubovoľnú podmnožinu
    $S \subseteq \mathscr{S}$, bude platiť
    \begin{equation*}
        S \in \mathcal{B} \iff \forall \mbox{ konečné } S' \in
        \mathcal{B}: S' \subseteq S
    \end{equation*}
    %
     Potrebujeme ukázať, že operácia inklúzie $\Psi, \Psi \subseteq
     \mathcal{B}\times\mathcal{B}$ je čiastočné usporiadanie, čiže je
    %
    \begin{itemize}
        \item reflexívna
            \begin{equation*}
                \forall S \in \mathcal{B}: (S,S) \in \Psi
            \end{equation*}
        \item antisymetrická
            \begin{align*}
                &(S_1,S_2) \in \Psi \land (S_2,S_1) \in \Psi
                    \then S_1 = S_2, \quad \mbox{t.j.} \\
                &S_1 \subseteq S_2 \land S_2 \subseteq S_1 \then S_1 = S_2
            \end{align*}
        \item tranzitívna. (dokážeme analogicky)
    \end{itemize}
    %
    $\mathcal{B}$ je teda čiastočne usporiadaná inklúziou $\Psi$.
    Navyše
    \begin{equation*}
        \emptyset \in \mathcal{B}, \mbox{ lebo } T \union \emptyset = T
    \end{equation*}
    Ak by totiž $\emptyset$ nebola v $\mathcal{B}$, potom by $T$ bola sporná.

    Teraz ukážeme, že množina $\mathcal{B}$ má konečnú vlastnosť:
    \begin{equation*}
        S \in \mathscr{S}, S \in \mathcal{B} \iff \forall \textrm{
        konečnú }S' \subseteq S \textrm{ platí } T \union S' 
            \textrm{ je bezosporná, teda }S' \in \mathcal{B}
    \end{equation*}
   
   \begin{itemize}
       \item[$\Rightarrow:$]
            Nech $S \in \mathcal{B}$. Potom $T \union S$ je bezosporná.
            Ak teda $S' \subseteq S$ a  $S'$ je konečná, teória 
            $T \union S'$ bude tiež bezosporná $\then S' \in \mathcal{B}$
        \item[$\Leftarrow:$]
            Predpokladajme, že pre každú konečnu podmnožinu $S' \subseteq S$
            je $T \union S'$ je bezosporná.
            Chceme ukázať $S' \in \mathcal{B} \then S \in \mathcal{B}$.

            Tvrdenie dokážeme sporom.
            Predpokladajme, že $S \not \in B$.
            Potom $T \union S$ je sporná,
            čiže pre ľubovoľnú dokázateľnú formulu A je dokázateľné
            $\neg(A \implies A)$.

            Nech $A_1, A_2, \dots, A_n$ je dôkaz 
                $\provable \neg(A \implies A)$.
            Nech $B_1, B_2, \dots, B_m$ sú tie formuly,
            ktoré sa v tom dôkaze vyskytujú a patria do množiny $S$.
            Tvrdíme, že ich počet $m \ge 1$.
            Prečo? Inak by bola priamo $T$ sporná.
            Zoberme teraz ale konečnú množinu $S'=\{B_1, \dots, B_m\}$.
            Zjavne $T \union S'$ je sporná teória.
            To je ale spor s predpokladmi.
    \end{itemize}

    \begin{lema}[Princíp maximality]
        Každá neprázdna neprázdna podmnožina $\mathcal{P}(\mathscr{S})$ s konečnou vlastnosťou 
        má maximálny prvok vzhľadom na inklúziu.
    \end{lema}

    Nech $S_0$ je maximálny prvok množiny $\mathcal{B}$ vzhľadom na inklúziu.
    Položme rozšírenie $T' = T \union S_0$. Ukážeme, že $T'$ je úplná
    teória, t.j. že pre ľubovoľnú uzavretú formulu $A$ je dokázateľná
    v teórii $T'$ buď $A$ alebo $\neg A$.

    Uvažujme sporom - nech $T'$ nie je úplná teória. 
    Potom existuje uzavretá formula $A$ taká,
    že $T'\not \provable A$ a $T' \not \provable \neg A$.
    Lenže z tohoto je evidentné (keďže $T'$ je bezosporná), že aj
    $T'' = T' \union \{\neg A\}$ je bezosporná teória.\footnote{
        \fixme{kde sme to dokazali}
        Odvoláme sa na predtým dokázaný dôsledok: ak $A$ je uzavretá,
        potom $T \provable A \iff T \union \{\neg A'\}$ je sporná teória.
    }
    To je ale v spore s tým, že $S_0$ je maximálny prvok $\mathcal{B}$.

    Dosiahli sme zúplnenie teórie. Pre teóriu $T$ sme získali $T'$,
    ktorá je bezosporná a je úplným rozšírením $T$ na rovnakom jazyku.
%%% }}}    
\end{dokaz}

\startFIXME

Keď máme $T$, $L$, vieme to rozšíriť na $T_H$ s jazykom $L(C)$. Keď máme $T$ bezospornú, $T(H)$ s jazykom $L(C)$ je tiež bezosporná. Dokážem vytvoriť Hankinovskú úplnú $T_H'$, $L(C)$ a pre ňu model $m$.

$T, L$ - bezosp.

$T', L'$ - úplná hakinovská

$T', L' - m'$

$T'$ je bezosp. rozšírenie $T$, potom $T \provable A \iff T' \vdash A$ Potom platí, že keď zooberiem ľub. $A \in T$, potom  $m' \vdash A$. Pre ľubov. rozšírenie $T$, $T'$ $A\in T, T' \vdash A \implies m' \models A, m, v T$.

\paragraph{Definícia.}

$T,L$, $T',L'$, $L'$ je rozš. $L$. Nech $m'$ je realizácia jaz. $L'$, redukovaním štruktúry $m'$ na jaz. $L$ získame realiz. $T$ v jaz. $L$.

$m$: univerzum $m$ bude to isté univerzum ako univerzum $m'$.

$m$ obsahuje iba také relácie a zobrazenia, ktoré realizujú špecálne symboly jazyka $L$ v realizácii $m'$. Tzn. ak $f$ je ľub. n-ány funkčný symbol $L$, $f_m'$ je zobrazenie, ktoré realizuje $f$ v $m'$, potom zostáva realizáciou $f$ v $m$.

Podobne s predikátom $P$.

$m$ vzniklo redukciou $m'$ na jazyk $L$. $(m: m' \triangle L)$
\footnote{na prednáške bolo miesto $\triangle$ použíté niečo ako $\land$, ale jedna nožička bola dlhšia ako druhá}.
$m'$ je expanzia realizácie $m$

\paragraph{Poznámka.}
$m$ je redukcia $m'$ na jazyk $L$ ($m'\triangle L$), $A$,$L$. Chceme ukázať $m \models A \iff m' \models A$. Všimnime si, že univerzum, t.j. aj ohodnotenia a symboly, sú tie isté. Preto, ak $m' \models A$, potom aj $m \models A$ a obrátene. Teraz dokážeme, že každá bezosporná teória má model.

\paragraph{Veta.}
Nech $T'$ je rozšírenie teórie $T, L$. Ak $m$ je model $T'$ a $m:m'\triangle L$, tak $m$ je model $T$.

\paragraph{Dôkaz} To, čo sme tu už porozprávali. $m'$ je model $T$',
$A \in T$, $A$ je axioma $T$, $m' \models A \iff M \models A$.

Nasledujú dve vety o kompaktnosti a ich dôsledky. Dávajú nám konečnú charakterizáciu splniteľnosti.

\stopFIXME

\section{Veta o kompaktnosti}

\begin{veta}[O kompaktnosti]
    Nech $T$ je množina formúl jazyka $L$, $A$ je formula $L$. Potom
    $T \models A \iff \exists \mbox{konečná }T'\subseteq T: T' \models A$.
\end{veta}

\begin{dokaz}
    Podľa G\"odelovej vety o úplnosti v predikátovej logike platí
    $T \models A \iff T \provable A$.
    Nech $A_1, A_2, \dots , A_n$ je konečná postupnosť formúl -
    dôkaz $A$. Zoberme $T' = \{B_1, B_2, \dots, B_m\}$, kde $B_i$ sú axiómy
    $T$ nachádzajúce sa v dôkaze $A$.

    Potom môžeme písať $T' \provable A \iff T' \models A$.
\end{dokaz}

\begin{priklad}
    Na základe vety o kompaktnosti ukážeme, že teóriu telies
    charakteristiky 0 nedokážeme v teórii prvého rádu zapísať
    pomocou konečného počtu axióm.

    Nech $T$ je teória telies charakteristiky 0 (s rovnosťou),
    $L=\{0,1,+,*\}$.
    
    Označme si pre $p\in N, p\ge 1$ výraz $p \times 1$ ako
    \begin{equation*}
        p \times 1 : \underbrace{1 + (1 + ( 1 + \dots (1 + 1}_{p}) \dots )
    \end{equation*}
    Ďalej si označme formulu $P_p$ ako $P_p: p \times 1 = 0$.
    Telesá charakteristiky $p$ môžeme popísať ako
    \begin{equation*}
     \models \neg P_1 \land \neg P_2 \land \dots \land \neg P_{p-1}
        \land P_p
    \end{equation*}
    Ak teleso nemá konečnú charaktetistiku, tak jeho charakteristika
    je 0.

    Nech $A$ je formula teórie telies charakteristiky 0. Pre $A$ platí
    \begin{equation*}
     T \models A \iff T' \subseteq T, T'\mbox{ konečná}
    \end{equation*}
    Teória $T'$ obsahuje len konečne veľa axióm $\neg P_p$.
    Označme $m$ najväčší index axiomy, ktorá patrí to $T'$.
    Potom formula $A$ je splnená aj vo všetkých telesách dostatočne
    veľkej konečnej charakteristiky 
    (presnejšie povedané charakteristiky $>m$).
    Ak teda formula $A$ je teorémou telies charakteristiky $0$,
    tak platí aj vo všetkých telesách dostatočne veľkej charakteristiky.

    Uvažujme teraz axiómy $A_1, \dots A_n$ teórie telies s
    charakteristikou 0.
    \begin{equation*}
        T' \models A_i \implies T' \models A_1 \land A_2 \land \cdots
        \land A_n \equiv B
    \end{equation*}
    Formula $B$ teda platí vo všetkých telesách charakteristiky 0.
    Lenže ako sme už povedali, tým pádom platí aj vo všetkých telesách
    dostatočnej charakteristiky. To je ale spor, pretože axiómy telies
    charakteristiky 0 nemôžu platiť v telesách konečnej
    charakteristiky.

    Záverom teda je, že logikou prvého rádu nemôžeme 
    konečným počtom axiom charakterizovať teóriu telies.
    V logike prvého rádu totiž kvantifikujeme len individuá.
\end{priklad}

\begin{poznamka}
    Ak zavedieme ďalší druh premenných, prirodzené čísla,
    nekonečne veľa axiom typu $\neg P_p$ môžeme nahradiť jednou formulou:
    \begin{equation*}
        (\forall p)(p\times 1\not=0)
    \end{equation*}
    Toto sa v literatúre terminologicky nazýva
    \emph{slabá logika druhého rádu}. V nej
    vieme konečne axiomatizovať teóriu telies charakteristiky 0.
    Tým pádom 
    \emph{v slabej logike druhého rádu neplatí veta o kompaktnosti}.
\end{poznamka}

\begin{poznamka}
    Poznáme doteraz tieto logiky:
    \begin{itemize}
    \item logika 1. rádu -- premenné len pre indivíduá
    \item logika 2. rádu -- premenné pre množiny indivíduí
    \item slabá logika 2. rádu -- premenné pre indivíduá a prirodzené
        čísla
    \end{itemize}
\end{poznamka}

\medskip \noindent
Mali sme Godelovu vetu v dvoch tvaroch. Prvý bol
$ T \models A \iff T \provable A$, druhý ``Teória je bezosporná, keď
má model''. Preto aj veta o kompaktnosti bude mať druhý tvar.

\begin{veta}[2. tvar vety o kompaktnosti]
%%% {{{
    Nech $T$ je množina formúl jazyka $L$. 
    Model teórie $T$ existuje práve vtedy, keď každá konečná podmnožina
    $T' \subseteq T$ má model.
%%% }}}    
\end{veta}

\begin{dokaz}
%%% {{{
    Z 2. G\"odelovej vety vyplýva, že $T$ má model $\iff$ je bezosporná.
    Z lindenbauma zase, že $T$ je bezosporná $\iff$ každá konečná
    $T' \subseteq T$ je bezosporná, čo opäť z G\"odelovou vety
    $\iff$ každá konečná podmnožina $T'$ má model.
%%% }}}
\end{dokaz}

\begin{priklad}
    Pomocou vety o kompaktnosti tohto tvaru dokážeme zaručiť,
    že pre Peanovu aritmetiku existujú aj neštandardné modely.
    Zopakujme si najskôr jej axiómy:

    \begin{enumerate}
    \item $\provable \neg (S(x) = 0)$
    \item $\provable (S(x) = S(y)) \implies (x=y)$
    \item $\provable (x+0) = x$
    \item $\provable (x+S(y)) = S(x+y)$
    \item $\provable (x * 0) = 0$
    \item $\provable (x * S(y)) = ((x*y)+x)$
    \par \noindent Doteraz sme dostali Robinsonovu aritmetiku. Peanovu
    aritmetiku dostávame pridaním axiómy indukcie:
    \item $\provable A_x[0] \implies (\forall x)\Big(A \implies (
            A_x[S(x)] \implies (\forall x)A)\Big)$
    \end{enumerate}
    Táto aritmetika má model 
    $\mathcal{N} = \langle 0,1,S,+,*,\mathbb{N}^+ \rangle$.

    Ukážeme, že Peanova aritmetika má aj neštandardné modely.
    Definujme množinu ``numerálov'' $\bar{n}$ -- termov jazyka $L$ ako
    \begin{itemize}
    \item $\bar{0} = 0$
    \item $\overline{n+1} = S(\bar{n}) = 
        \underbrace{S(S( \cdots S}_{(n+1)\mbox{-krát}}(0) \cdots ))$
    \end{itemize}

    K $L$ zostrojíme $L_c$ tak, že do $L$ pridáme konštantu $c$ a
    teóriu rozšírime o axiómy
    \begin{equation*}
        \provable C_n: c \not= \bar{n}
    \end{equation*}

    Vidíme, že každá konečná podmnožina $T_c$ má model,
    ktorý vznikne expanziou štandardného modelu $n$.
    Konštantu $c$ realizujeme indivíduom $n$, ktoré nepatrí do
    indivíduí použitých v axiómach.
    Čiže každá $T' \subseteq T_c$ má model a na základe vety o kompaktnosti 
    aj $T_c$ má model. Teraz si môžeme všimnúť, že model $T_c$ bude rôzny
    od štandardného modelu a nebude s ním ani izomorfný.
    Vieme teda, že neštandardný model existuje, ale nevieme ho skonštruovať.
\end{priklad}


\chapter{Dokazovanie formúl -- Metóda rezolvent}
\section{Metóda rezolvent}

Celú túto kapitolu sa budeme venovať sémantike formúl logiky prvého
rádu.

Keď uvažujeme výrokovú logiku, interpretujeme formulu $A$ funkciou
$\bar{v}$ -- valuáciou.
Formulu sme nazvali tautológiou, ak $\bar{v}(A)$, ak pre každú
interpretáciu dostaneme pravdu - treba preskúmať $2^n$ interpretácii.
V predikátovej logike je to ešte horšie.

Zopakujme si, čo už vieme. V predikátovej logike máme funkčné a
predikátové symboly. Univerzum budem označovať netradične ako $D$.
Znakom $M$ budeme totiž označovať jadro(maticu) formuly v prenexnom
tvare.

$n$-árny funkčný symbol bude funkcie $f(x_1,\dots,x_n): D^n \rightarrow D$.
$n$-árny predikátový symbol bude funkcia
$P(x_1,\dots,x_n): D^n \rightarrow \{0,1\}$.

Vo formule rozlišujeme voľné a viazané premenné.
Ak má formula voľné premenné, neviem určiť jej pravdivostnú hodnotu,
iba ak za všetky voľné premenné dosadím konštanty.

Ďalej, ku každej formule A viem zostrojiť takú,
ktorá je v prefixovom (prenexnom) tvare -- t.j. je 
prefix + jadro: $(Q_1 x_+)\cdots(Q_n x_n)M$.

Na upravenie formuly do tohoto tvaru používame prenexné operácie.
Skutočnosť, že premenná $x$ má voľný výskyt vo formule $A$ budeme
značiť ako $A(x)$, v opačnom prípade budeme písať jednoducho $B$.

Prenexné operácie:
\begin{align*}
    (Qx)A(x) \lor B &\equiv (Qx)(A(x) \lor B) \tag{1a}\\
    (Qx)A(x) \land B &\equiv (Qx)(A(x) \land B) \tag{1b}\\
    \neg (\forall x) A(x) &\equiv (\exists x) \neg A(x) \tag{2a}\\
    \neg (\exists x) A(x) &\equiv (\forall x) \neg A(x) \tag{2b}\\
    (\forall x) A(x) \land (\forall x) B(x) &\equiv 
        (\forall x) (A(x) \land B(x)) \tag{3a}\\
    (\forall x) A(x) \lor (\forall x) B(x) &\equiv
        (\forall x) (A(x) \lor B(x)) \tag{3b}
\end{align*}
\begin{poznamka}
    S predchádzajúcimi operáciami 3a,3b treba pracovať pozorne. Im veľmi
    podobné úpravy totiž neplatia:

    \begin{align*}
      (\forall x) A(x) \lor (\forall x) B(x)  & \not \equiv
        (\forall x) (A(x) \lor B(x)) \tag{x1}\\
      (\exists x) A(x) \land (\exists x) B(x) & \not \equiv
        (\exists x) (A(x) \land B(x)) \tag{x2}
    \end{align*}
    Zoberme si napríklad $D=\{1,2\}$. Ak položíme
    $A(1)=1, A(2)=0, B(1)=0, B(1)=1$, dostaneme, že pravá strana x1 bude
    platiť, zatiaľ čo ľavá strana nie. Pre x2 naopak.
\end{poznamka}

Pretože môžeme substituovať za premenné --
    $(\forall x) B(x) \equiv (\forall z) B(z)$, môžeme predchádzajúce
prenexné operácie zhrnúť do nasledujúcich všeobecných transformácií.
\begin{align*}
    (Q_1 x)A(x) \lor (Q_2 x)B(x) \equiv 
        (Q_1 x)(Q_2 z)(A(x) \lor B(z)) \tag{4a} \\
    (Q_3 x)A(x) \land (Q_4 x)B(x) \equiv
        (Q_3 x)(Q_4 z)(A(x) \land B(z)) \tag{4b}
\end{align*} 
kde v oboch prípadoch $z$ je premenná, ktorá sa nevyskytuje voľne v A
(a ani v pôvodnom B(x)).
Tento najvšeobecnejší tvar ale zbytočne pridáva premenné a preto je
vhodný iba v prípadoch, keď nefungujú operácie 1 až 3.
\subsection{Algoritmus na zostrojenie prenexného tvaru}
\startFIXME
\paragraph{Krok 1}
Odstránenie implikácií a ekvivalencií:
\begin{align*}
    A \leftrightarrow B &\equiv (A \implies B) \land (B \implies B) \\
    A \implies B &\equiv \neg A \lor B
\end{align*}

\paragraph{Krok 2}
Odstránenie dvojitej negácie a presun negácie k formule.
\begin{align*}
    \neg (\neg A) &\equiv A \\
    \neg (A \lor B) &\equiv (\neg A \land \neg B) \\
    \neg (A \land B) &\equiv (\neg A \lor \neg B) \\
    \neg (\forall x) A(x) &\equiv (\exists x) \neg A(x) \\
    \neg (\exists x) A(x) &\equiv (\forall x) \neg A(x)
\end{align*}

\paragraph{Krok 3}
Premenovanie viazaných premenných, ak je to potrebné.

\paragraph{Krok 4}
Použijeme zákony:
\begin{align*}
    (Q x) A(x) \lor B &\equiv (Q x)(A(x) \lor B) \\
    (Q x)A(x) \land B &\equiv (Q x)(A(x) \land B) \\
    (\forall x) A(x) \land (\forall x)B(x) &\equiv 
        (\forall x)(A(x) \land B(x)) \\
    (\exists x)A(x) \lor (\exists x)B(x) &\equiv 
        (\exists x) (A(x)\lor B(x)) \\
    (Q_1 x) A(x) \lor (Q_2 x)B(x) &\equiv
        (Q_1 x)(Q_2 z) (A(x)\lor B(z)) \mbox{ $z$ nie je v $A(x)$}. \\
    (Q_3 x) A(x) \land (Q_4 x)(B(X) &\equiv 
        (Q_3 x)(Q_4 z)(A(x) \land (B(z))
\end{align*}

\begin{priklad}
    \begin{align*}
        (\forall x)(\forall y) \left( (\exists z)(P(x,z) \land
            P(y,z)\right) \implies (\exists u) Q(x,y,u)) \equiv \\
        (\forall x)(\forall y)\left(\neg (\exists z)(P(x,z)\land 
            P(y,z))\lor(\exists u)Q(x,y,u)\right) \equiv \\
        (\forall x)(\forall y)\left((\forall z)(\neg P(x,z)\lor 
            \neg P(y,z)\right)\lor(\exists u)Q(x,y,u)) \equiv \\
        (\forall x)(\forall y)(\forall z)(\exists u)(\neg P(x,z) \lor 
            \neg P(y,z) \lor Q(x,y,u))
    \end{align*}
\end{priklad}

\section{Herbrandova veta - história}
\startFIXME
\par Leibniz (1646 -- 1716), Peáno (1900), Hilbertova škola (1920), Herbrand
(1930). Fundamentálny výsledok, ktorý dal odpoveď na otázku, či existuje
všeobecná procedúra, ktorá vie zistiť, či je daná formula tautológia, dali
Church a Turing (1936, nezávisle od seba) -- ukázali, že v predikátovej logike
nie je problém rozhodnuteľnosti riešiteľný.
\par Herbrandova metóda je založená na procedúre vyvrátenia. Ide o nájdenie
interpretácie, pri ktorej daná formula nie je splnená. Ak formula nie je
tautológia, v konečnom počte krokov sa procedúra zastaví, pokiaľ ale nie je,
procedúra pokračuje vo výpočte.

\par Gilmore (1960), Davis a Putnem (1960) -- navrhovali procedúry, ktoré by
boli schopné overovať tautologickosť formúl; fungovali ale iba pre jednoduchšie
formuly.

\par V rokoch 1965 -- 1968 Robinson zaviedol pojem \emph{resolventa} (v podstate
iným spôsobom zapísane pravidlo \emph{modus ponens}.
\stopFIXME

\section{Skolemovské štandardné formuly}

Budeme sa zaoberať procedúram založených na Herbrandovej myšlienke vyvrátenia 
(nájdenie interpretácie, ktorá vyvracia danú formulu).
Davis a Putnam pracovali s formulou v prefixovom
normálnom tvare a vychádzali z nasledujúcich hľadísk:
\begin{enumerate}
%%% {{{
    \item Každá formula logiky 1. rádu
        môže byť vyjadrená v prefixovom normálnom tvare.

    \item Pretože matica (jadro) neobsahuje kvantifikátory, môže byť
        vyjadrená v konjunktívnej normálnej forme.

    \item Zachovajúc nesplniteľnosť formúly môžeme z nej eliminovať
        existenčné kvantifikátory pomocou tzv. skolemovských funkcií (keď máme
        premennú vyjadrenú v prefixovom normálnom tvare a eliminujeme z nej
        existenčný kvantifikátor, dostaneme formulu, ktorá je nesplniteľná
        vtedy, keď je nesplniteľná pôvodná formula -- nesplniteľnosť
        transformovanej formuly sa robí ľahšie).
%%% }}}        
\end{enumerate}

Transformácia formuly:
Majme formulu $A: (Q_1 x_1) (Q_2 x_2) \ldots (Q_n x_n) M$,
kde $M$ je konjunktívnom normálnom tvare.
Nech $Q_r$ je existenčný kvantifikátor v prefixe $(Q_1 x_1) \ldots
(Q_n x_n)$, $1 \le r \le n$.

Uvažujme, že vľavo od $Q_r$ nestojí žiaden veľký kvantifikátor.
Vtedy, ak z prefixu vynechám $(Q_r x_r)$ a $x_r$ v $M$ nahradím
novou konštantou $c_r$ a mám formulu, ktorá je nesplniteľná práve
vtedy, keď bola nesplniteľná pôvodná formula.

Na druhej strane, naľavo od $Q_r$ sa môžu vyskytovať všeobecné
kvantifikátory, $Q_{s_1}, Q_{s_2}, \dots Q_{S_m}$,
kde $1 \le s_1 < s_2 < \dots < s_m < r$.
Vtedy opäť odstránime $(Q_r x_r)$ z prefixu a
za premennú $x_r$ dosadíme do $M$ term 
$f(x_{s_1}, \ldots, x_{s_m})$. Túto funkciu voláme Skolemova funkcie a
konštanty $c$ voláme Skolemove konštanty.

Tento postup môžem zopakovať pre všetky existenčné kvantifikátory a
výsledok nazveme štandardnou skolemovskou formou.

\begin{priklad}
%%% {{{
    Nájdite štandardnú formu pre formuly:
    \begin{equation*}
        (\exists x) (\forall y) (\forall z) (\exists u) (\forall v) (\exists w)
        P(x,y,z,u,v,w)
    \end{equation*}
    Premennú $x$ budeme nahrádzať konštantou $c$ (ktorá nie je v predikáte),
    premennú $u$ nahradíme $f(y,z)$ a
    premennú $w$ nahradíme funkciou $g(y,z,v)$. Výsledkom je
    \begin{equation*}
        (\forall y)(\forall z)(\forall w) P(c,y,z,f(y,z),v,g(y,z,v))
    \end{equation*}
%%% }}}    
\end{priklad}

\begin{priklad}
%%% {{{
    Nájdite štandardnú formu pre formulu:
    \begin{equation*}
        (\forall x)(\exists y) (\exists z) 
        \left[ (\neg P(x,y) \land Q(x,z))\lor R(x,y,z)\right]
    \end{equation*}
    Využijeme distributívne zákony pre logické spojky a dostaneme
    \begin{equation*}
        (\forall x) (\exists y) (\exists z)
        \left[ \Big(\neg P(x,y) \lor R(x,y,z)\Big) \land
            \Big(Q(x,z)\lor R(x,y,z)\Big)\right]
    \end{equation*}
    a po Skolemizácii máme
    \begin{equation*}
        (\forall x) \left[
            \Big(\neg P(x,f(x)) \lor R(x,f(x),g(x)) \Big) \land
            \Big( Q(x,g(x))\lor R(x,f(x),g(x) \Big) \right]
    \end{equation*}
%%% }}}    
\end{priklad}


\begin{definicia}[Literál]
    Premenná, alebo jej negácia.
\end{definicia}
\begin{definicia}[Klauzula]
    Disjunkcia literálov.
    Špeciálne \emph{jednotková klauzula} (tiež jednoliterálna klauzula)
    obsahuje jeden literál. 
\end{definicia}

\begin{definicia}[Prázdna klauzula]
    Klauzala neobsahujúca žiaden literál, (označujeme ju $\varepsilon$).
    Prázdna klauzula je nesplniteľná, pretože neobsahuje žiaden
    literál, ktorý by mohol byť splniteľný.    
\end{definicia}
\begin{poznamka}
    Ak sa nad tým, zamyslíme, definovali sme si nesplniteľnosť
    prázdnej klauzuly tak, aby výraz $\varepsilon \lor K$ bol
    ekvivalentný výrazu $K$.
\end{poznamka}

\begin{poznamka}
    Štandardná skolemovská forma sa dá zadať aj v množinovom zápise.
    Takýto zápis budeme často využívať.
\end{poznamka}

\begin{poznamka}
    Formula môže mať viacero štandardných foriem.
\end{poznamka}

\begin{priklad}
%%% {{{
    Uvažujme klauzulu $P \land Q \land \neg R$.\footnote{
    Toman na tomto mieste píše $P \lor Q \lor \neg R$, som si však
    takmer istý, že to má byť opačne.
    }
    Ako vieme ju zapísať ako
    množinu literálov v konjunkcii nasledovne
    $\{P, Q, \neg R\}$.
    Podobne, pre štandardnú formu
    \begin{equation*}
        (\forall x) \left[
            \Big(\neg P(x,f(x)) \lor R(x,f(x),g(x)) \Big) \land
            \Big( Q(x,g(x))\lor R(x,f(x),g(x) \Big) \right]
    \end{equation*}
    z minulého príkladu vieme spraviť ekvivalentný množinový zápis ako
    \begin{equation*}
        \left\{
            \neg P(x,f(x)) \lor R(x,f(x),g(x)), \quad
            Q(x,g(x))\lor R(x,f(x),g(x) 
        \right\}
    \end{equation*}
%%% }}}    
\end{priklad}

\begin{veta}
    Nech $\mathscr{S}$ je množina klauzúl, ktoré tvoria štandardnú
    formu formuly $A$. Potom $A$ nie je splniteľná $\iff$ $\mathscr{S}$
    nie je splniteľná.
\end{veta}

\begin{dokaz}
%%% {{{
    Uvažujme formulu $A$ tvaru $A: (Q_1 x_1) \dots (Q_n x_n) M[x_1, \dots x_n]$.
    Nech $Q_r$ je prvý existenčný kvantifikátor, čiže $A$ je tvaru
    $(\forall x_1) (\forall x_2) \dots (\forall x_{r-1}) (\exists x_r)
     (Q_{r+1} x_{r+1}) \dots (Q_n x_n) M[x_1, \dots x_n]$.
    Existenčný kvantifikátor nahradíme novou $(r-1)$-árnou funkciou $f$ a
    dostávame formulu $A_1$ v tvare
    \begin{equation*}
        A_1: (\forall x_1) \dots (\forall x_{r-1}) (Q_{r+1} x_{r+1})  \dots (Q_n x_n)
                M[x_1, \dots, x_{r-1}, f(x_1, \dots, x_{r-1}), x_{r+1}, x_n]
    \end{equation*}

    \noindent Tvrdíme, že $A$ nie je splniteľná práve vtedy keď $A_1$ nie je splniteľná:
    \begin{itemize}
    %%% {{{
    \item[$\Rightarrow:$] Sporom. Nech $A_1$ je splniteľná.
        Potom $\forall x_1,\dots,x_{r-1}$ platí formula $A_1$, čize
        existuje taká interpretácia $I$, že $I$ vyhovuje $A_1$ ( $A_1$ je pravdivé v $I$).
        To znamená, že pre všetky $x_1 , \dots x_{r-1}$ existuje také $x_r$,
        ktoré je rovné $f(x_1, \dots, x_{r-1})$.
        Tá istá interpretácia teda vyhovuje aj pôvodnej formule $A$.

    \item[$\Leftarrow:$] Opäť sporom. Nech $A$ je splniteľná.
        Teda existuje interpretácia $I$, ktorá vyhovuje formule $A$.
        Potom pre každé $x_1, \dots, x_{r-1}$ existuje také $x_r$, že platí 
        $M[x_1, \dots, x_{r-1}, x_r, \dots, x_n]$.
        Rozšírme $I$ o $(r-1)$-árny funkčný 
        symbol $f$ tak, že $f(x_1, \dots, x_{r-1}) = x_r$. Potom $I'$
        vyhovyje formule $A_1$ a teda dostávame spor.
    %%% }}}
    \end{itemize}

    \noindent Pokračujme ďalej v dôkaze:
    Označme pôvodnú formulu ako $A_0$. Potom môžeme dosiahnuť štandardnú formu
    $S$ ako postupnosť $A_0, A_1, \dots, A_m$, kde každý krok je zámena prvého
    existenčného kvantifikátora vo formule $A_i$.
    Ak $A_i$ nie je splniteľná, tak ani $A_{i+1}$ nie je splniteľná.
    Tým pádom $A_0$ nie je splniteľná $\iff$ $A_1$ nie je splniteľná $\iff
    \dots \iff$ $A_m\equiv S$ nie je splniteľná.
    No a $S$ nie je splniteľná práve vtedy keď $\mathscr{S}$ nie je
    splniteľná.
%%% }}}
\end{dokaz}

\begin{poznamka}
    Ak $A$ \emph{nie je splniteľná}, tak splniteľnosť $A$ je
    ekvivalentná splniteľnosti $\mathscr{S}$, teda môžeme písať
    $A \equiv \mathscr{S}$.
\end{poznamka}

\begin{poznamka}
    Ak $A$ je splniteľná, tak tvrdenie neplatí.
    Majme formulu $A$, ktorá je $A: (\exists x) P(x)$.
    Tým pádom $S$ je $S: P(a)$ ($x$ sme nahradili skolemovskou
    konštantou $a$).
    Vezmime interpretáciu s univerzom $D = \{1, 2 \}$,
    pričom $a$ priradíme v $I$ hodnotu $1$.
    Hodnoty pre $P$ budú $P(1)=false, P(2)=true$.
    Potom ale $A$ je splniteľná v $I$ (existuje také $x$, menovite
    $2$), lenže $P(a)$ nie je splniteľná.
    Preto $A \not \equiv S$.
\end{poznamka}

\begin{poznamka}
    Z predchádzajúcej poznámky sa ukazuje, že metódy dokazovania
    formúl budú založené na vyvrátení negácie.
\end{poznamka}

\begin{poznamka}
    Majme formulu $A: A_1 \land A_2 \land \ldots \land A_n$.
    Pre každé $A_i$ si vytvoríme štandardnú formulu $S_i$.
    Potom štandardná forma pre $A$ je 
    $\mathscr{S}=\mathscr{S_1} \union \mathscr{S_2} 
                \union \ldots \union \mathscr{S_m}$,
    kde $\mathscr{S_i}$ je množina klauzúl štandardnej formu $S_i$.
    Formula $A$ nie je splniteľná práve vtedy, keď nie je 
    splniteľná množina klauzúl $\mathscr{S}$.
\end{poznamka}

\begin{priklad}[Z teórie grúp]
    Uvedieme príklad, kde sa pokúsime zapísať pomocou množiny
    klauzúl tvrdenie z teórie grúp. 

    \par Uvažujme grupu a na nej operáciu $\circ$.
    Platia nasledovné axiomy:
    \begin{itemize}
    %%% {{{
        \item [$A_1:$] $x, y \in G \Rightarrow x \circ y \in G$ --
                vlastnosť uzavretosti.
        \item [$A_2:$] $x, y, z \in G \Rightarrow 
                        (x \circ y) \circ z = x \circ (y \circ z)$ --
                asociatívnosť operácie.
        \item [$A_3:$] $x \in G \Rightarrow x \circ e = e \circ x = x$ --
                existencia neutrálneho prvku.
        \item [$A_4:$] $x \in G \Rightarrow \exists x^{-1} \in G :
                            x \circ x^{-1} = e = x^{-1} \circ x$ --
                existencia inverzného prvku.
    %%% }}}
    \end{itemize}

    Chceme ukázať, že ak pre všetky $x \in G$ platí $x \circ x \in G$,
    potom je grupa komutatívna. Prvou časťou úlohy je formalizovanie zadania:
    Nech $P(x,y,z)$ označuje $x \circ y = z$ a nech  $i(x) =
    x^{-1}$.\footnote{ % o tom, ze to ma Toman zle
        %%% {{{
        Na tomto mieste by som poznamenal, že to má Toman zle.
        Totiž, nahradením funkcie predikátom sme stratili jednoznačnosť
        funkčnej hodnoty, teda, momentálne môže platiť
        $P(x,y,z) \land P(x,y,z') \land z \ne z'$. Jednoduchou
        interpretáciou, ktorá spĺňa všetky axiómy je napríklad
        $P \equiv true$ a evidentne takáto realizácia nie je teóriou grúp.
        Riešenie je buď pridať piatu axiómu obmedzujúcu $P$ alebo naozaj
        používať funkčný symbol namiesto predikátového. Každopádne, zvyšok
        tohoto príkladu bude podľa Tomana.
        %%% }}}
    }

    Prepísané axiómy vyzerajú nasledovne:
    \begin{align*}
    %%% {{{
        A_1':\ &(\forall x) (\forall y) (\exists z) P(x,y,z) \\
        A_2': 
            \begin{split}
                & \phantom{\land (} (\forall x) (\forall y) (\forall z)
                (\forall u) (\forall v) (\forall w)
                [ (P(x,y,u) \land  P(y,z,v) \land P(u, z, w))
                    \implies P(x,v,w)] \\ 
                &\land (\forall x) (\forall y)(\forall z)
                (\forall u)(\forall v) (\forall w)
                [(P(x,y,u) \land P(y,z,v) \land P(x,v,w)) \implies P(u,z,w)]
            \end{split} \\
        A_3':\ &(\forall x) P(x,e,x) \land (\forall x) P(e,x,x) \\
        A_4':\ &(\forall x) P(x,i(x),e) \land (\forall x) P(i(x),x,e)
    %%% }}}
    \end{align*}

    Späť k nášmu tvrdeniu. Chceme ukázať
    \begin{align*}
    %%% {{{
        &(\forall x \in G : x \circ x \in G) \then
            (\forall u,v \in G: u \circ v = v \circ u) \\
        B':\ & (\forall x) P(x,x,e) \implies  
            (\forall u) (\forall v) [P(u,v,w) \implies P(v,u,w)]
    %%% }}}
    \end{align*}

    Teraz potrebujeme ukázať, že
    $A_1' \land A_2' \land A_3' \land  A_4' \implies B'$.
    Vytvoríme negáciu výroku a ukážeme, že nie je splniteľná:
    \begin{align*}
        F':\ & \neg (A_1' \land A_2' \land A_3' \land A_4') \lor B' \\
      \neg F':\ & A_1' \land A_2' \land A_3' \land A_4' \land \neg B'
    \end{align*}

    Ideme vytvárať množinu klauzúl. Prevedením $A_1', A_2', A_3', A_4'$ do
    štandardnej formy dostávame
    \begin{align*}
    %%% {{{
        \mathscr{S}_1:\ & \Big\{ P(x,y,f(x,y)) \Big\} \\
        \mathscr{S}_2:\ 
            \begin{split}
                 &\Big\{ \neg P(x,y,u) \lor \neg P(y,z,v)
                    \lor \neg P(u,z,w) \lor P(x,v,w), \\
                 &\ \ \neg P(x,y,u)\lor \neg P(y,z,v) \lor \neg P(x,v,q)
                    \lor P(u,z,w) \Big\}
            \end{split} \\
        \mathscr{S}_3:\ & \Big\{ P(x,e,x), \ P(e,x,x) \Big\} \\
        \mathscr{S}_4:\ & \Big\{ P(x,i(x),e), \ P(i(x),x,e) \Big\}
    %%% }}}
    \end{align*}

    Ostáva nám vyjadriť $\neg B'$ ako množinu $\mathscr{B}$.
    \begin{align*}
    %%% {{{
        \neg B':\ &
            \neg [(\forall x) P(x,x,e) \implies
                    (\forall u) (\forall v)(\forall w)(P(u,v,w) \implies
                    P(v,u,w))] \iff \\
        & \neg [ \neg (\forall x) (P(x,x,e)
                \lor (\forall u) (\forall v) (\forall w)
                    (\neg P(u,v,w) \lor P(v,u,w))] \iff \\
        &(\forall x) (P(x,x,e) \land \neg (\forall u)(\forall v) (\forall w) 
          [\neg P(u,v,w) \lor (P(v,u,w)] \iff \\
        & (\forall x) P(x,x,e)
            \land (\exists u) (\exists v) (\exists w)
                    [P(u,v,w)\land \neg P(v,u,w)]
    %%% }}}
    \end{align*}

    Množina $\mathscr{B}$ je potom
    \begin{equation*}
        \mathscr{B}:\  \Big\{P(x,x,e),\ P(a,b,c),\ \neg P(b,a,c) \Big\}
    \end{equation*}

    Máme dokázať, že množina klauzúl 
    $\mathscr{S} = \mathscr{S}_1 \union \mathscr{S}_2 \union
        \mathscr{S}_3 \union \mathscr{S}_4 \union \mathscr{B}$
    nie je splniteľná. Teda pre každú oblasť a interpretáciu je
    nesplniteľná.
    Bolo by teda potrebné vyšetriť všetky oblasti.
    Toto je ale príliš zložité -- overiť niečo na všetkých možných
    oblastiach je nad naše sily.
    Nám ale bude stačiť vyšetrovať na \emph{Herbrandovskom univerze},
    ako si ukážeme v nasledujúcej stati.
\end{priklad}

\section{Herbrandovské univerzum}

\begin{poznamka}
    V nasledujúcom texte budeme používať a medzi sebou zamieňať
    nasledujúce výrazy s rovnakým významom:
    ``nie je splniteľná'', 
    ``je sporná'' a
    ``je protirečivá''. Taktiež, občas budeme zamieňať ich negácie
    medzi sebou.
\end{poznamka}

\begin{definicia}[Herbrandovské univerzum množiny klauzúl]
    Nech $H_0$ je množina konštánt, ktoré sa vyskytujú v $S$.
    Ak $S$ neobsahuje žiadnu konštantu, tak položíme $H_0=\{ a \}$,
    kde $a$ je nejaká konštanta.
    Ďalej definujme $H_{i+1}$ ako zjednotenie $H_{i}$ a množiny všetkých termov
    tvaru $f^{(n)}(t_1,\dots, t_n)$, kde $f^{(n)} \in S$ a
    $t_1, \dots, t_n \in H_i$.
    Množinu $H_i$ nazývame
    \emph{herbrandovské univerzum $i$-tej úrovne}.
    Herbrandovské univerzum množiny klauzúl definujeme ako zjednotenie
    cez všetky úrovne:
    \begin{equation*}
        H=\bigcup_{i=0}^{\infty} H_i
    \end{equation*}
\end{definicia}

\begin{priklad}
    Majme množinu klauzúl
    $S= \{ P(x),\ \neg P(x) \lor \neg P(f(x))\}$.
    Potom herbrandovské univerzá jednotilivých úrovní sú
    \begin{align*}
        H_0& = \{ a \} \\
        H_1& = \{ a,\ f(a) \} \\
        H_2& = \{ a,\ f(a),\ f(f(a)) \} \\
         \vdots\ \ &\\
        H_{\phantom{0}} &= \{ a,\ f(a),\ f(f(a),\ f(f(f(a))),\ \ldots \}
    \end{align*}
\end{priklad}

\begin{priklad}
    Nech $S=\{P(x) \lor R(x),\ R(z), T(y) \lor \neg W(y) \}$, teda
    množina $S$ neobsahuje žiadnu konštantu.
    Preto kladieme $H_0 = \{ a \}$. Dostávame
    $H_0=H_1=\dots=H=\{a\}$.
\end{priklad}

\begin{priklad}
    Uvažujme množinu klauzúl $S=\{P(f(x),a g(y), b)\}$. Potom
    \begin{align*}
        H_0 &= \{ a,\ b\} \\
        H_1 &= \{ a,\ b,\ f(a),\ g(a),\ f(b),\ g(b) \} \\
        H_2 &= \{ a,\ b,\ f(a),\ g(a),\ f(b),\ g(b),\
                f(f(a)),\ f(g(a)),\ f(f(b)),\ f(g(b)),\
                g(f(a)),\ g(f(b)),\ g(g(a)),\ g(g(b)) \}
    \end{align*}
\end{priklad}

\begin{definicia}[výraz]
    Pod pojmom výraz budeme chápať 
    term, množinu termov,
    klauzulu, množinu klauzúl,
    atóm, množinu atómov,
    literál, množinu literálov.
\end{definicia}

\begin{definicia}[podvýraz]
    Podvýraz výrazu $F$ je ľubovoľný výraz, ktorý sa vyskytuje v $F$.
\end{definicia}

\begin{definicia}[základný výraz]
    Ľubovoľný výraz, ktorý neobsahuje premenné, sa nazýva základný výraz
    (základný term, základný atóm, základná klauzula, základný literál, ...).
\end{definicia}

\begin{definicia}[základná inštancia]
    Základnou inštanciou klauzuly $C$ z množiny klauzúl $S$ je klauzula,
    ktorú dostaneme zámenou všetkých premenných
    za prvky Herbrandovského univerza.
\end{definicia}

\begin{definicia}[Herbrandovská báza]
    Nech $S$ je množina formúl. Potom množina atómov
    tvaru $p^{(n)}(t_1, \ldots, t_n)$ pre všetky $n$-árne predikáty,
    ktoré sa vyskytujú v $S$ a všetky termy $t_i$ z Herbrandovského univerza
    nazývame Herbrandovskou bázou $S$.
    Sú to atómy takého tvaru, že sa v nich nevyskytuje žiadna
    premenná.
\end{definicia}

\begin{priklad}
    Majme množinu $S = \{P(x),\ Q(f(x)) \lor R(y) \}$ a majme klauzulu
    $C: P(x)$. 
    Herbrandovské univerzum $H$ je
    $H = \{ a,\ f(a),\ f(f(a)),\ f(f(f(a))),\ \dots \}$.
    Základné inštancie pre $C$ sú $P(a), P(f(a)), P(f(f(a))), \dots$.
\end{priklad}

Položme si otázku, čo znamená interpretovať množinu $S$ na
Herbrandovskom univerze $H$. Musíme poznať kodnoty konštánt,
interpretáciu funkčných a predikátových symbolov. Budeme uvažovať
špeciálnu interpretáciu, takzvanú $H$-interpretáciu, pri ktorej
nebudeme interpretovať predikátové symboly (necháme si to ako keby na
neskôr).

\begin{definicia}[H-interpretácia]
    Nech $S$ je množina klauzúl. Ďalej nech $H$ je herbrandovské
    univerzum pre množinu klauzúl $S$ a $I$ je interpretácia v množine
    klauzúl $S$ nad $H$.
    Hovoríme, že interpretácia $I$ je herbrandovská interpretácia 
    (alebo tiež H-interpretácia), ak platí: 

    \begin{enumerate}
	\item Interpretácia $I$ zobrazuje všetky konštanty na samé
            seba, t.j. konštante $a \in S$ priradí tú istú konštantu $a \in H$.

	\item Ak $f^{(n)}$ je $n$-árny funkčný symbol a
            $h_1, \dots, h_n$ sú prvky herbrandovského univerza $H$.
            Potom funkciu $f$ budeme v $I$ realizovať ako funkciu,
            ktorá zobrazuje $(h_1, \dots, h_n) \in H^n$ na element
            $f^{(n)}(h_1,\dots,h_n) \in H$.
    \end{enumerate}
\end{definicia}

\begin{poznamka}
    Ako sme už spomínali, nekladieme žiadne obmedzenia 
    na interpretáciu predikátov.

    Uvažujme ako príklad množinu $A = \{ A_1, A_2, \dots \}$.
    Nech je to herbrandovská báza pre množinu klauzúl $S$.
    Herbrandovskú interpretáciu určíme tak, že $I$ zadáme ako
    $I=\{ m_1, m_2, \ldots \}$, kde
    $m_i$ bude buď $A_j$ (ak $A_j$ je pravdivé) alebo
    $\neg A_j$ (ak $A_j$ je nepravdivé).
\end{poznamka}

\begin{priklad}
    Nech $S=\{ P(x) \lor Q(x),\ R(f(y)) \}$.
    Herbrandovské univerzum je
    $H=\{a,\ f(a),\ f(f(a)),\ \ldots \}$.
    V $S$ sa vyskytujú unárne predikáty $P,Q,R$.
    Herbrandovská báza je potom
    \begin{equation*}
        A: \{ P(a),\ Q(a),\ R(a),\ 
            P(f(a)),\ Q(f(a)),\ R(f(a)),\ \dots\}.\footnote{
            Všimnime si, že v danej množine je aj $R(a)$, hoci
            by sme možno očakávali, že to musí začínať $R(f(a))$
            }
    \end{equation*}
    Môžeme mať nasledovné interpretácie:
    \begin{equation*}
        I_1: \{ P(a),\ Q(a),\ R(a),\ 
        P(f(a)),\ Q(f(a)),\ R(f(a)),\ \dots \}
    \end{equation*}
    teda, predikáty sú vždy pravdivé. Alebo
    \begin{equation*}
        I_2: \{ \neg P(a),\ \neg Q(a),\ \neg R(a),\ 
            \neg P(f(a)),\ \neg Q(f(a)),\ \neg R(f(a)),\ \dots \}
    \end{equation*}
    čiže predikáty nie sú nikdy pravdivé.
    Ďalšia možná interpretácia je
    \begin{equation*}
        I_3: \{ P(a),\ Q(a),\ \neg R(a),\ 
            P(f(a)),\ Q(f(a)),\ \neg R(f(a)),\ \dots \}
    \end{equation*}
    V zásade pre každú možnú kombináciu si vieme vytvoriť
    interpretáciu.
\end{priklad}


\medskip
Ďalšou úlohou bude k ľubovoľnej interpretácii nad ľubovoľným univerzom
priradiť Herbrandovskú interpretáciu a zaviesť príslušné tvrdenia a definície.

\begin{priklad}
    \label{prikl:interpretacia}
    Majme množinu klauzúl $S = \{ P(x),\ Q(y, f(y,a)) \}$.
    Nech je oblasť interpretácie $D = \{1,2\}$ a interpretujme
    konštanty, funkčné symboly a predikátové symboly podľa tabuľky
    \ref{tab:priklad-interpretacia}

    \begin{table}[h]
        %%% {{{
        \centering
        \begin{tabular}{|r||c|c|c|c|c|}
            \hline
            symbol & a & f(1,1) & f(1,2) & f(2,1) & f(2,2) \\
            \hline
            hodnota & 2 & 1 & 2 & 2 & 1 \\
            \hline
        \end{tabular}
        
        \medskip
        \begin{tabular}{|r||c|c|c|c|c|c|}
            \hline
            predikát & P(1) & P(2) & Q(1,1) & Q(1,2) & Q(2,1) & Q(2,2) \\
            \hline
            hodnota & true & false & false & true & false & true \\
            \hline
        \end{tabular}
        \caption{Interpretácia $I$ z príkladu
          \ref{prikl:interpretacia}}
        \label{tab:priklad-interpretacia}
        %%% }}}
    \end{table}

    Ideme určiť H-interpretáciu $I^*$, ktorá akýmsi spôsobom bude
    prislúchať našej interpretácii $I$. Najsôr si určíme bázu:
    \begin{equation*}
        A=\{ P(a),\ Q(a,a),\ P(f(a,a)),\ Q(a,f(a,a)),
            Q(f(a,a), a),\ \ldots \}
    \end{equation*}

    Hodnoty pre príslušné predikáty $I^*$ určíme pomocou zadaných tabuliek
    pre $I$:

    \begin{align*}
        P(a) &= P(2) = false \\
        Q(a,a) &= Q(2,2) = true\\
        P(f(a,a)) &= P(f(2,2)) = P(1) =true \\
        Q(f(a,a),a) &= Q(f(2,2),2) = Q(1,2) =true \\
        Q(a,f(a,a)) &= Q(1,f(2,2)) = Q(2,1) =false \\
        Q(f(a,a),f(a,a)) &= Q(f(2,2), f(2,2)) = Q(1,1) =false \\
    \end{align*}

    Interpretácia $I^*$ je teda
    \begin{equation*}
        I^* = \{ \neg P(a),\ Q(a,a),\ P(f(a,a)),\ 
            \neg Q(a,f(a,a)),\ Q(f(a,a),a),\ \ldots \}
    \end{equation*}
\end{priklad}

\begin{priklad}
    \label{prikl:interpretacia2}
    Môže nastať aj iný prípad -- majme množinu klauzúl $S$, ktorá neobsahuje
    konštantu: $S=\{P(x),\ Q(y,f(y,z)) \}$.
    Máme danú interpretáciu $I$ s oblasťou $D=\{1,2\}$ popísanú
    tabuľkou \ref{tab:prikl-interpretacia2}

    \begin{table}[h]
        %%% {{{
        \centering
        \begin{tabular}{|r||c|c|c|c|}
            \hline
            symbol & f(1,1) & f(1,2) & f(2,1) & f(2,2) \\
            \hline
            hodnota  & 1 & 2 & 2 & 1 \\
            \hline
        \end{tabular}
        
        \medskip
        \begin{tabular}{|r||c|c|c|c|c|c|}
            \hline
            predikát & P(1) & P(2) & Q(1,1) & Q(1,2) & Q(2,1) & Q(2,2) \\
            \hline
            hodnota & true & false & false & true & false & false \\
            \hline
        \end{tabular}
        \caption{Interpretácia $I$ z príkladu
          \ref{prikl:interpretacia2}}
        \label{tab:prikl-interpretacia2}
        %%% }}}
    \end{table}
   
    Tejto interpretácii budú zodpovedať dve H-interpretácie
    $I_1^*$ a $I_2^*$, pričom v prvej bude $a$ interpretované ako $1$ 
    a v druhej ako $2$.

    \begin{align*}
        I_1^* &= \{ P(a),\ \neg Q(a,a),\ P(f(a,a)),\ \neg Q(a,f(a,a)),\
            \neg Q(f(a,a),a),\ \dots \} \\
        I_2^* &= \{ \neg P(a),\ \neg Q(a,a),\ P(f(a,a)),\ \neg Q(a,f(a,a)),\
            Q(f(a,a),a),\ \dots \}
    \end{align*}
\end{priklad}

Teraz si formálne zavedieme, čo to znamená priradiť nejakej interpretácii $I$
zodpovedajúcu H-interpretáciu $I^*$:

\begin{definicia}[zodpovedajúca H-interpretácia]
    Majme interpretáciu $I$ pre množinu klauzúl $S$ na oblasti
    interpretácií $D$. Interpretácii $I$ priradíme zodpovedajúcu
    H-interpretáciu $I^*$ nasledovne:
    Nech $h_1, h_2, \ldots, h_n$ sú prvky Herbrandovského univerza
    $H$.
    Nech každé $h_i$ sa zobrazí v interpretácii $I$ na prvok
    $d_i \in D$.
    Zoberme predikát $P^{(n)}$ z množiny klauzúl $S$.
    Definujeme splniteľnosť (a pravdivosť) predikátu $P$ v $I^*$
    nasledovne:
    $P^{(n)}(h_1,\dots,h_n)$ je pravdivý v $I^*$ práve vtedy keď je
    $P^{(n)}(d_1,\dots,d_n)$ pravdivý v $I$.
\end{definicia}

\begin{lema}
    \label{lema:h-interpretacia}
    Majme interpretáciu $I$ na oblasti $D$. Nech táto interpretácia
    vyhovuje množine klauzúl $S$. Potom ľubovoľná H-interpretácia $I^*$,
    ktorá je priradená (zodpovedá) $I$, taktiež vyhovuje množine klauzúl $S$.
\end{lema}

\begin{dokaz}
    Majme množinu klauzúl $S=\{C^{(1)}, C^{(2)}, \ldots , C^{(n)}\}$.
    Nech klauzula $C^{(i)}$ je nasledujúca disjunkcia literálov
    $C^{(i)} = L^{(i)}_1 \lor L^{(i)}_2 \lor \ldots \lor
        L^{(i)}_{r_i}$ pre $i=1,\ldots,n$.
    Ľubovoľný literál $L^{(i)}_j$ je tvaru atomická formula alebo
    jej negácia. Predpokladajme že interpretácia $I$ na univeze $D$
    vyhovuje množine klauzúl $S$. Potom musí byť pravdivý aspoň jeden
    literál v každej klauzule $C^{(i)}$.
    Lenže pre tieto literály
    $L^{(i)}_j$ platí
    $L^{(i)}_j = P^{(n)}(d_1,\dots,d_n)$ je pravdivé v $I$ a z toho
    dostávame $P^{(n)}(h_1, \dots, h_n)$ je pravdivé v $I^*$.
    Teda, ak $I$ vyhovuje množine klauzúl $S$, bude jej vyhovovať aj
    $I^*$.
\end{dokaz}

\begin{veta}
    Množina klauzúl $S$ nie je splniteľná práve vtedy,
    keď $S$ je nepravdivá pri všetkých H-interpretáciach $S$.
\end{veta}


\begin{dokaz}
    \noindent
    \begin{itemize}
    \item[$\Rightarrow:$] Nech množina klauzúl $S$ nie je splniteľná.
        Potom je nepravdivá pre ľubovoľnú interpretáciu na ľubovoľnej
        oblasti.
        Teda aj pre ľubovoľnú H-interpretáciu na H-univerze.

    \item[$\Leftarrow:$] Nech množina klauzúl $S$ je nepravdivá pre
        ľubovoľnú H-interpretáciu množiny klauzúl $S$.
        Pre spor predpokladajme, že existuje interpretácia $I$ 
        s oblasťou $D$ pre množinu klauzúl $S$, ktorá
        vyhovuje $S$.
        Uvažujme $I^*$, ktorá je priradená (zodpovedá) interpretácii
        $I$ pre množinu klauzúl $S$.
        Podľa lemy \ref{lema:h-interpretacia} ale vieme,
        že ak vyhovuje $I$, vyhovuje aj $I^*$, čo je spor.
    \end{itemize}
\end{dokaz}

Podarilo sa nám teda objaviť také univerzum (H-univerzum), pre ktoré
stačí vyšetriť splniteľnosť množiny klauzúl $S$ na všetkých
interpretáciach a budeme vedieť povedať (ne)splniteľnosť formuly na
ľubovoľnej interpretácii v ľubovoľnom univerze. Posunuli sme sa o krok
bližšie k rozhodovaniu platnosti $S$.

\begin{poznamka}
    V nasledujúcom texte budeme používať iba H-interpretácie.
    Preto sa dohodneme (na skrátenie a zjednodušenie zápisu), že
    ich budeme nazývať iba interpretácie a označovať ako $I$.
\end{poznamka}

\begin{poznamka}
    \noindent
    \begin{enumerate}
	\item Majme klauzulu (disjunkciu literálov) $C$
            a nech $C'$ je jej základná inštancia, t.j. 
            každá premenná bola nahradená prvkom H-univerza.
            $C'$ je splniteľná v interpretácii $I$
	    práve vtedy, keď existuje základný literál $L' \in I$,
            ktorý je splniteľný.
            Teda $C'$ je splniteľná $\iff$ $C' \intersect I \ne 0$.            

	\item Klauzula $C$ je splnená v interpretácii $I$ $\iff$ každá jej
	základná inštancia $C'$ je splnená v $I$.

	\item Klauzula $C$ je odmietnutá (vyvrátená) v $I$ $\iff$ existuje
	aspoň jedna taká základná inštancia $C'$, ktorá 
        nie je splniteľná (teda je vyvrátená) v $I$.

	\item Množina klauzúl $S$ nie je splniteľná $\iff$ 
        pre každú interpretáciu $I$ existuje aspoň jedna základná inštancia 
        $C'$ klauzuly $C\in S$, ktorá nie je splniteľná v $I$.
    \end{enumerate}
\end{poznamka}

\begin{priklad}
    Uvažujme klauzulu $C=\neg P(x) \lor Q(f(x))$ a
    interpretácie $I_1,I_2,I_3$ definované nasledovne:
    \begin{align*}
        I_1 &= \{ \neg P(a),\ \neg Q(a),\ \neg P(f(a)),\ 
            \neg Q(f(a)),\ \neg P(f(f(a))),\ \neg Q(f(f(a))),\ \ldots \} \\
        I_2 &= \{ P(a),\ Q(a),\ P(f(a)),\ Q(f(a)),\ P(f(f(a))),\ 
            Q(f(f(a))),\ \ldots \} \\
        I_3 &= \{ P(a),\ \neg Q(a),\ P(f(a)),\ \neg Q(f(a)),\ 
            P(f(f(a))),\ \neg Q(f(f(a))),\ \ldots \} \\
    \end{align*}
    Všímajme si klauzulu $C$ a jednotlivé interpretácie:
    $I_1$ vyhovuje $C$ (inak povedané $C$ je splnená) a zabezpečujú to
    $\neg P(a),\ \neg P(f(a)),\ \neg P(f(f(a))), \dots$.
    Podobne, $C$ je splnená v $I_2$ a zabezpečujú to
    $Q(a),\ Q(f(a)),\ Q(f(f(a))), \dots$. Naopak, $C$ nie je
    splniteľná v $I_3$.
\end{priklad}

\begin{priklad}
    Uvažujme množinu klauzúl $S=\{P(x),\ \neg P(x)\}$ a
    interpretácie $I_1 = \{ P(x) \}, I_2 = \{ \neg P(x) \}$.

    Množina $S$ nie je splnená ani jednou interpretáciou.
\end{priklad}


\chapter{Dokončenie úvodu do matematickej logiky}

\todo{prenexne, skolemov tvar}

\chapter{Matematická logika}

\section{Pravdivosť a dokázateľnosť}
\subsection{Veta o úplnosti}

\paragraph{Definícia:} Nech $L$ je jazyk prvého rádu a $A$ je formula jazyka
$L$. Hovoríme, že formula $A$ je \emph{logicky platná}, označujeme $\models A$,
ak je splnená v ľubovoľnej realizácii jazyka $L$.
$$
\begin{array}{lll}
m & \models & A[e] \\
  & \models & A \\
 ~ & \models & A \\
\end{array}
$$

\paragraph{Definícia:} Nech $L$ je jazyk, prvého rádu a $T$ je množina formúl
jazyka $L$. Hovoríme, že $T$ je teória 1. rádu predikátovej logiky s jazykom
$L$ (množina formúl $T$ je množina axióm teórie).

\paragraph{Definícia:} Nech $T$ je teória v jazyku $L$, $m$ je realizácia jazyka
$L$. Hovoríme, že $m$ je modelom teórie $T$, ak pre každú formulu $A$ patriacu
$T$ platí $\models A$.
\par Hovoríme, že formula $A$ je sémantickým dôsledkom (vetou teórie) množiny formúl $T$,
označujeme $T \models A$, ak $A$ je splnená v každom modeli teórie $T$.

\paragraph{Príklad:}
$\varphi, x, y, z, <$
\begin{enumerate}
	\item $(x,y) \in \varphi  \implies (y,x) \notin  \varphi$
	\item $(x,y) \in \varphi \land (y,z) \in \varphi \implies (x,z) \in \varphi$
	\item $x \neq y \implies (x,y) \in \varphi \lor (y,x) \in \varphi$
\end{enumerate}

je modelom pre teóriu ostrého usporiadania.
Ďalšie príklady: neostré usporiadanie: antisymetrickosť, tranzitívnosť,
reflexívnosť, trichotomockosť.


\paragraph{Príklad:} Teória grúp.

\begin{enumerate}
	\item $((x+y)+z) = (x+(y+z))$
	\item $(x+0) = (0+x) = x$
	\item $x, -x, x+(-x) = 0 = (-x)+x$
\end{enumerate}


\paragraph{Príklad:} Elementárna aritmetika.
$$+, \cdot, 0, S(x) = x+1, N^+ $$
\begin{enumerate}
	\item $\neg S(x) = 0$
	\item $S(x) = S(y) \implies x=y$
	\item $x+0 = x$
	\item $x+S(y) = S(x+y)$
	\item $x \cdot 0 = 0$
	\item $x\cdot S(y) = ((x\cdot y) + x)$
\end{enumerate}

Pätica v zátvorkách je \emph{štandardný model} elementárnej aritmetiky:
$$<N^+, +, \cdot, S, 0>$$

\par
Pojem platnej formuly je prirodzené zovšobecnenie pojmu \emph{logicky platnej
formuly}.

Ďalším cieľom je stotožniť dokázateľné formuly s tautológiami.

\paragraph{Veta (o korektnosti):} Ak $T$ je teória v jazyku $L$ a ak formula $A$
je taká, že $T \vdash A$, potom $T \models A$.
\paragraph{Dôkaz:} Nech $A_1, A_2, \ldots A_n$ je odvodenie (dókaz) formuly $A$
z predpokladov $T$ (v teórii $T$). $A_1, A_2, \ldots A_n(A)$.
\par $m$ nech je ľubovoľný model teórie $T$. Ukážeme (indukciou podľa dĺžky
dôkazu), že platí $m \models A_i$ pri predpoklade, že pre $A_j, j < i$ platí
$m \models A_j$.i

Do dôkazu sa $A_i$ môže dostať niekoľkými spôsobmi:
\begin{enumerate}
	\item $A_i \in T$, m -- model $T$, potom $m \models A_i$
	\item $A_i$ sa dostane do dôkazu ako schéma axiom predikátovej logiky:
	\begin{enumerate}
		\item $A_i$ je axioma výrokovej logiky -- je poskladaná z
		atomických formúl a logických spojok $\neg$ a $\implies$. Potom
		$A_i$ je tautológia výrokovej logiky (ak formula je tautoĺógia,
		jej pravdivostná hodnota nezávisi od ohodnotenia premenných, a
		$m \models A_i$, teda $m \models A_i[e]$).
		
		\item $A_i$ je axioma špecifikácie, teda je tvaru $A_i: (\forall
		x) B \implies B_x[t]$, $t$ je substitúcia za $x$ do $B$. $m
		\models A_i$.
		\par
		Zaujíma nás prípad, kedy $(\forall x) B$ je pravdivý. To
		znamená, že pre ľubovoľné ohodnotenie platí $B[e(x/m)]$, teda
		$e(x)=m$. Tvrdenie zo zimného semestra: $$m \models A_{x_1,
		\ldots x_n}[t_1, \ldots t_n][e] \Leftrightarrow m \models
		A[e(x_1/m_1, \ldots x_n/m_n)]$$ a $$t_i[e] = m_i.$$ 

		Ak položím tvrdenie $B_x[t][e]$, zmení sa to na $B[e(x/m)]$ a
		táto formula je pravdivá v $m \models B[e(x/m)]$.

		\item $A_i: (\forall x) (B \implies C) \implies (B \implies
		(\forall x) C)$ a $x$ nie je voľná v $B$. Mali by sme dokázať,
		že platí $m \models A_i$. Formula  $m \models (B \implies
		C)[e(x/m)]$ platí, pozeráme sa na $(B \implies C)$. To je
		ekvivalentné s $\neg B \lor C$. Dôležitý je tiež predpoklad, že
		$x$ nie je voľná v $B$, a teda nezávisí od ohodnotenia viazanej
		premennej. Keď je ľavá časť axiomy pravdivá, musí byť aj pravá
		(TODO: spresniť!) $m \models A_i$.
	\end{enumerate}
	\item $A_i$ je niektorá axioma rovnosti:
	\begin{enumerate}
		\item $A_i: x=x$, potom $m\models x=x$, $A_i[e]$ a $m=m$.
		\item $A_i: x_1 = y_1 \implies x_2 = y_2 \implies \ldots \implies
		f(x_1, \ldots, x_n) = f(y_1, \ldots, y_n).$. Zaujíma nás prípad,
		keď  $e(x_i/m_i)$ a $e(y_i/m_i)$. Získavame potom zobrazenie na
		$f_m(m_1, \ldots m_n)$ a rovnosť platí.
		\item $A_i: x_1 = y_1 \implies x_2 = y_2 \implies \ldots \implies
		P(x_1, \ldots, x_n) = P(y_1, \ldots, y_n).$. Zaujíma nás opäť
		prípad, keď $x_i$ aj $y_i$ reprezentujeme rovnakým indivíduom:
		$e(x_i/m_i)$ a $e(y_i/m_i)$.  $P_m(m_1, \ldots, m_n)$, $(m_1,
		\ldots m_n) \in P$ -- buď je v relácii, alebo nie je.
	\end{enumerate}
	\item Odvodzovacie pravidlá, $A_i: A_{j_1}, A_{j_2}$, pričom $j_1 < i$ a
	$j_2 < i$.
	\begin{enumerate}
		\item $A_{j_2}. A_{j_1} \implies A_j$, $A_{j_1}$ a $m \models
		A_{j_1}$, $m \models A_{j_2}$, potom $m \models A_i$.
		\item $A_i: \forall A_j$, potom $m \models A_j[e(x/m)]$ a $m
		\models A_j$.
	\end{enumerate}

\end{enumerate}

\paragraph{Pokračovanie vety o korektnosti:}

\paragraph{Príklad:} Uvažujme elementárnu aritmetiku, ktorá má svoj štandardný
model. Uvažujme formulu $x=0$ v $N$. 

$$
\begin{array}{l}
	e(x) = m \neq 0 \\
	A: x=0 \\
	m \notmodels A[e] \\
\end{array}
$$

$$
\begin{array}{l}
	A: \neg x=0 \\
	e(x)=0 \\
	m \notmodels A[e] \\
\end{array}
$$

Formula ani jej negácia nie sú dokázateľné v elementárnej aritmetike.


\paragraph{Poznámka:} Vetu o dedukcii v predikátovej logike nemožno vysloviť pre
ľubovoľnú formulu.

$$ 
\begin{array}{ll}
	A: & \neg x=0 \\
	A: & \neg y=0 \\
\end{array}
$$

$$
	\neg x=0 \vdash \neg y=0
$$

Môžem podľa vety o dedukcii napísať, že $ \vdash \neg x=0 \implies y=0$?

$$
\begin{array}{l}
	A \implies B \\
	e:	\\
	e(x)=0	\\
	e(x)= m \neq 0	\\
	m \notmodels (A \implies B)[e]	\\
\end{array}
$$

\paragraph{Dôsledok:} Ak teória $T$ v jazyku $L$ má model $m$, potom $T$ je
bezposporná.
\paragraph{Dôkaz:} Nech $m$ je model $T$ a $A$ je ľubovoľná uzavretá formula
jazyka $L$. Potom práve jedna z formúl $A$, $\neg A$ je pravdivá v modeli $m$.
Tá, ktorá nie je pravdivá, nie je ani dokázateľná (podľa vety o korektnosti).

\paragraph{}
\par Tento výsledok hovorí, ze ak máme vyšetriť bezospornosť nejakej teórie,
treba nájsť jej model. Keď si zoberieme predikátovú logiku:

\par Každá realizácia jazyka predikátovej logiky je modelom (neskôr nám bude
stačiť zobrať model, ktorý má jedno jediné indivíduum).  

\par
Syntaktický prístup ku dôkazu bezospornosti: Teóriu, ktorej bezospornosť chceme
dokázať, transformuje pomocou istých pravidiel do teórie, o ktorej vieme, že je
bezsosporná  (my budeme predikátovú logiku transformovať do výrokovej logiky).

\par $L$ -- jazyk prvého rádu, $c$ -- konštanta taká, že $c \notin L$. $L'$
vznikne z $L$ pridaním $c$. Ľubovoľný term nahradíme konštantou $c$, z formuly
vynechávame všetky kvantifikátory, vynechávame premenné bezprostredne spojené s
kvantifikátorom.

\par Každej formuli $A$ na jazyku $L$ priradíme formulu $A^*$ na jazyku $L'$:
\begin{enumerate}
	\item ak A je tvaru $A: P(t_1, \cdots t_n)$, tak $A^*: P(c, c, \cdots,c)$.
	\item $A: B \square C$, potom $A^*: B^* \square C^*$.
	\item $A: \neg B$, potom $A^*: \neg B^*$
	\item $A: (Qx) B$, potom $A^*: B^*$.
\end{enumerate}

Ak $L$ je jazyk bez rovnosti a $\vdash^L A$, potom $A^*$ je tautológia. Ak $L$
je jazyk s rovnosťou a $\vdash^L A$, potom $A^*$ je tautologický dôsledok $c=c$.

\par Tvrdíme, že pre žiadnu formulu $A$ nie je $\vdash A$ aj $\vdash \neg A$. Ak
by to platilo, dostali by sme sa do sporu, že vo výrokovej logike je $\vdash
A^*$ aj $\vdash \neg A^*$.


\section{Veta o úplnosti}
\paragraph{Veta (G\"odel, 1):} Nech $T$ je teória v jazyku $L$ a ak $A$ je
ľubovoľná formula jazyka $L$, potom $T \vdash A \iff T \models A$.

\paragraph{Veta (G\"odel, 2)} Teória $T$ je bezosporná práve vtedy, ked $T$ má
model.

\paragraph {Veta (o korektnosti):}
	$$T \vdash A \implies T \models A$$.

\paragraph{Príklady vety o korektnosti:}
Ak teória $T$ má model, tak je bezosporná.

\paragraph{Dôkaz:}
$\mbox{Veta 2} \implies \mbox{Veta 1}$. Veta o dedukcii mala nasledovný
dôsledok:
\par Ak máme teóriu $T$, $A$ je jej formula a $A'$ je uzáver formuly $A$, tak $T
\vdash A \iff T \union \{ \neg A' \} \mbox{ je sporná teória}$, podľa vety 2 $T
\union \{ \neg A' \}$ nemá model.

\par Toto znamená, že v každom modeli teórie $T$ je pravdivý uzáver $A'$. $T
\models A' \Rightarrow T \models A$. Z platnosti vety 2 vyplýva veta 1.

\par Budeme sa snažiť zostrojiť model pre teóriu, ktorá je bezosporná. Majme
bezospornú teóriu v jazyku $L$. Potrebujeme relácie, zobrazenia a $M$ --
univerzum. K dispozícii máme len syntaktické prostriedky teórie. Do úvahy
prichádzajú iba termy bez premenných -- majú jednoznačnú realizáciu (sami sebe
budú realizáciou).

$$A[e] \iff T \vdash A $$.

Možné problémy pri konštrukcii modelu:
\begin{enumerate}
	\item Jazyk $L$ neobsahuje žiadne konštanty (a teda žiadne termy bez
		premenných).
	\item Ak jazyk $L$ bude jazyk s rovnosťou, môže sa stať, že v teórii $T$
		bude $T \vdash t=s$, ale $t$ a $s$ sú rôzne termy bez
		premenných.
	\item Nech $m$ je ľubovoľná realizácia jazyka $L$ a $A$ je uzavretá
		formula jazyka $L$. Potom práve jedna z formúl $A$, $\neg A$ je
		pravdivá, ale žiadna z nich nemusí byť dokázateľná v $T$.
	\item Môže sa stať, že formula $(\exists x)B$ je dokázateľná v teórii
		$T$, ale pre žiaden term $t$ bez premenných formula $B_x[t]$ nie
		je dokázateľná v $T$. To znamená, že nie je ani pravdivá a podľa
		Tarského definície pravdepodobnosti $(\exists)B$ je term $T$.
\end{enumerate}

Ako odstránime tieto nedostatky:
Odstránenie bodu 2 je jednoduché --  faktorizáciou. Body 1 a 4 riešime
rozšírením o konštanty. Tretí bod sa rieši tzv. rošírením teórie (Henkinovským).
Budú to tzv.  konzervatívne teórie (na pôvodnom jazyku nezískame žiadne nové
teorémy).

\paragraph{Definícia:} Hovoríme, že teória $T$ s jazykom $L$ je \emph{úplná}, ak
je bezosporná teória a pre ľubovoľnú uzavretú formulu $A$ na jazyku $L$ buď $A$,
buď $\neg A$ je dokázateľná v $T$.

\par Označím $T_h(m)$ všetky pravdivé uzavreté formuly, tak $T_h$ je úplná
(nemusíme brať otvorené formuly -- napr. $x=0$ v aritmetike nemusí byť pravdivé,
lebo závisí od ohodnotenia $x$).

\paragraph{Definícia:} Hovoríme, že teória $T$ s jazykom $L$ je \emph{Henkinova}, ak
pre ľubovoľnú uzavretú formulu $(\exists x)B$ jazyka $L$ platí $$ T \vdash
(\exists x)B \implies B_x[c]$$ pre nejaké $c$ -- konštantu.

\paragraph{Lema:} Ak $T$ je úplná a Henkinova teória, tak potom $T$ má model.
\paragraph{Dôkaz:} Nech $L$ je jazyk teórie $T$, $\tau$ je množina všetkých
termov jazyka $L$ bez premenných. Na množine $\tau$ definujem reláciu
ekvivalencie nasledovne:
$$
	t_1, t_2 \in \tau
	t_1 \equiv t_2 \iff T \vdash t_1 = t_2
$$
Rovnosť je reflexívna, symetrická, tranzitívna, teda týmto spôsobom definovaná
relácia je relácia ekvivalencie a rozdeľuje nám množinu $\tau$ na triedy
ekvivalencie:

$$
	[t] = \{ s \in \tau, t \equiv s\}
$$

Univerzum budú tvoriť termy bez premenných.

\paragraph{Lema:} Ak $T$ je úplná a Henkinova teória, tak potom $T$ má model.
\paragraph{Dôkaz:} $M$ bude mnomžina všetkých termov bez premenných. Na tejto
množine definujeme reláciu $\equiv$ tak, že 
$$ t_1 \equiv t_2,~t_1, t_1 \in \tau,~T \vdash t_1=t_2 $$
$$[t]=\{ s\in \tau, s \equiv t \} $$

Nech $f$ je $n$-árny funkčný symbol. $[t_1], [t_2], \ldots, [t_n] \in M$, pričom
platí $f_m([t_1], \ldots, [t_n]) = [f(t_1, \ldots t_n)]$. $P$ je $n$-árny
predikátový symbol rôzny od $=$. $([t_1], \ldots [t_n]) \in P_m \iff T \vdash
P(t_1, \ldots, t_n)$.

\paragraph{Poznámka:} Doteraz sme pracovali s $[\ldots]$. Výsledok sa nezmení,
ak zoberiem $s_1 \in [t_1], \ldots s_n \in [t_n]$. Nezávisí to od reprezentanta
danej triedy.
\par
$$ s_1, \ldots s_n $$
$$ s_i \in [t_i], i = 1, \ldots n $$
$$ [t_1], \ldots, [t_n]) \in P_m \iff T \vdash P(s_1, \dots s_n) $$

\par
$$
\begin{array}{rll}
x_1, x_2, &\ldots,& x_n \mbox{-- premenné}.	\\
e(x_i) &=& t_i, i=1,\ldots,N	\\
t[e] &=& [t_{x_1,\ldots,x_n}][t_1,\ldots, t_n]	\\
t[e] &=& [t] \\
\end{array}
$$

$A$ -- ľubovoľná uzavretá formula z $L$, potom $m \models A \iff T \vdash A$.
$P(t_1, t_2, \ldots t_n)$ -- atomická formula,  uzavretá, $t_1. \ldots t_n$ sú
termy bez premenných.

$m \models A \iff ([t_1], \ldots, [t_n]) \in P_m \iff T \vdash P(t_1,\ldots
t_n)$ pre ľubovoľné ohodnitenie.

\par $A: t_1 = t_2$. $m \models t_1 = t_2 \iff [t_1] = [t_2] \iff T \vdash t_1 =
t_2$.

\par $A: \neg B$. Indukčný predpoklad: pre $B$ bolo tvrdenie dokázané ($m\models
B \iff T \vdash B)$. Formula $A$ je pravdivá v $m$ práve vtedy, keď formula $B$
nie je pravdivá $B$ nie je pravdivá, $T \nvdash B, T \vdash \neg B$.

\par $A. B \implies C$, pričom $B$ a $Cc$ sú dokázané podľa indukčného
predpokladu. Ak platí $m \models B \implies C$, potom $m \models \neg B \lor m
\models C$. $T \vdash \neg B \lor T \vdash C$.

$$ T \vdash B \implies C $$
$$ \vdash C \implies (B \implies C)$$
Opačná implikácia:
$$T\vdash B\implies C\Rightarrow m \models B \implies C$$

Teória $T$ je úplná teória a $B \implies C$ je uzavretá. Pre každú uzavretú
formulu platí, jedna z možností $T \vdash \neg B$ (potom $m \models \neg B$)
alebo $T \vdash B$, a z toho vyplýva, že $T \vdash C$, a teda $m \models C$.
Tvrdenie je teda dokázané.

$$m \models A \iff T \vdash A$$
\par Predpokladajme, že formula $A$ je v tvare $A: (\forall x) B$. Pre každú
inštanciu formuly $B$ tvrdenie platí, $\vdash (\forall x) B \implies B_x[t]$
(axióma špecifikácie). Potom aj $T \vdash B_x[t]$ (lebo teória je Henkinova). $m
\models B[e(x/t)]$.

\par Ak je pravdivá negácia, potom $T \vdash \neg A$, čo znamená, že $(\exists
x) \neg B$. $T \vdash (\exists x) \neg B \implies B_x[c]$, a to znamená, že $T
\vdash \neg B_x[c]$.

\par Nech $A$ je ľubovoľná formula z teórie $T$. Pre ňu platí, že $T \vdash A$
(predpoklad). Potrebujem dokázať, že $T \vdash A \iff m \models A$. Zoberiem si
$A'$ -- uzáver formuly $A$. Na $A'$ sa vzťahuje tvrdenie, ktoré sme už dokázali,
a teda $m \models A' \iff T \vdash A'$. Podľa vety o uzávere:
$m \models A \iff  m\models A' T \vdash A' \iff T \vdash A$

\par Pre ľubovoľnú teóriu, ktorá je idealizovaná (Henkinova a úplná), viem
zostrojiť model.

\section{Veta o kompaktnosti}

Budeme definovať rozšírenia (ktorá ani nepomôžu, ani neuškodia :-). (Henkinova
veta, Lindenbaumova veta)

\paragraph{Definícia:} Jazyk $L'$ je \emph{rozšírením jazyka} $L$, ak každý
špeciálny symbol jazyka $L$ (prípadne aj symbol $=$) je v jazyku $L'$ obsiahnutý
s rovnakým významom a s rovnakou árnosťou.

\paragraph{Príklad:} $L'$ s rovnosťou a špeciálnymi symbolmi $<$ a $0$. Jazyk
$L$ má symboly $<$ a $=$. Jazyk $L'$ je rozšírením jazyka $L$

\paragraph{Definícia:} Majme teóriu $T'$ s jazykom $L'$. Hovoríme, že je
rozšírením teórie $T$ s jazykom $L$, ak platí, že $L'$ je rozšírením jazyka $L$
a každá formula $A$ na jazyku $L$, $T \vdash A, T' \vdash A$.
\paragraph{Veta (Henkinova):}

$$ \vdash (\exists x) \neg B \implies \neg B_x[c] $$
$$ T \vdash  \neg B_x[c] $$
$$ m \models  (\forall x) B \iff T (\forall c) B $$

\par Relačná štruktúra $<M, P_m, t_m>$ -- model pre množinu formúl $T$. Vezmeme
ľubovoľnú formulu $A \in T$ a platí $m \models A$, potom $T$ je modelom.

\par $T \vdash A$. $A'$ -- uzáver A. $$ m \models A' \iff T \vdash A'$$. Ak v
modeli $m$ je splnená formula $A'$, potom musí byť splnená aj $A$. Štruktúra,
ktorú sme uvažovali, je modelom pre našu teóriu.

\paragraph{Záver:} Pre idealizovanú teóriu viem vytvoriť model.

\par Ďalšou úlohou je rozšíriť bezospornú teóriu do úplnej a Henkinovskej.

\paragraph{Príklad (čo je to model):} Teória grúp -- je to teória s rovnosťou a
používa špeciálny symbol $+$ a existuje v ňom neutrálny prvok. Axiomy:

$$\forall x,y,z (x+y)+z = x+(y+z)$$
$$(x+0) = (0+x) = x$$
$$(x+(-x)) = ((-x)+x) = 0$$.

\par Ak si vezmeme relačnú štruktúru $m=<N^+,0,+>$, táto realizuje jazyk teórie
grúp, ale nie je modelom (nie všetky axiomy nie sú splnené). Pokiaľ ale vezmeme
$m=<Z,0,+>$, táto množina tvorí model teórie grúp.  

\par Máme daný jazyk $L$, jazyk $L'$, ktorý je rozšírením $L$ (obsahuje každý
symbol z $L$ s rovnakou hodnotou a árnosťou). Vezmime jazyk $L= \{0, 1, +,
\cdot \}$. Ak nad týmto jazykom budeme pridávať axiómy, môžeme získať teórie ako
teória okruhov, teória oborov integrity, teória telies, \ldots . Zistíme, že
teória telies je rozšírenie teórie oborov integrity a tie sú rozšírením teórie
okruhov. 

\par Ak máme teóriu $T$, $T'$ je rozšírenie $T$ a je bezosporná. Potom teória
$T$ je taktiež bezosporná. Ak by sme pripustili, že teória $T$ je sporná, potom
by sa našla taká formula, že $T \vdash A$ aj $T \vdash \neg A$, z čoho vyplýva
$T' \vdash A$ aj $T' \vdash \neg A$.

\paragraph{Definícia:} Majme teóriu $T$ a teóriu $T'$. Hovoríme, že $T'$ je
\emph{konzervatívne rozšírenie} teórie $T$, ak $T'$ je rozšírenie $T$ a pre
každú formulu $A$ na jazyku $L$, pre ktorú platí, že v $T'$ je dokázateľná tá
formula implikuje, že v $T$ je taktiež dokázateľná.

\paragraph{Tvrdenie:} $T$, $T'$ je konzervatívne rozšírenie $T$, potom $T$ je
bezosporná vtedy a len vtedy, keď $T'$ je bezosporná. Vyplýva to z toho, že ak
zoberieme konzervatívne rozšírenie, pre ľubovoľnú formulu pôvodného jazyka
platí, že:
$$ T' \vdash_L A \Rightarrow T \vdash A$$.

\paragraph{Veta (Henkinova):} K ľubovoľnej teórii $T$ môžeme zostrojiť Henkinovu
teóriu $T_H$, ktorá je konzervatívnym rozšírením teórie $T$.

\paragraph{Dôkaz:} Teóriu $T$ rozšírime a priradíme k nej teóriu $T_H$ tak, že
jednak rozšírime jazyk teórie a jednak pridáme axiómy. Nech teória $T$ má jazyk
$L$. Budeme tvoriť jazyk rozšírený jazyk $L_H$ (budeme pridávať konštanty):

\textbf{<skontroluj, malé písmená>}

\par Vezmime si formulu $A$ napísanú v jazyku $L$, ktorá má jedinú premennú $x$
priradíme jej premenné $c_A$ a $c_{\neg A}$ a nasledovné axiomy: 

$
\begin{array}{lll}
\mbox{(*)} & (\exists x) A \implies A_x[c_A] & (H1) \\
\mbox{(**)} & A_x{c_{\neg A}} \implies (\forall x) A & (H2) \\
\end{array}
$

Stačila by nám aj jedna konštanta:
$$B: \neg A$$
$$ (*) (\exists x) \neg A \implies \neg A_x[c_{\neg A}]$$
$$\neg \neg A_x[c_{\neg A}] \implies (\forall x) \neg \neg A$$
$$A_x[c_{\neg A}] \implies (\forall x) A$$

Alebo:
$$B: \neg A$$
$$\neg A_x[c_{??} \implies (\forall x) \neg A$$
$$(\exists x) A \implies A_x[c_A]$$
\textbf{</skontroluj, malé písmená>}

\par $C_1$ bude množina všetkých $c_A$ a $c_{\neg A}$. Týmito konštantami
rozšírime jazyk. Induktívne pokračujeme ďalej a podobne vytvárame množinu $C_2$
(zoberieme ľubovoľnú formulu $A$ nad rozšíreným jazykom, ktorá mala jedinú voľnú
premennú $x$ a obsahovala aspoň jednu konštantu prvého rádu).
Pridané konštanty budeme nazývať konštanty druhého rádu. Pokračujeme až po
$C_n$.

\par Vezmeme formuly, ktoré majú jedinú voľnú premennú $x$ a obsahujú nejakú
pridávanú konštantu rádu $n$.

\par
$$C_1 \union C_2 \union \ldots C_n \union C_{n+1} = L(C)$$. 
Teóriu $T_H$ bude tvoriť jazyk $L(C)$ a pridáme všetky axiomy, ktoré sme pridali
a tie, ktoré tam už boli. Je ihneď zrejmé, že $T_H$ je rozšírenie teórie $T$.
Potrebujeme ukázať, že $T_H$ konzervatívne rozšírenie teórie $T$.

$T_H$ -- konzervatívne rozšírenie, teda keď zoberiem formulu $A$ na jazyku $L$,
musí platiť:
$$ T_H \vdash A \implies T \vdash A$$

Nech platí, že v $T_H$ je dokázateľná teória, $B_1, B_2, \ldots B_n$ -- všetky
axiomy z $H$ a môžem predpokladať, že všetky sú jednoduché. Platí, že:
$$ T, B_1, \ldots B_n \vdash A $$. Môžeme použiť vetu o dedukcii:

$$ T \vdash B_1 \implies B_2 \implies \ldots B_n \implies A $$.

Môžeme ich poprehadzovať tak, že $B_1$ obsahuje koštantu maximálneho rádu. Nech
$B_1$ je tvaru $(\exists x) D \implies D_x[c_D]$, pričom $C_D$ nie je je použitá
v $B_2, \ldots, B_n$. Použijeme vetu o konštantách a dostávame, že $T \vdash
((\exists x) D \implies D_x[c_D] ) \implies (B_2 \implies \ldots \implies B_n
\implies A)$.

\par Ak $T \vdash A \implies B$, potom $T\vdash (\exists x) A \implies B$.
$$ T\vdash (\exists w) ((\exists x)D \implies D_x[w]) \implies (B_2 \implies
\ldots \implies B_n \implies A)$$

Použili sme pravidlo zavedenia existenčného kvantifikátora. Ďalej použijeme
prenexnú operáciu $(Qx) (A\implies B) \iff (A\implies (Qx)B)$ a $x$ nie je voľná
v $A$. V $T$ je dokázateľné:
$$T \vdash ((\exists x) D \implies (\exists w)D_x[w]) \implies B_2 \implies \ldots \implies B_n \implies A$$

Použijeme vetu o variantoch a dostávame:
$$ T \vdash ( \exists x ) D \implies (\exists w) D_x[w] $$
$$ T \vdash B_2 \implies B-3 \implies \ldots \implies B_n \implies A $$
$$ T \vdash A$$

Ukázali sme, že $T_H$ s jazykom $L(C)$ je konzervatívne rozšírenie jazyka $L$.

\paragraph{Veta (Lindenbaum):} Ak $T$ je teória s jazykom $L$ a $T$ je
bezosporná, potom existuje úplné rozšírenie $T'$ s jazykom $L$.

Ak teória $T$ bude bezosporná, potom aj $T_H$ bude bezosporná.

\paragraph{Dôkaz:} Nech $\mathcal{S}$ je množina všetkých uzavretých formúl jazyka $L$.
Ďalej definujeme $\mathcal{B} \{ S \mid S \subseteq \mathcal{S}, T \union S
\mbox{je bezosporná}\}$. $\mathcal{B}$ je čiastočne usporiadaná na $\subseteq$.
Množina $\mathcal{B}$ má konečnú vlastnosť -- keď zoberiem ľubovoľnú podmnožinu
$S \subseteq \mathcal{S}$, bude platiť $$ S \in \mathcal{B} \iff S' \in
\mathcal{B}, S' \mbox{ je konečná}$$


\section{Prednáška 10.3.2009}

\par\textbf{Zapísal:} P. Struhár

\par Dostali sme sa do toho  kroku, že máme bezospornú teóriu, ktorú je treba zúplniť.
To sme urobili \emph{Lindenbaumovou vetou\footnote{bolo už na predchádzajúcej prednáške}}.

\paragraph{Veta (Lindenbaum):} Ak $T$ je teória s jazykom $L$ a $T$ je
bezosporná, potom existuje úplné rozšírenie $T'$ s jazykom $L$.

Ak teória $T$ bude bezosporná, potom aj $T_H$ bude bezosporná.

\paragraph{Dôkaz:} Nech $\mathcal{S}$ je množina všetkých uzavretých formúl jazyka $L$.
Ďalej definujeme $\mathcal{B} \{ S \mid S \subseteq \mathcal{S}, T \union S
\mbox{ je bezosporná}\}$. $\mathcal{B}$ je čiastočne usporiadaná na $\subseteq$.
Množina $\mathcal{B}$ má konečnú vlastnosť -- keď zoberiem ľubovoľnú podmnožinu
$S \subseteq \mathcal{S}$, bude platiť $$ S \in \mathcal{B} \iff S' \in
\mathcal{B}, S' \mbox{ je konečná}, S' \subseteq S$$

 Potrebujeme ukázať, že $\mathcal{S} \subseteq \mathcal{B}\times\mathcal{B}$ je čiastočné usporiadanie, čiže je

\begin{itemize}
\item reflexívna
\item antisymetrická
\item tranzitívna.
\end{itemize}

Reflexívnosť:
$$ \forall S \in \mathcal{B}: (S,S) \in \mathcal{S}$$

Antisymetrickosť:
$$ (S',S") \in \mathcal{S} \land (S",S') \in \mathcal{S} \implies S' = S"$$
$$ t.j. S' \subseteq S" \land S" \subseteq S' \implies S' = S"$$

Tranzitívnosť dokážeme analogicky.

$\mathcal{B}$ je čiastočne usporiadaná inklúziou.
$$\emptyset \in \mathcal{B}, \mbox{ lebo } T \union \emptyset = T$$
Ak by $\emptyset$ nebola v $\mathcal{B}$, potom by $T$ bola sporná.

Množina $\mathcal{B}$ má konečnú vlastnosť:
 $$S \in \mathcal{S}, S \in \mathcal{B} \iff \forall \textrm{konečná }S' \subseteq S, T \union S' \textrm{ je bezosporná, teda }S' \in \mathcal{B}.$$

$\Rightarrow:$

$$ S \in \mathcal{B} \implies T \union S  \mbox{ je bezosporná}$$
$$ S' \subseteq S, S' \mbox{ je konečná}, T \union S' \mbox{ je bezosporná}\implies S' \in \mathcal{B}$$

$\Leftarrow:$Obrátene: Predpokladajme, že $S'$ je ľubovoľná konečná podmnožina $S$, $T \union S'$ je bezosporná, $S' \in \mathcal{B} \implies S \in \mathcal{B}$.

Tvrdenie dokážeme sporom. Predpokladajme, že $S$ nepatrí do $B$. Potom $T \union S$ je sporná, čiže pre ľubovoľnú dokázateľnú formulu A, $\vdash \neg(A \implies A)$.

Nech $A_1, A_2, \cdots A_n$ je dôkaz (odvodenie) $\neg(A \implies A)$. Nech $B_1, B_2, \cdots B_m$ sú tie, ktoré sa v tom dôkaze vyskytujú a patria do $S$. Tvrdíme, že $m>=1$. Prečo? Inak by bola $T$ sporná.

$S'=\{B_1 ... B_m\}$ je konečná a nenachádza sa v množine $\mathcal{B} \implies$ SPOR (Jedine, kde mohol vzniknúť, je pri predpoklade $S \notin \mathcal{B} )\footnote{poznámky Lukáša H.}$

\subsection{Princíp maximality}

\paragraph{Definícia princípu maximality.}

Každá neprázdna neprázdna podmnožina $\mathcal{P}(\mathcal{S})$ s konečnou vlastnosťou má maximálny prvok vzhľadom na inklúziu.

Nech $S_0$ ja maximálny prvok množiny $\mathcal{B}$ vzhľadom na inklúziu. Položme rozšírenie $T'$ (ktoré hľadáme)  $T' = T \union S_0$. Nech $A$ je ľubovoľná uzavretá, potom v $T'$ buď platí
$T' \vdash A $ alebo $T' \vdash \neg A$.

Nech $T'$ nie je úplná teória. Čiže ani jedno z predošlých nie je dokázateľné. Čiže $\neg A$ sa nenachádza ani v $T$ ani $S_0$. 

Uvažujme, že $T' \union \{\neg A\}$ je bezosporná teória. To je spor s tým, že $S_0$ je maximálny prvok z $\mathcal{B}$ vzhľadom na inklúziu. Pozn.: Odvoláme sa na predtým dokázaný dôsledok: máme $A$, $A'$ je uzáver. Potom $T \vdash A \iff T \union \{\neg A'\}$ je sporná teória.

Dosiahli sme zúplnenie teórie. Máme $T$, získali sme $T'$, ktorá je bezosporná a je úplným rozšírením.

Keď máme $T$, $L$, vieme to rozšíriť na $T_H$ s jazykom $L(C)$. Keď máme $T$ bezospornú, $T(H)$ s jazykom $L(C)$ je tiež bezosporná. Dokážem vytvoriť Hankinovskú úplnú $T_H'$, $L(C)$ a pre ňu model $m$.

$T, L$ - bezosp.

$T', L'$ - úplná hakinovská

$T', L' - m'$

$T'$ je bezosp. rozšírenie $T$, potom $T \vdash A \iff T' \vdash A$ Potom platí, že keď zooberiem ľub. $A \in T$, potom  $m' \vdash A$. Pre ľubov. rozšírenie $T$, $T'$ $A\in T, T' \vdash A \implies m' \models A, m, v T$.

\paragraph{Definícia.}

$T,L$, $T',L'$, $L'$ je rozš. $L$. Nech $m'$ je realizácia jaz. $L'$, redukovaním štruktúry $m'$ na jaz. $L$ získame realiz. $T$ v jaz. $L$.

$m$: univerzum $m$ bude to isté univerzum ako univerzum $m'$.

$m$ obsahuje iba také relácie a zobrazenia, ktoré realizujú špecálne symboly jazyka $L$ v realizácii $m'$. Tzn. ak $f$ je ľub. n-ány funkčný symbol $L$, $f_m'$ je zobrazenie, ktoré realizuje $f$ v $m'$, potom zostáva realizáciou $f$ v $m$.

Podobne s predikátom $P$.

$m$ vzniklo redukciou $m'$ na jazyk $L$. $(m: m' \triangle L)$
\footnote{na prednáške bolo miesto $\triangle$ použíté niečo ako $\land$, ale jedna nožička bola dlhšia ako druhá}.
$m'$ je expanzia realizácie $m$

\paragraph{Poznámka.}
$m$ je redukcia $m'$ na jazyk $L$ ($m'\triangle L$), $A$,$L$. Chceme ukázať $m \models A \iff m' \models A$. Všimnime si, že univerzum, t.j. aj ohodnotenia a symboly, sú tie isté. Preto, ak $m' \models A$, potom aj $m \models A$ a obrátene. Teraz dokážeme, že každá bezosporná teória má model.

\paragraph{Veta.}
Nech $T'$ je rozšírenie teórie $T, L$. Ak $m$ je model $T'$ a $m:m'\triangle L$, tak $m$ je model $T$.

\paragraph{Dôkaz} To, čo sme tu už porozprávali. $m'$ je model $T$',
$A \in T$, $A$ je axioma $T$, $m' \models A \iff M \models A$.

Nasledujú dve vety o kompaktnosti a ich dôsledky. Dávajú nám konečnú charakterizáciu splniteľnosti.

\paragraph{Veta} \emph{O kompaktnosti.}
 $T$ je množina formúl jaz. $L$, $A$ je formula $L$, $T \models A \iff T' \models A, T' \subseteq T, T$ je konečná.

\paragraph{Dôkaz}


G\"odel: $T \models A \iff T \vdash A$

$A_1, A_2 ... A_n(A)$ je konečná postupnosť formúl - konečne veľa axiom z $T$.

$T' = \{B_1, B_2 ... B_m\}$, $B_i$ - axiomy v $T$ nachádzajúce sa v dôkaze $A$

Môžeme napísať $T' \vdash A \iff T' \models A$.

\paragraph{}
Na základe vety o kompaktnosti ukážeme, že teóriu telies charakteristiky 0 nedokážeme v teórii prvého rádu napísať charakterizovať pomocou konečného počtu axiom. Komut. grupa. distr. zákony ma krát a plus ... rád prvku - konečný, najmenší takýto prvok ... teleso konečnej charakteristiky. A tak ďalej.


\section{Prednáška 12.3.2009}

\par \textbf{Zapísal:} P. Struhár

\par

\paragraph{Veta o kompaktnosti} bola minule.

\paragraph{}
Nech $A$ je splnená vo všetkých telesách charakteristiky 0.
 $$ T' \models A \iff T'' \subseteq T', T''\mbox{ konečná} $$
 $$ T \models A $$
$T''$ obsahuje len konečne veľa axiom $\neg P_p$. Nech m je najväčší index
axiomy, ktorá patrí to $T''$. Potom formula $A$ je splnená aj vo všetkých telesách dostatočne
veľkej konečnej charakteristiky ($>m$). Ak teda formula je teorémou telies
charakteristiky $0$, tak platí aj vo všetkých telesách dostatočne veľkej charakteristiky.
$$ T' \models A_i \implies T' \models A_1 \land A_2 \land \cdots \land A_n $$
Logikou prvého rádu nemôžem konečným počtom axiom charakterizovať teóriu telies.
V logike prvého rádu kvantifikujem len individuá.

\paragraph{Veta}
Každá konečná množina formúl, ktorá je splnená vo všetkých telesách charakteristiky 0,
je splnená aj vo všetkých telesách dostatočne veľkej konečnej charakteristiky.

\paragraph{Poznámka}
Ak zavedieme ďalší druh premenných, prirodzené čísla, nekonečne veľa axiom
typu $\neg P_p$, nekonečne veľa axiom môžme nahradiť jednou formulou:
$$ (\forall p)(p\times 1\not=0)$$
Toto v literatúre terminologicky nazývajú \emph{slabá logika druhého rádu}. V nej
vieme konečne axiomatizovať teóriu telies charakteristiky 0.
\emph{Veta o kompaktnosti neplatí v slabej logike druhého rádu.}

Mali sme Godelovu vetu v dvoch tvaroch:
$$ T \models A \iff T \vdash A$$
Druhý tvar: Teória je bezosporná, keď má model. Z druhého tvaru môžeme sformulovať a dokázať
\emph{vetu o kompaktnosti}.

\paragraph{Veta o kompaktnosti}
Nech $T$ je množina formúl jazyka $L$. Model teórie $T$ existuje $\iff$ každá konečná podmnožina
$T' \subseteq T$, má model.

\paragraph{Dôkaz}
Vieme, že $T$ má model $\iff$ je bezosporná. Viď aj Lindenbaum -- konečná vlastnosť.
Každá $T'$ je bezosporná a teda má model.

Pomocou vety o kompaktnosti tohto tvaru dokážeme zaručiť, že pre Peanovu aritmetiku
existujú aj neštandardné modely. Najprv si však ukážeme Robinsonovu aritmetiku.

\begin{enumerate}
\item $\neg S(x) = 0$
\item $S(x) = S(y) \implies x=y$
\item $x+0 = x$
\item $x+S(y) = S(x+y)$
\item $x\cdot 0 = 0$
\item $x\cdot S(y)= (xy)+x$
Ak pridáme axiomu
\item $A_x[0] \implies ((\forall x)(A \implies A_x[S(x)]) \implies (\forall x)A)$,
\end{enumerate}
dostaneme Peanovu, pričom mám model $n = \langle N^+,0,D,+,-\rangle$

Ukážeme, že Peanova aritmetika má aj neštandardné modely. Zavedieme označenie:

$L$, definujeme $\bar{n}$ - numerály:
\begin{itemize}
\item $\bar{0} = 0$
\item $\overline{n+1} = S(\bar{n}) = S(S \cdots (S(0) \cdots ))$
\end{itemize}

K $L$ zostrojíme $L_c$ (do $L$ pridáme konstantu $c$).
$$\mbox{teória }P_c,\mbox{ jazyk }L_c$$
$$C_n: c \not= \bar{n}$$
$$\mbox{zoberieme }T \subseteq P_c, P_c \mbox{ obsahuje len konečne veľa numerálov.}$$

Vidíme, že každá konečná podmnožina má model, ktorý vznikne expanziou štandardného modelu $n$.
Konštantu c realizujeme indivíduom, ktoré nepatrí do $t$ (čoho?).
Čiže každá $T$ má model a na základe vety o kompaktnosti aj $P_c$ má model. Teraz, keď si všimneme,
že model $P_c$ bude rôzny od štandardného modelu a nebude ani izomorfný, pretože tam máme konštantu $c$.
Každá konečná podmnožina bude mať model. To je zabezpečené, že vznikne expanziou $n$
a bude obsahovať konštantu, ktorá sa predtým nevyskytovala.
$m\triangle L = n$ vznikne redukciou L. Dostaneme $n'$, ktorý bude obshovať konštantu c,
čím bude neštandardný a nebude izomorfný. Vieme, že existuje, ale nevieme ho skonštruovať.

\subsection{Metóda rezolvent}

Budeme pracovať s formulami, ktoré budú v prefixovom tvare a v konjuktívnej normálnej
forme. Keď uvažujeme výrokovú logiku, interpretujeme formulu $A$ funkciou $\bar{v}$
a zobrazením $v$. Formulu sme nazvali tautológiou, ak $\bar{v}(A)\cdots$, ak pre každú
interpretáciu dostaneme pravdu - $2^n$ interpretácii. V predikátovej logike ešte horšie úplným preberaním.

Zopakujme si, čo už vieme. V predikátovej logike -- funkčné a predikátové symboly. Univerzum budem označovať $D$.
M bude matica -- jadro konj. norm. formy (?).

Keď sme mali n-árny funkčný symbol $f: D^n \rightarrow D$. N-árny predikátový symbol
$P: D^n \rightarrow \{0,1\}$.

Rozlišovali sme voľné a viazaná premenné. Keď nemajú formuly voľné premenné, neviem
určiť pravdivostnú hodnotu, iba ak tam dosadím konštantu (?).

Ku každej formule A viem zostrojiť takú, ktorá je v prefixovom (prenexnom) tvare -
prefix + jadro: $(Q_1 x_+)\cdots(Q_n x_n)M$.

Prenexné opperácie:
$$(Qx)A(x) \vee B \equiv (Qx)(A(x) \vee B) \mbox{ - $x$ nie je voľné v $B$} $$
$$ (Qx)A(x) \land B \equiv (Qx)(A(x) \land B) \mbox{ - $x$ nie je voľné v $B$} $$
$$ \neg (\forall x) A(x) \equiv (\exists x) \neg A(x) $$
$$ \neg (\exists x) A(x) \equiv (\forall x) \neg A(x) $$
$$ (\forall x) A(x) \land (\forall X) B(x) \equiv (\forall x) (A(x) \land B(x)) $$
$$ (\forall x) A(x) \vee (\forall x) B(x) \equiv (\forall x) (A(x) \vee B(x)) $$
Platí toto?:
$$(\forall x) A(x) \vee (\forall X) B \equiv (\forall x) (A(x) \vee B(x))$$
Môže platiť pravá strana a ľavá nie, keď $A(a)=1, A(b)=0, B(a)=0, B(B)=1$.
Podobne neplatí
$$(\exists x) A(x) \land (\exists X) B(x) \equiv (\exists x) (A(x) \land B(x))$$

Vezmem si $(\forall x)B(x) a (\forall z) B(z)$:
$$(\forall x) B(x) \equiv (\forall z) B(z) $$
$$\cdots$$
$$ (Q_1 x)A(x) \vee (Q_2 x)B(x) \equiv (Q_1 x)(Q_2 z)(A(x) \vee B(z)) $$
$$ (Q_3 x)A(x) \land (Q_4 x)B(x) \equiv (Q_3 x)(Q_4 z)(A(x) \land B(z)) $$
proste opakovanie z minulého semestra.

\subsection{Algoritmus na zostrojenie prenexného tvaru}
\paragraph{Krok 1} Odstránenie implikácií a ekvivalencií:
$$ 
	A \iff B \equiv (A \implies B) \land (B \implies B) \\
	A \implies B \equiv \neg A \lor B
$$

\paragraph{Krok 2} Odstránenie dvojitej negácie a presun negácie k formule.
$$
\begin{array}{lll}
	\neg (\neg A) &\equiv& A \\
	\neg (A \lor B) &\equiv& (\neg A \land \neg B) \\
	\neg (A \land B) &\equiv& (\neg A \lor \neg B) \\
	\neg (\forall x) A(x) &\equiv& (\exists x) \neg A(x) \\
	\neg (\exists x) A(x) &\equiv& (\forall x) \neg A(x)
\end{array}{lll}
$$

\paragraph{Krok 3} Premenovanie viazaných premenných, ak je to potrebné.
\paragraph{Krok 4} Použijeme zákony:
$$
\begin{array}{lll}
	(Q x) A(x) \lor B &\equiv& (Q x)(A(x) \lor B) \\
	(Q x)A(x) \land B &\equiv& (Q x)(A(x) \land B) \\
	(\forall x) A(x) \land (\forall x)B(x) &\equiv& (\forall x)(A(x) \land
	B(x))\\
	(\exists x)A(x) \lor (\exists x)B(x) &\equiv& (\exists x) (A(x)\lor B(x))
	\\
	(Q_1 x) A(x) \lor (Q_2 x)B(x) &\equiv& (Q_1 x)(Q_2 z) (A(x)\lor B(z))
	\mbox{z nie je v A(x)}. \\
	(Q_3 x) A(x) \land (Q_4 x)(B(X) &\equiv& (Q_3 x)(Q_4 z)(A(x) \land (B(z))
	\\
\end{array}
$$

\paragraph{Príklad} $$
\begin{array}{l}
(\forall x)(\forall y) \left( (\exists z)(P(x,z) \land
P(y,z)\right) \implies (\exists u) Q(x,y,u)) \equiv \\
(\forall x)(\forall
y)\left(\neg (\exists z)(P(x,z)\land P(y,z))\lor(\exists u)Q(x,y,u)\right)
\equiv \\
(\forall x)(\forall y)\left((\forall z)(\neg P(x,z)\lor \neg
P(y,z)\right)\lor(\exists u)Q(x,y,u)) \equiv \\
(\forall x)(\forall y)(\forall z)(\exists u)(\neg P(x,z) \lor \neg
P(y,z) \lor Q(x,y,u)) \\
\end{array}
$$

\section{Herbrandova veta}
\par Leibniz (1646 -- 1716), Peáno (1900), Hilbertova škola (1920), Herbrand
(1930). Fundamentálny výsledok, ktorý dal odpoveď na otázku, či existuje
všeobecná procedúra, ktorá vie zistiť, či je daná formula tautológia, dali
Church a Turing (1936, nezávisle od seba) -- ukázali, že v predikátovej logike
nie je problém rozhodnuteľnosti riešiteľný.
\par Herbrandova metóda je založená na procedúre vyvrátenia. Ide o nájdenie
interpretácie, pri ktorej daná formula nie je splnená. Ak formula nie je
tautológia, v konečnom počte krokov sa procedúra zastaví, pokiaľ ale nie je,
procedúra pokračuje vo výpočte.

\par Gilmore (1960), Davis a Putnem (1960) -- navrhovali procedúry, ktoré by
boli schopné overovať tautologickosť formúl; fungovali ale iba pre jednoduchšie
formuly.

\par V rokoch 1965 -- 1968 Robinson zaviedol pojem \emph{resolventa} (v podstate
iným spôsobom zapísane pravidlo \emph{modus ponens}.

\section{Skolemovské štandardné formuly}
Procedúry založené na Herbrandovej myšlienke vyvrátenia (nájdenie interpretácie,
ktorá vyvracia danú formulu). Davis a Putnam pracovali s formulou v prefixovom
normálnom tvare a vychádzali z nasledujúcich hľadísk:
\begin{enumerate}
	\item Každá formula môže byť vyjadrená v prefixovom normálnom tvare.
	\item Pretože matica (jadro) neobsahuje kvantifikátory, môže byť
	vyjadrená v konjunktívnej normálnej forme.
	\item Zachovajúc nesplniteľnosť formúly môžeme z nej eliminovať
	existenčné kvantifikátory pomocou skolemovských funkcií (keď máme
	premennú vyjadrenú v prefixovom normálnom tvare a eliminujeme z nej
	existenčný kvantifikátor, dostaneme formulu, ktorá je nesplniteľná
	vtedy, keď je nesplniteľná pôvodná formula -- nesplniteľnosť
	transformovanej formuly sa robí ľahšie).
\end{enumerate}

Transformácia formuly. Majme formulu $A: (Q_1 x_1) (Q_2 x_2) \ldots (Q_n x_n)
M$, $M$ -- konjunktívny normálny tvar. $Q_r$ je existenčný kvantifikátor v
prefixe $(Q_1 x_1) \ldots (Q_n x_n)$, $1 \leq r \leq n$ vľavo od $Q_r$ nestojí
žiaden veľký kvantifikátor.
$A': Q_r$, $M$, za $x_r$ dosadíme $c$ ktorá v $M$.
$Q_{s_1}, Q_{s_2}, \ldots Q_{S_m}$ (zoznam všeob. kvantifikátorov) vľavo od
$Q_r$,  $1 \leq s_1 \leq s_2, \ldots \leq s_m < r$.

Za premennú $x_r$ dosadíme do $M$ term $f(x_{s_1}, \ldots, x_{s_m})$. Tieto
funkcie voláme Skolemove funkcie a $c$ voláme Skolemove konštanty.

\paragraph{Príklad.} Nájdite štandardnú formu pre formuly:
$$ (\exists x) (\forall y) (\forall z) (\exists u) (\forall v) (\exists w)
P(x,y,z,u,v,w)$$.
Premennú $x$ budeme nahrádzať konštantou $c$ (ktorá nie je v predikáte),
premennú $u$ nahradíme $f(y,z)$ a premennú $w$ nahradíme funkciou $g(y,z,v)$ .

$$ (\forall y)(\forall z)(\forall w) P(x,y,z,f(y,z),v,g(y,z,v)) $$

\paragraph{Príklad:} Nájdite štandardnú formulu: $$ (\forall x)(\exists y)
(\exists z) \left((\neg P(x,y) \land Q(x,z))\lor R(x,y,z)\right) $$.

$$ (\forall x) (\exists y) (\exists z)(\neg P(x,y) \lor R(x,y,z)) \land
(Q(x,z)\lor R(x,y,z))$$

$$ (\forall x) (\neg P(x,f(x)) \lor R(x,f(x),g(x)) \land Q(x,g(x))\lor
R(x,f(x),g(x)))$$

\paragraph{Definícia:} \emph{Literál} -- premenná, alebo jej negácia. 
\paragraph{Definícia:} \emph{Klauzula} -- disjunkcia literálov. \emph{Jednotková
klauzula} (tiež jednoliterálna klauzula) obsahuje jeden literál. 

\paragraph{Definícia:} \emph{Prázdna klauzula} neobsahuje žiaden literál
(označujeme ju $\square$). Prázdna klauzula je nesplniteľná.  

\par (niečo na tabuli)
$\{\neg P(x,y), R(x,y,z) \}$ \\
$P \lor Q$, $\{P, Q\}$

$y: f(x)$
$z: g(x)$

\paragraph{Veta:} Nech $\mathcal{S}$ je množina klauzúl, ktoré tvoria štandardnú
formu formuly $A$. Potom $A$ nie je splniteľná práve vtedy, keď $\mathcal{S}$
nie je splniteľná.

\paragraph{Dôkaz:} $A: (Q_1 x_1) \ldots (Q_m x_n) M[x_1, \ldots x_n]$. $Q_r$ je
existenčný kvantifikátor, potom z $A$ dostávame formulu $A_1$ v tvare $$A_1:
(\forall x_1) \ldots (\forall x_{r-1}) (Q_{r+1} x_{r+1})  \ldots (Q_n x_n)
M[x_1, \ldots x_{r-1}, f(x_1, \ldots, x_{r-1}), x_{r+1}, x_n]$$. $A$ -- nie je
splniteľná práve vtedy, keď $A_1$ nie je splniteľná.

\par Ak $A$ nie je splniteľná, potom ani $A_1$ nie je splniteľná. Nech $A_1$ je
splniteľná. Potom existuje taká interpretácia $I$, že $I$ vyhovuje $A_1$ ($A_1$
je pravdivé v $I$). To znamená, že pre všetky $x_1 , \ldots x_{r-1}$ existuje
také $x_r$, ktoré je rovné $f(x_1, \ldots, x_{r-1})$. Tá istá interpretácia teda
vyhovuje aj pôvodnej formule $A$.

\par Obrátené tvrdenie. Predpokladajme, že $A_1$ nie je splniteľná a teda
formula $A$ nie je splniteľná. Nech, že $A$ je splniteľná. Potom pre každé $x_1,
\ldots x_{r-1}$ existuje také $x_r$, že platí $M[x_1, \ldots x_{r-1}, x_r,
\ldots]$. Rozšírme $I$ na $I'$ tak, že $f(x_1, \ldots x_{r-1}) = X_R$. Ak $A_1$
nie je splniteľná, tak ani $A_2$ nie je splniteľná. Pôvodnú formulu označíme
$A_0$, a $A: A_0, \ldots A_m$, pričom platí $A_0 \iff A_1 \iff \cdots \iff A_m$.
$A_m: A'$
$A_0 \iff A'$


\par Toto platí iba vtedy, ak formula nie je splniteľná. Majme formulu $A$, $A:
(\exists x) P(x)$, a $S: P(a)$. Vezmime interpretáciu: $D = \{1, 2 \}$, pričom
$a: 1$. Hodnoty pre $P$. $P(1)$ -- nepravda, $P(2)$ -- pravda.

(opakovanie z minula) 

\paragraph{Poznámka} Majme formulu $A: A_1 \land A_2 \land \ldots \land A_n$.
Pre každé $A_1$ si vytvoríme štandardnú formulu $S_1, S_2, \ldots S_m$ Potom
štandardná forma pre $A$ je $S_1 \union S_2 \union \ldots \union S_m$ (kde $S_i$
je množina klauzúl). Formula $A$ nie je splniteľná práve vtedy, keď nie je
splniteľná táto množina klauzúl.

\paragraph{Príklad} Uvedieme príklad, ktorý sa pokúsime zapísať pomocou množiny
klauzúl -- príklad z teórie grúp. 

\par Uvažujme grupu a na nej operáciu $\circ$. Platia nasledovné axiomy:
\begin{itemize}
\item $A1: x, y \in G, x \circ y \in G$ -- vlastnosť uzavretosti.
\item $A2: x, y, z \in G, (x \circ y) \circ z = x \circ (y \circ z)$ -- asociatívnosť
operácie.
\item $A3: x \in G, x \circ e = e \circ x = x$ -- existencia neutrálneho prvku.
\item $A4: x \in G, \exists x^{-1} \in G, x \circ x^{-1} = e = x^{-1} \circ x$ --
existencia inverzného prvku.
\end{itemize}
Chceme ukázať, že ak pre všetky $x \in G$ platí $x \circ x \in G$, potom grupa
je komutatívna. Prvou časťou úlohy je formalizovanie zadania:

$$ P(x,y,z), x \circ y = z, i(x) = x^{-1}$$
Potom axioma $A1$ vyzerá nasledovne:
$$(A'1) (\forall x) (\forall y) (\exists z) P(x,y,z)$$


$$ 
\begin{array}{l}
A'2:(\forall x) (\forall y) (\forall z) (\forall u) (\forall v) (\forall w) (P(x,y,u) \land  P(y,z,v)
\land P(u, z, w) \implies  
P(x,v,w)) \\ \land (\forall x) (\forall y)(\forall z)
(\forall u)(\forall v) (\forall w)
P(x,y,u) \land P(y,z,v) \land P(x,v,w) \implies P(u,z,w)
\end{array}
$$


$$A'3: (\forall x) P(x,e,x) \land (\forall x) P(e,x,x)$$
$$A'4: (\forall x) P(x,i(x),e) \land (\forall x) P(i(x),x,e)$$

Záver vety:
$$B: x \circ x, x \in G, u \
\circ v = v \circ u; u,v \in G$$
$$B': (\forall x) P(x,x,e) \implies  (\forall u) (\forall v) (P(u,v,w) \implies
P(v,u,w))$$.

Teraz potrebujeme ukázať, že $A_1' \land A_2' \land A_3' \land  A_4' \implies
B'$. Vytvoríme negáciu výroku a ukážeme, že nie je splniteľná:
$$F: \neg (A_1' \land A_2' \land A_3' \land A_4') \lor B'$$
$$\neg F:A_1' \land A_2' \land A_3' \land A_4' \land \neg B'$$.

Ideme vytvárať množinu klauzúl. Máme axiomy $A_1'$, $A_2'$, $A_3'$, potom
množiny sú:
$$ S_1': \{ P(x,y,f(x,y)) \} $$
$$ S_2': \{ \neg P(x,y,u) \lor \neg P(y,z,v) \lor \neg P(u,z,w) \lor P(x,v,w),
\neg P(x,y,u), \lor \neg P(y,z,v) \lor \neg P(x,v,q) \lor P(u,z,w)\}$$
$$S_3': \{ P(x,e,x), P(e,x,x) \}$$
$$S_4': \{ P(x,i(x),e),P(i(x),x,e) \}$$


TOTO NIEKTO SKONTROLUJTE:
$$ \neg B': \neg ((\forall x) P(x,x,e) \implies (\forall u) (\forall v)(\forall
w)(P(u,v,w) \implies P(v,u,w)))$$.
$$ \neg ( \neg (\forall x) (P(x,x,e) \lor (\forall u) (\forall v) (\forall w)
(\neg P(u,v,w)) \lor P(v,u,w))) \iff (\forall x) (P(x,x,e) \land \neg (\forall
u)(\forall v) (\forall w) (\neg P(u,v,qň \lor (P(v,u,w) \iff (forall z) P(x,x,e)
\land (\exists u) (\exists v) (\exists w) (P(u,v,w)\land \neg (P(v,u,w)$$

$$T: \{P(x,x,e), P(a,b,c), P(b,a,c) \} $$

Máme dokázať, že množina klauzúl nie je splniteľná, a teda pre každú oblasť a
interpretáciu je nesplniteľná -- bolo by treba vyšetriť všetky oblasti. Nám ale
bude stačiť vyšetrovať na \emph{Herbrandovskom univerze}.

\section{Herbrandovské univerzum}
\paragraph{Definícia} \emph{Herbrandovské univerzum množiny klauzúl} $H_0$ je
množina konštánt, ktoré sa vyskytujú v $S$. Ak $S$ neobsahuje žiadnu
konštantu, tak kladieme $H_0=\{ a \}$, kde $a$ je nejaká konštanta. Ďalej
definujeme pre $i=1,2,\ldots$ $H_{i+1}$ ako zjednotenie $H_{i}$ a množiny
všetkých termovtvaru $f^{(n)}(f_1,\ldots, f_n)$, kde $f^{(n)} \in S$ a $f_1,
\ldots, f_n \in H_1$. $H_i$ nazývame \emph{herbrandovské univerzum $i$-tej
úrovne}. $H=\union_{i=0}{\infty} H_i$.

\paragraph{Príklad} $S= \{ P(x), \neg P(x) \lor \neg P(f(x))\}$. 

$$
\begin{array}{lll}
H_0& = &\{ a \} \\
H_1& = &\{ a, f(a) \} \\
H_2& = &\{ a, f(a), f(f(a)) \} \\
\vdots & \vdots & \vdots \\
H &=& H_\infty = \{ a, f(a), f(f(a), f(f(f(a))), \ldots \}.
\end{array}
$$

\paragraph{Príklad} $S=\{P(x) \lor R(x), R(z), T(y) \lor \neg W(y) \}$ --
množina neobsahuje žiadnu konštantu, a preto kladieme $H_0 = \{ a \}$. $\{ a \}
=  H = H_0 = H_1$.

\paragraph{Príklad} Uvažujme množinu klauzúl $S=\{P(f(x),a g(y), b)\}$. Potom
$$
\begin{array}{lll}
H_0 &=& \{ a, b\} \\
H_1 &=& \{ a, b, f(a), g(a), f(b), g(b) \} \\
H_2 &=&  \{ a, b, f(a), g(a), f(b), g(b), f(f(a)), f(g(a)), f(f(b)), f(g(b)),
g(f(a)), g(f(b), g(g(a)), g(g(b)) \} \\
\end{array}
$$

\paragraph{Definícia} Nech $S$ je množina formúl. Potom množina atómov
tvaru $p^{(n)}(t_1, \ldots, t_n)$ pre všetky $n$-árne predikáty, ktorá sa
vyskytujú v $S$ a všetky termy $t_i$ z Herbrandovského univerza
nazývame \emph{Herbrandovskou  bázou} $S$. (všimnite si, že atómy sú
takého tvaru, že sa v nich nevyskytuje žiadna premenná).

Ľubovoľný výraz, ktorý neobsahuje premenné, sa nazýva \emph{základný výraz}
(základný term, základný atóm, základná klauzula, základný literál, ...).

\paragraph{Definícia} \emph{Základná inštancia} klauzuly C z množiny klauzúl
$S$ je klauzula, ktorú dostaneme zámenou premenných prvkami z
Herbrandovského univerza.

\paragraph{Príklad} Majme $S = \{P(x), Q(f(x)) \lor R(y) \}$, a nech
$C: P(x)$. $H = \{ a, f(a), f(f(a)). f(f(f(a))), \ldots \}$. Základné inštancie
pre $C: P(a), P(f(a)), P(f(f(a))), \ldots$.

\paragraph{Definícia} $S$ je množina klauzúl, $H$ je herbrandovské
univerzum pre množinu klauzúl $S$, $I$ je interpretácia v množine
klauzúl $S$ nad $H$. Hovoríme, že interpretácia $I$ je \emph{herbrandovská
interpretácia}, ak platí: 

\begin{enumerate}
	\item interpretácia $I$ zobrazuje všetky konštanty na samé seba (t.j.
	konštante $a \in S$ priradí tú istú konštantu $a \in H$.
	\item $f^{(n)}$ je $n$-árny funkčný symbol a $h_1, \ldots h_n$ sú prvky
	herbrandovského univerza $H$. V $I$ označujeme ako $f^{(n)}$ funkciu,
	ktorá zobrazuje $(h_1, \ldots h_n)$ z $H_*$ do $H$. $$f^{(n)}(h_1,
	\ldots, h_n)$$
\end{enumerate}

\paragraph{Poznámka} Poznamenávame, že nekladieme žiadne ohraničenia na initerpretáciu predikátov.

Uvažujme množinu $A = \{ A_1, A_2, \ldots A_n, \ldots \}$. Nech je to
herbrandovská báza pre množinu klauzúl $S$. Herbrandovskú interpretáciu určíme
tak, že  $I$ zadáme ako $I=\{ m_1, m_2, \ldots, m_n, \ldots \}$. 

\par $m_j$ bude buď $A_j$ (ak $A_j$ je pravdivé) alebo $\neg A_j$ (ak $A_j$ je
nepravdivé).

\paragraph{Príklad} $S=\{ P(x) \lor Q(x), R(f(y)) \}$. Herbrandovské univerzum
vyzerá: $H=\{a, f(a), f(f(a)), \ldots \}$. V $S$ sa vyskytujú predikáty $P$,
$Q$, $R$ --  všetky sú unárne. Herbrandovská báza je potom $A: \{ P(x), Q(x),
 P(f(x)), Q(f(x)), R(f(x)) \}$
 
$I_1: \{ P(x), Q(x), P(f(x)), Q(f(x)), R(f(x)) \}$

Čo znamená toto? :
$  \{ \neg P(a), \neg Q(a), \neg P(f(a)), \neg Q(f(a)), \neg R(f(a)), \ldots \}$


Ďalšou úlohou bude k ľubovoľnej interpretácii nad ľubovoľným univerzom
priradiť Herbrandovskú interpretáciu a zaviesť príslušné tvrdenia a definície.

\paragraph{Príklad} $S=\{ P(x) \lor Q(x). R(f(y)) \}$.  Jedna z možných
Herbrandovských interpretácií je nasledovná:
$$ I^* = \{ \neg P(x), \neg Q(a), \neg R(a), \neg P(f(a)), \neg Q(f(a)), \neg
R(f(a)), \ldots \}$$
$$ I_2^* = \{ P(a), Q(a), \neg R(a), P(f(a)), Q(f(a)), \neg R(f(a)), \ldots \}$$

\paragraph{Príklad} Majme množinu $S = \{ P(x), Q(y, f(y,a)) \}$. Oblasť
interpretácie $D = \{1,2\}$.

\begin{tabular}{|l|c|c|c|c|}
\hline
a & f(1,1) & f(1,2) & f(2,1) & f(2,2) \\
\hline
2 & 1 & 2 & 2 & 1 \\
\hline
\end{tabular}

\begin{tabular}{|l|c|c|c|c|c|}
\hline
P(1) & P(2) & Q(1,1) & Q(1,2) & Q(2,1) & Q(2,2) \\
\hline
T & F & F & T & F & T \\
\hline
\end{tabular}

$I^*$, H-im ??
Herbrandovská báza: $$A=\{ P(a), Q(a,a), P(f(a,a)), Q(a,f(a,a)), Q(f(a,a), a,
\ldots \}$$.

Hodnoty pre príslušné predikáty určíme pomocou zadaných tabuliek:
$$
\begin{array}{lll}
P(a) &=& P(2) \mbox{neplatí} \\
Q(a,a) &=& Q(2,2) \mbox{platí} \\
P(f(a,a)) &=& P(f(2,2)) = P(1) \mbox{platí} \\
Q(f(a,a),a) &=& Q(f(2,2),2) = Q(1,2) \mbox{platí} \\
Q(a,f(a,a)) &=& Q(1,f(2,2)) = Q(2,1) \mbox{platí} \\
Q(f(a,a),f(a,a)) &=& Q(f(2,2), f(2,2)) = Q(1,1) \mbox{neplatí} \\
\end{array}
$$

$$I^* = \{ \neg P(a), Q(a,a), P(f(a,a)), \neg Q(a,f(a,a)), Q(f(a,a),a), \ldots
\}$$

\par
Ďalšia vec, čo sa môže stať -- majme množinu klauzúl, ktorá neobsahuje
konštantu: $S=\{P(x), Q(y,f(y)) \}$. Máme danú interpretáciu $I$ s oblasťou
$D=\{1,2\}$, oblasti priradím dve interpretácie $I_1^*$ a $I_2^*$, pričom v
prvej bude $a$ interpretované ako $1$ a v druhej ako $2$.

\paragraph{Definícia} Majme interpretáciu $I$, množinu klauzúl $S$ na oblasti
interpretácií $D$. $I^*$ je H-interpretácia priradená ku $I$, vytvorená
nasledovne: Nech $h_1, h_2, \ldots h_n$ patria do Herbrandovského univerza $H$,
$I(h_1) \mapsto d_i \in , i=1 \ldots n$. $P^{(n)}$ patrí do množiny klauzúl $S$.
Ak $P^{(n)}(d_1, \ldots d_n)$ je splnený a pravdivý v $I$, potom
$P^{(n)}(h_1,\ldots,h_n)$ je splnený a pravdivý v $I^*$.

\paragraph{Lema} Majme interpretáciu $I$ na oblasti $D$. Nech táto interpretácia
vyhovuje množine klauzúl $S$. Potom ľubovoľná H-interpretácia $I^*$, ktorá je
priradená (zodpovedá) $I$ taktiež vyhovuje množine klauzúl $S$.

\paragraph{Dôkaz} Majme množinu klauzúl $S=\{C_1, C_2, \ldots , C_n\}$. Nech
klauzula (disjunkcia literálov) $C_i = L_{i_1} \lor L_{i_2} \lor \ldots \lor
L_{i_{r_i}}$, $i=1,\ldots,n$. Ľubovoľný literál $L_{i_j}$ je tvaru $P^{(n)}(d_1,
\ldots, d_n)$. Množina klauzúl $C_1 \land C_2 \land \ldots \land C_n$ je
pravdivá práve vtedy, keď je pravdivá každá klauzula.
$$I(h_i) = d_i$$
$$L^*_{i_j}: P^{(n)}(h_1, \ldots, h_n)$$

\paragraph{Veta} Množina klauzúl $S$ nie je splniteľná práve vtedy, keď $S$ je
nepravdivá pri všetkých H-interpretáciách $S$. 

\paragraph{Dôkaz} Jedna implikácia (zľava doprava) je očividná:

\subparagraph{Zľava doprava} Nech $S$ nie je splniteľná, potom je nepravdivá pre ľubovoľnú interpretáciu
na ľubovoľnej oblasti, a teda aj pre ľubovoľnú H-interpretáciu na H-univerze.

\subparagraph{Obrátená implikácia} Nech množina klauzúl $S$ je nepravdivá pre
ľubovoľnú H-interpretáciu množiny klauzúl $S$. Pre spor predpokladajme, že
existuje interpretácia $I$ s oblasťou $D$ pre množinu klauzúl $S$, ktorá
vyhovuje $S$. Uvažujme $I^*$, ktorá je priradená (zodpovedá) interpretácii pre
množinu klauzúl $S$. Podľa lemy, ktorú sme dokázali, ak vyhovuje $I$, vyhovuje
aj $I^*$, čo je spor. $\Box$

\paragraph{Poznámka}
\begin{enumerate}
	\item Základná inštancia $C'$ (neobsahuje premenné), $C$ je splniteľná v interpretácii $I$
	práve vtedy, keď existuje základný literál $L'$:
	$$ L' \in C'$$
	$$ C' \cap I \neq \emptyset, L' \in I = \{m_1, m_2, \ldots \} $$

	\item Klauzula $C$ je splnená v interpretácii $I$ $\iff$ každá jej
	základná inštancia $C$ je splnená v $I$.

	\item Klauzula $C$ je odmietnutá (vyvrátená) v $I$ $\iff$ keď existuje
	aspoň jedna taká základná inštancia $C'$, ktorá je odmietnutá v $I$
	(vyvrátená).

	\item Množina klauzúl $S$ nie je splniteľná v $I$, ak existuje aspoň
	jedna základná inštancia $C'$ klauzuly $C$, ktorá nie je splniteľná v
	$I$. $$C \in S$$
\end{enumerate}


\paragraph{Príklad} Uvažujme klauzulu $C=\neg P(x) \lor Q(f(x))$ a
interpretáciu:
$$
\begin{array}{lll}
I_1 &=& \{ \neg P(a), \neg Q(a), \neg P(f(a)), \neg Q(f(a)), \neg P(f(f(a))),
\neg Q(f(f(a))), \ldots \} \\
I_2 &=& \{ P(a), Q(a), P(f(a)), Q(f(a)), P(f(f(a))), Q(f(f(a))), \ldots \} \\
I_3 &=& \{ P(a), \neg Q(a), P(f(a)), \neg Q(f(a)), P(f(f(a))), \neg Q(f(f(a))),
\ldots \} \\
\end{array}
$$

\paragraph{Príklad} Uvažujme množinu klauzúl $S=\{P(x), \neg P(x)\}$ a
interpretácie:
$$
\begin{array}{lll}
I_1 &=& \{ P(x) \} \\
I_2 &=& \{ \neg P(x) \} \\
\end{array}
$$

Množina nie je splnená ani jednou interpretáciou.

\section{Sémantické stromy}
Splniteľnosť množiny a konštrukcia sémantického stromu budú pre nás ekvivalentné
pojmy :).

\paragraph{Definícia} $A$ je nejaký atóm (elementárna formula), potom $A$ a
$\neg A$ nazývame \emph{kontrárna dvojica}.

\par  Ak klauzula obsahuje kontrárnu dvojicu, potom je tautológia.

\paragraph{Definícia} Nech $S$ je množina klauzúl, $A$ je Herbrandovské báza pre
množinu $S$. Pod \emph{sémantickým stromom} pre množinu klauzúl $S$ je dole
rastúci strom, pričom každej hrane jej pripísaná množina atómov alebo negácií
atómov (teda vlastne množina literálov) z Herbrandovskej bázy, pričom platí:
\begin{enumerate}
	\item Z každého vrchola stromu vychádza konečný počet hrán (označujeme
	ich $l_1, l_2, \ldots l_n$. Nech $Q_i$ je konjunkcia všetkých literálov
	pripísaných hrane $l_i$, $i=1,2,\ldots,n$, potom $Q_1 \lor Q_2 \lor
	\ldots \lor Q_n$ je tautológia.
	\item Pre každý vrchol $v$ je $I(v)$ zjednotenie všetkých literálov
	pripísaných hranám vetvy, ktorá vedie do vrcholu $v$. Potom $I(v)$
	neobsahuje kontrárne dvojice.
\end{enumerate}


\paragraph{Definícia} Majme $A=\{ A_1, A_2, \ldots A_n \}$ je Herbrandovská báza
pre množinu $S$. Hovoríme, že sémantický strom prislúchajúci množine klauzúl $S$
je úplný, ak pre každý vrchol $v$ platí, že pre každé $i$ končiace vo vrchole,
$I(v)$ obsahuje $A_i$ alebo $\neg A_i$ (ale nie obe naraz). (TODO: check) 

\paragraph{Príklad} $A: \{ P, Q, R \}$, viď obrázok
\ref{fig-semantic_tree_example}.
\begin{figure}[h]
	\caption{Príklad sémantického stromu}
	\centering\includegraphics{img/xx/2009-03-24.1.mps}
	\label{fig-semantic_tree_example}
\end{figure}



\begin{figure}[h]
	\centering\includegraphics{img/xx/2009-03-26.1.mps}
	\caption{Príklad stromu}
	\label{0326_tree1}
\end{figure}
\paragraph{Príklad}
$S=\{P(x), ... \}$, H-báza je $\{P($ ... TODO

\paragraph{Príklad} $S=\{P(x), Q(f(x)) \}$, herbrandovská báza pre množinu
klauzúl je $\{ P(a), Q(a), P(f(a)), Q(f(a)), \ldots \}$ -- je nekonečná, teda aj
prislúchajúci sémantický storm je nekonečný (obr. \ref{0326_tree2}). %pic2
\begin{figure}[h]
	\centering\includegraphics{img/xx/2009-03-26.2.mps}
	\caption{Ďalší sémantický strom}
	\label{0326_tree2}
\end{figure}

Nech $v$ je vrchol a $I(v)$ je \todo{}. Jednotlivé vetvy sémantického stromu nám
preberajú všetky interpretácie sémantickej bázy. Nech $I'(v)$ je množina
základných inštancií odmietnutá (úsek zamieta množinu klauzúl), potom na tomto
mieste môžeme strom \emph{odrezať}.

\paragraph{Definícia} Vrchol $v$ sémantického stromu pre množinu klauzúl $S$ sa
nazýva \emph{odmietajúcim}, ak  $I(v)$ odmieta niektorú základú inštanciu
klauzuly z množiny $S$, no ľubovoľný predchádzajúci vrchol $v'$ ($v'$ predchádza
$v$), $I(v')$ neodmieta žiadnu základnú inštanciu klauzúl z $S$.

\paragraph{Definícia}  Hovoríme, že sémantický strom $T$ pre množinu klauzúl $S$
je \emph{uzavretý}, ak každá vetva $T$ končí odmietajúcim vrcholom.

\paragraph{Definícia} Vrchol $v$ sémantického (TODO alebo uzavretého???) stromu
$T$ pre  množinu klauzúl $S$ nazývame \emph{akceptujúcim}, ak všetky nasledujúce
vrcholy vrchola $v$ sú odmietajúce.

\paragraph{Príklad} $S= \{ P, Q \lor R, \neg P \lor \neg Q, \neg R \lor \neg P
\}$. Herbrandovskú báza množinu $S$ vyzerá nasledovne: $\{P, Q, R\}$ (obr.
\ref{0326_tree3}, \ref{0326_tree_cut}).

\begin{figure}[h]
	\centering\includegraphics{img/xx/2009-03-26.3.mps}
	\caption{A ďalší sémantický strom\ldots}
	\label{0326_tree3}
\end{figure}
%pic3+4

\begin{figure}[h]
	\centering\includegraphics{img/xx/2009-03-26.4.mps}
	\caption{Orezaný sémantický strom}
	\label{0326_tree_cut}
\end{figure}


\paragraph{Príklad} $S=\{ P(x), P(x) \lor Q(f(x)), \neg Q(f(x)) \}$.
Herbrandovská báza pre množinu klauzúl $S$ vyzerá: $\{ P(a), Q(a), P(f(a)),
Q(f(a)), \ldots \}$ (obr. \ref{0326_tree_cut2}). %pic5


\begin{figure}[h]
	\centering\includegraphics{img/xx/2009-03-26.5.mps}
	\caption{Ďalší orezaný sémantický strom}
	\label{0326_tree_cut2}
\end{figure}

\paragraph{Poznámka} Dirichletov princíp:
$$T_1:$$
$$X, Y -- \mbox{konečné množiny}$$
$$ |X| < |Y| $$
$$ f: X \rightarrow Y$$
$$ \exists y \in Y_1\qquad x_1,x_2 \in X, x_1 \neq x_2 \qquad f(x_1) = f(x_2) =
y$$
Sporom. Ak by také $x_1$, $x_2$ neexistovali, potom pre každé $x_1, x_2 \in X$,
$x_1 \neq x_2$ a $f(x_1) \neq f(x_2)$, potom $f$ je injektívne zobrazenie, z
čoho vyplýva, že $|X| \leq |Y|$.
$$T_2: $$
$$ f: X \rightarrow Y,\qquad X\mbox{-- nekonečná}, Y\mbox{ -- konečná} $$
$$T_2: \exists y\in Y, \qquad y=\{x|x\in X, f(x)=y\}$$

\paragraph{Dôkaz} Nech $|Y|=n$. Množinu $X$ môžeme napísať ako $X=\Cup_{y\in
Y}A_y$ takých, že ak $y_1\neq y_2$, potom $A_{y_1} \cap A_{y_2} = \emptyset$.

$$ |X| = \left|\Cup_{y\in Y} A_y \right|$$

Nech $k_0$ je maximálna mohutnosť $|A_y|$, $y\in Y$, množina $X$ ale nemá
konečnú mohutnosť -- spor.
\section{Herbrandova veta}

\paragraph{Definícia} Strom chápeme ako usporiadanú dvojicu $(T,\leq)$ -- $T$ a
relácia, ktorá čiastočne usporadúva $T$:
\begin{enumerate}
	\item $T(n)=\{v\in T, v<n\}$ je dobre usporiadaná.
	\item $T$ má najmenší vrchol (koreň)
\end{enumerate}

Množina je dobre usporiadaná, ak každá jej neprázdna podmnožina má najmenší
prvok. $T$ je lineárne usporiadanie -- ľubovoľné dva prvky sú porovnateľné:
$a \neq b$, $a,b\in T$, pre $\{a, b\}$ musí platiť $a<b$ alebo $a>b$.

\par Ak mám prvky $u$ a $v$, prvok $v$ nasleduje bezprostredne po prvku $u$, ak
neexistuje $z$ také, že $u < z < v$.

\paragraph{Lema (K\"onig)} (je dôležitá na pochopenie dôkazu Herbrandovej vety)
Nech každý vrchol stromu s koreňom má konečné vetvenie (t.j. konečný stupeň
vetvenia) a strom $T$ je nekonečný, tak potom v ňom existuje nekonečne dlhá
vetva.

\paragraph{Dôkaz} $(T_1,\leq)$, $v \in T$. Označme $A_v = \{ u \in T | v <u
\}$ -- $A_v$ je množina tých vrcholov, ktoré ležia nad $v$. Majme vrchol $v$ a
nech $v_1, v_2, \ldots v_n$ sú bezprostrední nasledovníci $v$ (obr.
\ref{0326_successors}). $A_v = A_{v_1} \cup A_{v_2} \cup \ldots \cup A_{v_n}
\cup \{v_1, v_2, \ldots, v_n\}$. $A_v$ je nekonečná. Na základe uvedeného
tvrdenia aspoň jedna z množín $A_{v_i}$ musí byť nekonečná.

\begin{figure}[h]
	\centering\includegraphics{img/xx/2009-03-26.6.mps}
	\caption{Nasledovníci $v$}
	\label{0326_successors}
\end{figure}

\par Vyberieme si $x_0$ ako koreň (najmenší vrchol stromu $T$). $A_{x_0} =
T \backslash \{x_0\}$, $A_{x_0}$ je nekonečná. Vezmime postupne prvky $\{x_0,
x_1, \ldots , x_n, x_\{n+1\}\}$. Ku každému prvku viem vybrať nasledovný prvok a
počet nasledovníkov je nekonečný. (TODO: lepšia formulácia)

\paragraph{Veta (Herbrandova, variant 1)} Množina klauzúl $S$ nie je splniteľná
práve vtedy, keď ľubovoľnému úplnému sémantickému stromu pre množinu klauzúl
$S$ zodpovedá konečný uzavretý sémantický strom, t.j. ľubovoľná vetva úplného
stromu vedie do odmietajúceho vrchola.

\paragraph{Dôkaz} Predpokladajme, že množina klauzúl $S$ nie je splniteľná, nech
$T$ je úplný sémantický strom prislúchajúci $S$. $I_V$ je množina všetkých
literálov pripísaných vetve $V$ stromu $T$. Potom $I_V$ je interpretácia
množiny klauzúl $S$. Predpokladajme, že $S$ je nesplniteľná (teda nesplniteľná v
každej interpretácii). To znamená, že existuje nejaká základná inštancia klauzúl
$C'$ (pre klauzulu $C$), ktorá je v interpretácii $I_V$ odmietnutá. (základná
inštancia neobsahuje premenná, klauzuly sú konečné objekty). 
%Je odmietnutá vo vrchole, ktorý je konečne vzdialený od koreňa stromu.
\par Teraz potrebujeme zabezpečiť, aby strom bol konečný -- to musí, lebo ak by
existovala nekonečná vetva, bola by to interpretácia, ktorá by neodmietala
niektorú zo základných inštancií.

\par Obrátená implikácia. Ak množina klauzúl nie je splniteľná, existuje k
úplnému sémantickému stromu uzavretý konečný sémantický strom. Teraz
predpokladajme, že k úplnému sémantickému stromu pre množinu klauzúl
$S$ existuje konečný uzavretý sémantický strom (každá vetva stromu $T$ končí v
zamietajúcom vrchole), potom interpretácia $I_V$ odmieta každú základnú
inštanciu.


\paragraph{Veta (Herbrandova)}: Množina klauzúl $S$ nie je splniteľná $\iff$ keď
existuje konečná množina $S'$ základných inštancií klauzúl z $S$, ktorá nie je
splniteľná. FIXME

\paragraph{Dôkaz:} Predpokladajme, že množina klauzúl $S$ nie je splniteľná,
existuje $T'$ -- konečný uzavretý strom priradený ku stromu $T$.

\paragraph{Obrátene}: Predpoklajme, že existuije konečná $S'$ množina základných
inštancií klauzúl $S$ nie je splniteľná, $I'$, $S'$. Potom $S$ nie je
splniteľná, $I$ pre $S$. 

\par
Interpretácia $I$ obsahuje $I'$, teda keď $I'$ odmieta $S'$, potom $I$ odmieta
$S$. Každá klauzula je disjunkcia literálov. Ak existuje množina základných
inštancií, ktorá nie je splniteľná, potom celá množina nie je splniteľná.


\paragraph{Príklad} Majme množinu klauzúl $S=\{P(x), \neg P(f(a))\}$. Vezmime
množinu $S' = \{ P(f(a)), \neg P(f(a))\}$.

\paragraph{Príklad} Majme množinu klauzúl $S=\{\neg P(x) \lor Q(f(x),x),
P(g(b)), \neg Q(y,z) \}$. Táto množina nie je splniteľná -- opäť vezmime $$S' =
\{ \neg P(g(b)) \lor Q(f(g(b)), g(b)), P(g(b)), \neg Q(f(g(b)), g(b)) \}$$.

\paragraph{Príklad} Majme množinu klauzúl $S$:
$$
\begin{array}{l}
	\neg P(x,y,u) \lor \neg P(y,z,v) \lor \neg P(x,v,w) \lor P(u,z,w) \\
	\neg P(x,y,u) \lor \neg P(z,y,v) \lor \neg P(u,z,w) \lor P(x,v,w) \\
	P(g(x,y),x,y), P(x,h(x,y),y), P(x,y,f(x,y),\neg P(f(v),x,f(x)) \\
\end{array}
$$
\subparagraph{Návod} Zostrojte konečný sémantický strom.

\par Algoritmické metódy na zisťovanie nesplniteľnosti pracovali iba pre menšie
množiny klauzúl -- problém bol, že $S$ je množina disjunkcií, pokiaľ tam
dosadíme prvky z univerza, dostaneme konjunktívnu normálnu formu (prípadne
transformujeme na disjunktívnu normálnu formu a tú preverujeme). Ak máme $10$
formúl, potrebujeme preverovať $2^{10}$ konjunkcií -- zložitosť exponenciálne
rastie.

\paragraph{Príklad}
$$S=\{P(x), \neg P(a)\}, H_0 = \{a\}$$
$$S'=P(a) \land \neg P(a) = \square$$
$S$ nie je splniteľná.

\paragraph{Príklad}
$$S=\{P(a), \neg P(x) \lor Q(f(x)), \neg Q(f(a)) \}, H_0 = \{a\}$$
$ S_0'=P(a) \land (\neg P(a)\lor Q(f(a))) \land \neg Q(f(a)) = \\
P(a)\land \neg P(a) \land \neg Q(f(a)) \lor P(a) \land Q(f(a)) \land \neg Q(f(a)) = 
\\ \square \lor \square = \square$.

\paragraph{Pravidlá}
Nech $S$ je množina klauzúl.

\subparagraph{Pravidlo tautológie.} Vynecháme z $S$ všetky tautologické klauzuly.
Množina $S'$, ktorá nám zostane, nie je splniteľná práve vtedy, keď $S$ nie je
splniteľná.

\subparagraph{Pravidlo jednoliterálnych klauzúl.} Nech $L$ je nejaký literál.
Vynechajme z $S$ všetky klauzuly, ktoré obsahujú literál $L$. Nech $S'$ sú
klauzuly, ktoré nám zostanú po vynechaní. Môžu nastať dva prípady:

\begin{enumerate}
	\item $S' = \emptyset$, tak potom množina klauzúl $S$ je splniteľná --
	Stačí zobrať model, ktorý obsahuje $L$.
	\item $S' \neq \emptyset$. Vezmem si literál $\neg L$ a vynechám z
	množiny $S'$ všetky klauzuly, ktoré obsahujú $\neg L$, pričom dostanem
	množinu klauzúl $S''$. Ak sa $S'$ nachádza $\neg L$, po vynechaní
	namiesto nej zostane $\square$.
\end{enumerate}

$S$ nie je splniteľná $\iff$ $S''$ nie je splniteľná.

\subparagraph{Pravidlo čistých literálov} Literál $L$ základnej klauzuly z $S$
budeme nazývať \emph{čistým}, ak literál $\neg L$ sa nevyskytuje v žiadnej
základnej klauzule $S$. Vezmem si z $S$ literál $L$, ktorý je čistý, a
vynecháme z $S$ všetky základne klauzuly obsahujúce $L$. $S'$ je množina, ktorá
nám zostala. $S$ nie je splniteľná $\iff$ $S'$ nie je splniteľná.

\subparagraph{Pravidlo rezu} Ak množinu $S$ vieme vyjadriť v tvare $(A_1 \lor L)
\land \ldots \land (A_m \lor L) \land (B_1 \lor \neg L) \land \ldots \land (B_n
\lor \neg L) \land R$, pričom v $A_i$, $B_j$, sa $R$, $L$, $\neg L$, $i=1..m$,
$j=1..n$ už nevyskytujú.

$$S_1 = A_1 \land A_2 \land \ldots \land A_m \land R$$
$$S_2 = B_1 \land B_2 \land \ldots \land _n \land R$$

$S$ nie je splniteľná $\iff$ $S_1 \lor S_2$ nie je splniteteľná.

\paragraph{Dôkaz}
\begin{enumerate}
	\item Keď $S$ nie je splniteľné, tak ani $S_2$ nie je splniteľné --
	tautológie sú splnené pre ľubovoľné interpretácie. To, čo zostane,
	hovorí, či množina je alebo nie je splniteľná.
	\item Zostáva nám $S'$, ktoré vzniklo vynechaním $L$ z $S$. Ak
	$S'=\emptyset$, každá klauzula z $S$ obsahuje $L$ (je tvaru $L \lor
	\square$), stačí nám ohodnotenie, kedy $L$ je pravdivé, čo zabezpečí
	pravdivosť celého výroku. 
	\par
	Keď $S'$ nie je prázdna, vytvoríme $S''$ tak,
	že vynecháme všetky klauzuly obsahujúce $\neg L$. $S'$ nie je splniteľná
	práve vtedy, keď $S$ nie je splniteľná. Predpokladajme, že $S''$ nie je
	splniteľná a $S$ je splniteľná. Model $\mathcal{M}$ musí obsahovať $L$ a
	$S$ platí na $S''$.
	TODO
	V konjunktívnej normálnej forme musia byť platné doplnky $\neg L \lor
	\square$. Potom $S''$ je splniteľná, čo je spor.

	\par Obrátene, nech $S$ nie je splniteľná. Z toho vyplýva, že $S''$ nie
	je splniteľná. Nech $S''$ je splniteľná. Potom existuje model
	$\mathcal{M}''$ pre $S''$. Ak tomuto modelu pridám $L$, potom
	$\mathcal{M}'' \cup L$ je model pre $S$, čo je spor.

	\item Pravidlo čistých literálov -- Máme čistý literál $L$, všetky
	klauzuly, ktoré ho obsahujú, vyhodíme z $S$ a dostaneme $S'$. $S$ nie je
	splniteľná práve vtedy, keď $S'$ nie je splniteľná. Predpokladajme, že
	$S'$ nie je splniteľná, teda pre žiadnu interpretáciu $I'$ neplatí. $S\
	\subseteq S$. Musí platiť, že $I' \subseteq I$.

	\par Obrátene -- predpokladajme, že $S$ nie je splniteľná. Potom aj $S'$
	nie je splniteľná. Predpokladajme, že $S'$ je splniteľná. Potom musí mať
	model $\mathcal{M}'$, ktorý jej vyhovuje. Keď si vytvorím model
	$\mathcal{M}=\mathcal{M}'\cup L$, potom $\mathcal{M}$ je modelom v $S$.

	\item Pravidlo rezu. Mám množinu $S$, popísaným spôsobom som zostrojil
	$S_1$ a $S_2$. Chceme ukázať, že $S$ nie je splniteľná $\iff$ $S_1 \lor
	S_2$. Predpokladajme, že $S$ je splniteľná. Potom $S_1 \lor S_2$ je
	splniteľná, teda aspoň jeden člen disjunkcie je splniteľný. Keď pridám
	modelu $S_1$ negované $B$, dostávam model pre $S$. Taktiež, keď je
	splnená $S_2$, má nejaký model $\mathcal{M_2}$. Keď k nemu pridám $B$,
	dostávam model pre $S$.

	\par Obrátene, predpokladajme, že $S_1$ alebo $S_2$ nie je splniteľná. Z
	toho vyplýva, že $S$ nie je splniteľná. Pre spor predpokladajme, že $S$
	je splniteľná. Pozrime sa, ako vyzerá $S$ -- všetky konjunktívne členy
	musia byť splnené a teda buď platí $\neg L$ alebo $L$. Vždy dospievam do
	sporu.
\end{enumerate}

\paragraph{Záver} Najvšeobecnejšie je pravidlo rezu, ktoré pokrýva aj všetky
ostatné prípady.

\paragraph{Príklad} Majmne množinu $S=(P \lor Q \lor \neg R) \land (P \lor \neg
Q) \land \neg P \land R \land U$. Ukážte, že množina klauzúl nie je splniteľná.
Použijeme pravidlo 2 s $\neg P$. Dostávam $Q\land\neg R) \land \neg  \land R
\land U$. Na túto množinu klauzúl použijem pravidlo 2 s $\neg Q$. Dostávam $\neg
R \land R \land U$, čo je $\square \land U \square$.




\paragraph{Príklad} (7) $S=(P \lor Q) \land \neg Q \land ( \neg P \lor Q \lor \neg
R)$. Ukážte, že množina klauzúl $S$ je splniteľná.
\par Pravidlo 2 s $\neg Q$: $(P \lor Q) \land ( \neg P \lor Q \lor \neg R)$ \\
$P \land (\neg P \lor \neg R)$ \\
$(P \land \neg P) \lor (P \land \neg R)$ (distributívnosť) \\
$P \land \neg L, I=\{P,\neg Q, \neg R\}$

\paragraph{Príklad} (8) $S=(P\lor \neg Q) \land (\neg P \lor Q) \land (Q \lor
\neg R) \land (\neg Q \lor \neg R)$. Zistite, či množina je alebo nie je
splniteľná.
\par vytvoríme si mnonžiny podľa pravidla 4: \\
$$(P\land \neg Q) \land (Q \lor \neg R)  \land (\neg Q \lor \neg R)$$
$$(\neg P \lor Q) \land (Q \lor \neg R) \land (\neg Q \lor \neg R)$$

\par
$$S_1 = \neg Q \land (Q \lor \neg R) \land (\neg Q \lor \neg R)$$
$$S_2 = Q \land (Q \lor \neg R) \land (\neg Q \lor \neg R)$$

($S_1$ a $S_2$ sú množiny rezu)
Na množiny môžem použiť pravidlo 2 -- s $Q$ a $\neg Q$. Dostávam $\neg R \lor
R$, čo je splniteľná klauzula.

\paragraph{Príklad} (9) $S=(P \lor Q) \land (P \lor \neg Q) \land (R \lor Q)
\land (R \lor \neg Q)$. $P$ je čistý literál; použijeme pravidlo 3:
$S': (R \lor Q) \land (R \lor \neg Q)$. \\
$\fbox{P,R}$ \\

\section{Metóda rezolvent pre výrokovú logiku}

\paragraph{Definícia} Majme $C_1$ a $C_2$ -- ľubovoľné klauzuly, $L_1$ je
literál z $C_1$ a je kontrárny literálu $L_2$, ktorý sa vyskytuje v $C_2$. Ak z
$C_1$ vynecháme $L_1$ a z $C_2$ vynecháme $L_2$. Ak z tých častí, ktoré
zostávajú, zostrojíme disjunkciu $C_1 \lor C_2$, táto sa nazýva
\emph{rezolventa} $C_1$ a $C_2$.

\paragraph{Príklad}
$$
\begin{array}{ll}
C_1: P \lor R, & C_2: \neg P \lor Q \\
P \in C_1 & C_1': R \\
P \in C_2 & C_2': Q \\
\end{array}
$$

\paragraph{Príklad}
$$C_1: \neg P \lor Q\lor R$$
$$C_2: \neg Q \lor S$$
Kontrárna dvojica $\neg Q$, $Q$; rezolventa je teda:
$$C_1, C_2: \neg P \lor R \lor S$$

\paragraph{Príklad}

$$C_1: \neg P \lor Q$$
$$C_2: \neg P \lor R$$

\paragraph{Veta} Nech $C_1$ a $C_2$ sú klauzuly, $C$ je rezolventa klauzúl
$C_1$, $C_2$. Potom $C$ je logickým dôsledkom klauzúl $C_1$ a $C_2$
$$A_1, A_2, \ldots, A_n, B$$
$$
\begin{array}{ll}
A_1 \land A_2 \land \ldots \land A_n \implies B & \equiv 1	\\
\vdash A_1 \land A_2 \land \ldots \land A_n \implies B & \\
A_1 \land A_2 \land \ldots \land A_n \land \neg B & \equiv 0	\\
\end{array}
$$
\paragraph{Dôkaz} $C_1$, $C_2$, $C$


$$
\begin{array}{lll}
C_1 &= &L \lor C_1' \\
C_2 &=& \neg L \lor C_2' \\
C   &=& C_1' \lor C_2'	\\
\end{array}
$$
$C_1$ a $C_2$ sú pravdivé a klauzula $C$ je logickým dôsledkom $C_1$ a $C_2$.

\paragraph{Poznámka} Ak $C_1$, $C_2$ sú jednotkové klauzuly, potom musia tvoriť
kontrárnu dvojicu $C_1: L$, $C_2: \neg L$ a rezolventou je prázdna klauzula --
$C: \square$


\paragraph{Definícia} Nech $S$ je množina klauzúl.
\emph{Rezolvenčným odvodením} $C$ z množiny $S$, nazývame konečnú
postupnosť klauzúl $C_1, C_2, \ldots, C_n (C)$ takú, že pre každé $i=1..n$ platí
$C_i$ je buď z $S$ alebo $C_i$ je rezolventa niektorých klauzúl $C_j$, $C_k$,
$j, k < i$.

\par Ak $C$ je prázdna klauzula, takémuto odvodeniu hovoríme \emph{zamietniue
odvodenia} a množina klauzúl $S$ nie je splniteľná.

\paragraph{Poznámka} Ak mám množinu $S$ a klauzulu $C$, hovoríme, že $C$ môžem
získať z $S$, ak existuje odvodenie $S, C_1, \ldots C_m (C)$.

\paragraph{Príklad} 
\begin{enumerate} 
	\item $ \neg P \lor Q $
	\item $ \neg Q $
	\item $ P $
	\item $ \neg P $
\end{enumerate}
Použitím 3 a 4 dostávame prázdnu klauzulu, teda množina klauzúl $S$ nie je
splniteľná.

\paragraph{Príklad}
S:
\begin{enumerate}
	\item $ P \lor Q $
	\item $ \neg P \lor Q $
	\item $ P \lor \neg Q $
	\item $ \neg P \lor \neg Q $
	\item $ Q $
	\item $ P $
\end{enumerate}

%\begin{verbatim}
%                     0
%		    / \
%		   /   \
%		  Q  neg Q
%		 /\       /\
%		/  \     /  \
%	       /    \   /    \
%	      P  neg P P    neg P
%	      (listy su jednotlive vyroky)
%\end{verbatim}
%
\begin{figure}[h]
	\centering\includegraphics{img/xx/2009-04-02.1.mps}
	\caption{Výrazy}
\end{figure}

$\neg A \lor V, A \lor C \vdash B \lor C$
$$A \implies B, \neg A \implies C \vdash \neg B \implies C$$
\begin{enumerate}
	\item $$\vdash (A \implies B) \implies (\neg B \implies \neg A)$$
	\item $$ \vdash (\neg B  \implies \neg A) \implies ((\neg A \implies C)
	\implies (\neg B \implies C))$$
\end{enumerate}
$$ A \implies B, \neg C \implies C \vdash \neg V \implies C$$
$$ A, A\implies B \vdash B$$

$$A \lor \square, \neg A \lor B \vdash$$
Pravidlo \emph{modus ponens} a pravidlo rezolventy sú ekvivalentné.

\paragraph{Príklad}
$$C_1: P(x) \lor Q(x)$$
$$C_2: \neg P(f(x)) \lor R(x)$$



\section{Substitúcia a unifikácia}

$$
\begin{array}{ll}
C_1  = P(x) \lor Q(x),& C_2 = \neg P(f(x)) \lor R(x)\\
C_1:  x \cdot f(a),&  C_2 \cdot : a\\
C_1':  P(f(x)) \lor Q(f(a)),& C_2':\neg P(f(a)) \lor R(a)\\
\end{array}
$$

Teraz môžeme definovať rezolventu; bude to $Q(f(a)) \lor R(a)$. Mohli by sme
postupovať aj takto:

$$
\begin{array}{ll}
C_1: & x: f(x) \\
C_1^*: & P(f(x)) \lor Q(f(x)) \\
C_2: & \neg P(f(x)) \lor R(x)\\
 & P(f(x))\\
 & \neg P(f(x)))\\
C_3: & Q(f(x)) \neg R(x)\\
\end{array}
$$

Formulu $C_3$ nazývam rezolventou pre klauzulu $C_1$ a $C_2$. Zmyslom
substitúcie a unifikácie je umožniť operácie, na ktoré sme boli zvyknutí vo
výrokovej logike -- tu bude situácia technicky náročnejšia.

\paragraph{Definícia} Pod \emph{substitúciou} rozumieme konečnú množinu tvaru:
$\{t_1/v_1, \ldots, t_n/v_n\}$, kde každé $v_i$ je premenná a všetky $v_i,
i=1..n$ sú navzájom rôzne. Každý term $t_i, i=1..n$ je rôzny od $v_i, i=1..n$.
ak $t_1, \ldots, t_n$ sú základné inštancie (teda termy bez premenných), tak
substitúciu nazývame \emph{základná substitúcia}. Ak množina je prázdna, tak
substitúciu budeme označovať $\varepsilon$. Na označovanie substitúcií budeme
používať grécke písmená.

\paragraph{Príklad} $\{ f(z)/x; y/z\}$, $\{ a/x, g(y)/y, f(g(y))/z\}$.

\paragraph{Definícia} Nech $\theta$ je ľubovoľná substitúcia a $E$ je nejaký
výraz. Nech $\theta = \{ t_1/v_1, \ldots t_n/v_n\}$. Potom $E\theta$ označuje
výraz, ktorý vznikne tak, že súčasne vo výraze $E$ nahradíme premennú $v_i$
termom $t_i, i=1..n$.

\paragraph{Príklad} Majme $\theta=\{a/x, f(b)/y, c/z\}$, $E = P(x, y, z)$. Potom
$E\theta = P(P(a, f(b), c)$.

\par Ďalšia operácia, ktorú budeme potrebovať, je operácia skladania
substitúcií.

\paragraph{Definícia} Majme $\theta = \{t_1/x_1, \ldots t_n/x_n\}$ a $\lambda =
\{ u_1/y_1, \ldots u_m/y_m \}$. Zloženie (kompozícia) substitúcií, označíme
$\theta \circ \lambda$:
$$\{t_1 \lambda/x_1, \ldots t_n \lambda/x_n, u_1/y_1m \ldots u_m/y_m \}$$ také,
že vyradí všetky členy $t_i\lambda/x_i$, pre ktoré platí, že $t_i \lambda = x_1$,
všetky $u_i/y_1$, pre ktoré $y_i \in \{x_1, x_2, \ldots x_n\}$.

\paragraph{Príklad} 
$$\theta \{t_1/x_1, t_2/x_2\} = \{ f(y)/x, z/y\}$$
$$\lambda = \{ u_1/y_1, u_2/y_2, u_3/y_3\} = \{ a/x, b/y, y/z\}$$

$$ \theta \circ \lambda = \{ t_1 \lambda / x_1, t_2\lambda/x_2, u_1/y_1,
u_2/y_2, u_3/y_3\} = \\ 
\{ f(b)x, y/y, a/x, b/y, y/z\}$$
$$ \theta \circ \lambda = \{ f(b)/x, y/z\}$$

\paragraph{Úloha} Dokážte, že skladanie substitúcií je asociatávna operácia.
Tiež dokážte, že $\varepsilon \circ \theta = \theta \circ \varepsilon$.
\par Zoberieme $\theta, \lambda, \mu$, potom platí $ \theta \circ(\lambda \circ
\mu) = (\theta \circ \lambda) \circ \mu$.

\paragraph{Definícia} Substitúciu $\theta$ nazveme \emph{unifikátorom} pre
množinu výrazov $E_1, E_2, \ldots, E_n$, ak platí $E_1\theta = E_2\theta =
\cdots = E_n\theta$. Množinu nazveme \emph{unifikovateľnou}, ak pre ňu existuje
unifikátor.

\paragraph{Príklad} Majme množinu $\{ P(a,y), P(x,f(b)\}$, $\theta = \{a/x,
f(b)/y\}$ -- unifikátor pre túto množinu. 

\paragraph{Definícia} Majme množinu výrazov $\{ E_1, E_2, \ldots, E_n\}$,
$\theta$ je ľubovoľný unifikátor pre množinu uvedených výrazov. Unifikátor
$\sigma$ pre množinu výrazov $\{E_1, E_2, \ldots, E_n\}$ nazveme
\emph{najvšeobecnejší unifikátor}, ak pre $\theta$ platí $\theta = \sigma \circ
\lambda$, kde $\lambda$ je vhodná substitúcia.

\par Uvažujme výrazy $P(a)$, $P(x)$ Pozerajme sa na ne ako na konečnú
postupnosť symbolov -- odlišujú sa akurát v treťom symbole. Toto je prvá
diferencia. Vo všeobecnosti môže byť týchto výrazov $n$.

\paragraph{Definácia} Nech $W$ je neprázdna množina výrazov. \emph{Diferenčnú
množinu}
pre množinu výrazov $W$ dostávame tak, že na prvú pozíciu (zľava) vypíšem všetky
tie výrazy, ktoré sa líšia. 

\paragraph{Príklad} Majme množinu $W = \{P(x,\underline{f(y,z)},
P(x,\underline{a}), P(x,\underline{g(h(k(x)))}) \}$
Diferenčnou množinou bude $D= \{ f(y,z), a, f(h(k(x)))\}$. 

\subsection{Unifikačný algoritmus}

Kroky:
\begin{enumerate}
	\item $k=0, W_0 = W, \sigma_0 = \Sigma$
	\item Ak $W_k$ je jednotková klauzula, algortimus zakončí svoju činnosť
	$\sigma_k$ je najvšeobecnejší unifikátor. V opačnom prípade nájdeme $D_k$
	(diferenčnú množinu) pre $W_k$.
	\item Ak existujú také premenné $v_k$ a $t_k$ v $D_k$, že $v_k$ je
	premenná, ktorá sa nevyskytuje v $t_k$, tak prejdeme ku kroku 4. V
	opačnom prípade algoritmus zakončuje svoji činnosť. Množina $W$ nie je
	unifikovateľná.
	\item Nech $W_{k+1} = W_k \{t_k/v_k\}$, $\sigma_{k+1} = \sigma_k \circ
	\{t_k/v_k\}$ ($W_{k+1} = W_k \sigma_{k+1}$)

	\item Vypíšeme hodnoty pre $k+1$ a prejdeme ku kroku 2.
\end{enumerate}


Ak je množina unifikovateľná, vždy existuje najvšeobecnejší unifikátor.

\paragraph{Príklad}
Nájdite najvšeobecnejší unifikátor pre množinu: $$W=\{ P(a,x,f(g(y))),
P(z,f(z),f(u)), \}$$

\begin{enumerate}
	\item $\sigma_0 = \varepsilon, W_0 = W$. Pretože $W_0$ nie je jednotková
	klauzula, $\sigma_0$ nie je najvšeobecnejší unifikátor.
	\item Musíme zostrojiť diferenčnú množinu $D_0 = \{a, z\}$. Existuje
	premenná $v_0 = z$, ktorá nie je obsiahnutá v terme $t$. $t_0 = a$.
	\item $$\sigma_1 = \sigma_0 \circ \{ t_0/v_0 \} = \epsilon \{a/z\} =
	\{a/z\}$$ $W_1 = W_0\{ t_0/v_0 \} = \{P(a,x,f(g(y)),
	P(z,f(z),f(u))\}\{a/z\} =  \\
	= \{P(a,x,f(g(y))), P(a, f(a), f(u))\}$

	\item $W_1$ nie je jednotková klauzula. Zostrojíme diferenčnú množinu
	$D_1$ pre $W_1$. $D_1 = \{x_1, f(a)\}$

	\item Z $D_1$ dostávame $v_1 = x$ a $t_1 = f(a)$.
	\item $$\sigma_2 = \sigma_1 \circ \{ t_1/v_1\}= \{a/z\} \circ \{ f(a)/x\}
	= \{a/z, f(a)/z \}$$
	$ W_2 = W_1 \{t_1/v_1\} = \{P(a,x,f(g(y))), P(a,f(a),f(u)\} \{ f(a)/x
	\} = \\
	= \{ P(a,f(a), f(g(y))), P(a,f(a), f(u), \}$. $W_2$ opäť nie je
	jednotková klauzula -- musíme vytvoriť diferenčnú množinu.

	\item $W_2$ nie je jednotková klauzula -- vytvárame diferenčnú množinu
	$D_2$ pre $W_2$. $D_2 = \{ g(y), u \}$. $v_2 = u$, $t_2 = g(y)$.

	\item $$\sigma_3 = \sigma_2 \circ \{t_2/v_2\} = \{a/z, f(a)/x\} \circ
	\{g(y)/u\}$$.
	$W_3 = W_2 \{ t_x/v_2\} = \{ P(a,f(a), f(g(y))), P(a,f(a),
	f(u))\}\{g(y)/u\} = \\
	= \{ P(a, f(a), f(g(y))), f(a, f(a), f(g(y)) \}$.
	$W_3$ je jednotková klauzula.
	$$ \sigma_3 = \{ a/z, f(a)/x, g(y)/u\}$$. $\sigma_3$ je najvšeobecnejší
	unifikátor pre množinu klauzúl $W$.
\end{enumerate}

\paragraph{Príklad}

\paragraph{Príklad} Zistitte, či je unifikovateľná množina $W=\{Q(f(x),g(x)),
Q(y,y)\}..$


	$$y\circ \sigma_0 \varepsilon \qquad \mbox{a} \qquad W_0 = W$$
$W_0$ nie je jednotková klauzula. Nájdeme diferenčnú množinu $D_0$ pre $W_0$:
$$D_0 = \{ f(a)/y \}$$
$$v_0, t_0 = f(a)$$

Nech $\sigma_1 = \sigma_0 \circ \{ t_1/\sigma_0\} = \varepsilon \circ  \{
f(a)y\}$

$W_1 = W_0 \{ t_0/\sigma_0\} = \{ Q(f(x),g(x)), Q(f(a),f(a))\}$

$W_1$ nie je jednotková klauzula, zostrojujeme $D_1$, diferenčnú množinu
pre $W_1$.


$$D_1 = \{ g(x), f(a)\}$$ nemáme prvok, ktorý by bol premennou. Algoritmus
ukončí svoju činnosť a $W$ nie je unifikovateľná.

\par Pri zisťovaní unifikovateľnosti vždy vytvárame množiny $W$ tvaru:
$$W\sigma_0, W\sigma_1, \ldots, W\sigma_k, \ldots$$, pričom v každom kroku sa
zmenší počet premenných o 1. Po konečnom počte krokov sa musí zastaviť.


\paragraph{Veta (unifikačná):} Ak $W$ je konečná neprádzna unifikačná množina
výrazov, tak unifikačný algoritmus vždy zakončuje svoju činnosť na druhom kroku
a posledné $\sigma_k$ bude najvšeobecnejší unifikátor.

\paragraph{Dôkaz:} $W$ je unifikovateľná množina, $\Theta$ označuje jej
ľubovoľný unifikátor. Indukciou ukážeme, že pre každé $k$ existuje taká
substitácia $\lambda_k$, že $\Theta = \sigma_k \circ \lambda_k$

\subparagraph{Báza indukcie} Nech $k = 0$. Máme ukázať, že $\Theta = \sigma_0
\circ \lambda_0$. V tomto prípade $\sigma_0 = \varepsilon$ a $\lambda_0 =
\Theta$.

\subparagraph{Indukčný krok} Indukčný predpoklad: $\Theta = \sigma_k \circ
\lambda_k$ $(0 \leq k \leq n)$. $W_n = W\sigma_n$. Ak $W_n$ je jednotková
klauzula, tak algoritmus zakončuje svoju činnosť na druhom kroku a $\sigma_n$ je
najvšeobecnejší unifikátor pre $W$. 
\par Nech $W_n$ nie je jednotková množina. Potom hľadám diferenčnú množinu $D_n
$ pre množinu $W_n$. Keď $D_n$ je diferenčná množina pre $W_n$, tak vo $W_n$
musí existovať premenná -- označme ju $v_n$. Ďalej nech $t_n$ je ľubovoľný výraz
z $D_n$ rôzny od $v_n$.

$$\Theta = \sigma_n \circ \lambda_n$$ 
$\Theta$ je unifikátor. Diferenčnú množinu $D_n$ unifikuje substitúcia
$\lambda_n$. $v_n \lambda_n = t_n \lambda_n$.

\par Ak by premenná $v_n$ bola obsiahnutá v $t_n$. Predpokladali sme, že $v_n$ a
$t_n$ sú rôzne a teda $v_n$ sa nemôže vyskytovať v $t_n$.

\par Prejdeme ku kroku 4 -- množine $W\sigma_{n+1}$ (množina $W_{n+1}$). Platí,
že $\sigma_{n+1} = \sigma-n \circ \{ t_n/v_n \}$. $\Theta = \sigma_{n+1} \circ
\lambda-{n+1}$.

\par Nech $\lambda_{n+1} = \lambda_n - \{t_n\lambda_n/v_n\}$. 

$$t_n\lambda_{n+1} = t_n (\lambda_n - \{t_n\lambda_n / v_n\}) -  t_n\lambda_n$$
$$\{ t_n / v_n\} \circ \lambda_{n-1} = \{t_n \lambda_{n+1}\}/v_n \cup
\lambda_{n+1} = \{t_n\lambda_n /v_n\} \cup \lambda_{n+1} = \{t_n\lambda_n/v_n\}
\cup \{ \lambda_n - \{ t_n\lambda_n/v_n\} = ??? \{ t_n/v_n \} \circ
\lambda_{n+1} = \lambda_n$$

$$\Theta = \sigma_n \circ \lambda_n$$.

\section {Metóda rezolvent pre logiku 1. rádu}

\paragraph{Definícia:} Nech $C$ je klauzula; ak dva alebo viac literálov (s
rovnakým znakom predikátu) Majú najvšeobecnejší unifikátor $\sigma$, tak
$C\sigma$ sa nazýva \emph{spojením $C$}. Ak $C\sigma$ je jednotková klauzula,
tak $C\sigma$ nazývame \emph{jednotkovým spojením $C$}.


\paragraph{Príklad} Uvažujme klauzulu $C$, ktorá vyzerá nasledovne: $C = \{ P(x)
\lor P(f(y)) \lor \neg Q(x)\}$. 

\paragraph{Definícia:} Nech $C_1$ a $C_2$ sú klauzuly (nazývame ich
predpoklady), ktoré nemajú spoločné premenné. Nech $L_1 \in C_1$ a $L_2 \in
C_2$ sú dva literály . Ak $L_1$ a $\neg L_1$ majú najvšeobecnejší unifikátor
$\sigma$, tak:
$$(C_1\sigma - L_1\sigma) \lor (C_2\sigma - L_2\sigma)$$

sa nazýva binárnou rezolventou ..... $C_1$ a $C_2$. Literály $L_1$ a $L_2$ nazývame
nadbytočné a môžu sa odstrániť.

\paragraph{Príklad} $C_1 = p(x) \lor Q(x), C_2 = \neg P(x) \lor R(x)$, $C_1$,
$C_2$ -- predpoklady. $C_2 = \neg P(a) \lor R(y)$.

$$L_1 = P(x), L_2 = \neg P(x)$$
$$\neg L_2 = P(x)$$
$$L_1, \neg L_2, \sigma = \{a/x\}$$

Binárna rezolventa $C_1$ a $C_2$ $(C_1\sigma - L_1\sigma) \lor (C_2\sigma -
L_2\sigma) = (\{P(a),Q(a))\} - \{P(a)\}) \lor (\{\neg P(a),R(y)(\}-\{\neg
P(a)\}) = \{Q(a)\lor R(y)\}$

\paragraph{Definícia} Rezolventou z predpokladov $C_1$ a $C_2$ je jedna z
nasledujúcich rezolvent:
\begin{enumerate}
	\item Binárna rezolventa $C_1$ a$C_2$
	\item Binárna rezolventa $C_1$ a spojenia $C_2$
	\item Binárna rezolventa spojenia $C_1$ a $C_2$
	\item Binárna rezolventa spojenia $C_1$ a spojenia $C_2$
\end{enumerate}

\paragraph{Príklad} $$C_1 = P(x) \lor P(f(y))\lor Rg(y))$$
$$ C_2 = \neg P(f(g(a)) \lor Q(b)$$
Spojenie pre $C_1$ vyzerá ako: $C_1': P(f(y)) \lor R(g(y))$.
Binárna rezolventa $C_1' a C_2$ bude vyzerať takto: $R(g(g(a))) \lor Q(b)$ --
rezolventa $C_1$ a $C_2$.


\section{Opakovanie}
	Ak množina klauzúl nie je splniteľná, potom metódou rezolvent z nej vždy
	môžeme dostať prázdnu klauzulu (a ak sa táto dostane množiny klauzúl,
	tak formula nie je splniteľná v žiadnej interpretácii). Ak máme nejakú
	klauzulu $C$, $C\sigma$ sme nazývali spojením klauzuly $C$. Definícia
	binárnej rezolventy. 

\par  Metódu rezolvent zaviedol roku 1965 Robinson, je efektívnejšia ako obe
varianty Herbrandovej metódy. 

\par \{ sleep...\}

\paragraph{Úplnosť metódy rezolvent}

\paragraph{Príklad} Majem množinu klauzúl $S$:
\begin{enumerate}
	\item $P$
	\item $\neg P\lor Q$
	\item $\neg P \lor \neg Q$
\end{enumerate}
Tejto množine klauzúl zodpoveda uzavretý sémantický strom.
Prislúchajúca herbrandovská báza je $\{P, Q\}$ (na tabuľu sa kreslí sémantický
strom pre $P$ a $Q$, usilovný čitateľ si ho isto domyslí). Každá vetva sa končí
odmietajúcim vrcholom, žiadna z tých interpretácií, ktoré končia v listoch, nie
je splniteľná. Tomuto stromu môžeme priradiť uzavretý podstrom (označíme ho
$T'$,  má odseknuté vetvy na miestach, kde sú podstromy odmietajúce)


\begin{verbatim}
            T
            /\
         P /  \ \neg P
       Q /\ nQ Q/\ nQ 
\end{verbatim}

\begin{verbatim}
            T'
            (1)
            /\
      (2)P /  \ \neg P
   (4) Q /\ nQ x (5)
       x   x
\end{verbatim}

\par $\neg P$ -- rezolventa $(4)$, $(5)$, $\neg P \cup S$.  $S\cup \{ \neg P
\}\cup \{ \square \}$.

\par Vznikli nám teda klauzuly.
$$(4) \neg P\qquad (2) (3)$$
$$(5) \Box\qquad (4)(1)$$

\par Strom sa po každej aplikácii pravidla postupne skracuje. 

\paragraph{Lema} Nech $C_1'$ a $C_2'$ sú inštancie $C_1$ resp. $C_2$ (v uvedenom
poradí). Ak $C'$ je rezolventa $C_1'$ a $C_2'$, tak potom existuje rezolventa
$C$ klauzúl $C_1$ a $C_2$, že $C'$ je inštancia $C$. 

\paragraph{Dôkaz} Ak je treba, premenujeme premenné $C_1$ a $C_2$. Nech $L_1'$ a
$L_2'$ sú literály, ktoré môžeme vynechať (sú nadbytočné). ďalej nech platí:
$$C' = (C_1' \nu - L_1'\nu) \cup ( C_2'\nu - L_2'\nu)$$
Pričom $\nu$ je najvšeobecnejší unifikátor pre $L_1'$ a $\neg L_2'$. $C_1'$,
$C_2'$ sú inštancie $C_1$ a $C_2$, a teda existuje substitúcia $\Theta$ taká, že
platí:

$$ C_1' = C_1 \Theta $$
$$ C_2' = C_2 \Theta $$

(Pozn.: $C_1$ a $C_2$ nemajú spoločné premenné). $L^1_i, L^2_i, \ldots, L^{r_i}_i,
i=1,2$ sú literály, ktoré v $C_1$ zodpovedajú $L_i'$, teda $L^1_i \Theta = L^2_i
\Theta = \cdots = L^{r_i}_i\Theta = L'_i (i=1,2)$. $D_i > 1$ dostaneme
najvšeobecnejší kvantifikátor $\lambda_i (i=1,2)$ pre $\{ L^1_i, L^2_i, \ldots,
L^{r_i}_i\}, L_i = L^1_i \lambda_i (i=1,2)$. $\lambda_i$ je najvšeobecnejší
unifikátor, tak pre vhodnú substitúciu $\xi$ platí:
$$ L_i' = L^1_i \Theta = L^1_i (lambda_i \circ \xi) = (L^1_i\lambda_i)\xi =
L_i \xi$$

$$L_i\xi = L_i'$$

$L_i$ .. spojení $C_i\lambda_i$, pre $C_i$, ak $r_i = 1$, ak $r_i = 1$, potom
$\lambda_i = \Sigma$, $L_i = L_i^1\lambda_i$.

$$ \lambda = \lambda_1 \cup \lambda_2$$
$$L_i' = \mbox{...} L_i$$

$L'_i, \neg L_2'$ -- unifikovateľné.
$L-1, \neg L_2$ -- unifikovateľné.

Označme $\sigma$ najvšeobecnejší unifikátor pre $L_1'$ a $\neg L_2'$.

$C= ((C_1\lambda)\sigma = L_1\sigma) \cup
((C_2\lambda_2)\sigma-L_2\sigma) = ((c_1\lambda)\sigma - (\{L^1_+, L^2_1,
\ldots, L^{r_i}_1 ...
= C_1(\lambda\circ \sigma) = \{ L^1_1, \ldots L^{r_1}_1 \} (\lambda\circ\sigma)
\cup C_2)(\lambda\circ\sigma) - \{L^1_2, \ldots
L^{r_i}_1\}(\lambda\circ\sigma)$

\par $C$ -- rezolventa $C_1$ a$C_2$, $C'$ je substitúcia $C$:
$C = (C_1' \nu = L_1' \nu) \cup (C_2'\nu = L_2'\nu) = (C_1\Theta)\nu -
(\{L^1_1, \ldots L^{r_i}_1\}\Theta)\nu)\cup ((C_2\Theta)\nu - \{L^1_2, L^2_2,
\ldots, L^{r_i}_2\}\Theta )\nu) = 
C_1(\Theta\circ\nu) - \{L^1_1, \ldots, L^{r_i}_1\} \Theta\circ\nu) \cup
(C_2(\Theta\circ\nu) - \{L^1_2, \ldots L^{r_i}_2\} \Theta\circ\nu)
$

$\lambda \circ \sigma$ jke všeobecnejšia ako $\theta \circ \nu$.


\paragraph{Veta (úplnosť metódy rezolvent)} Množina klauzúl $S$ nie je
splniteľná práve vtedy, keď existuje odvodenie prázdnej klauzuly $\Box$ z $S$.

\paragraph{Dôkaz} Predpokladajme, že z $S$ existuje odvodenie prázdnej klauzuly
$\square$. $R_1, R_2, \ldots R_n$ sú všetky rezolventy v odvodení (medzi nimi
niekde bude aj $\square$). Zoberiem $C_1$, $C_2$ -- ľubovoľné klauzuly z $S$ a
$C_1$ a $C_2$ bude príslúchať nejaká rezolventa. Ak sú klauzuly splniteľné, je
splniteľná aj rezolventa. To znamená, že $C_1$ a $C_2$ nemôžu byt splniteľné a
teda nemôže byť množina klauzúl (klauzuly, ktorých rezolventou je prázdna
klauzula, nebudú splniteľné nikdy). 

\subparagraph{Obrátené tvrdenie} Prepokladajme, že množina klauzúl $S$ nie je
splniteľná (máme ukázať, že ako rezolventa sa tam ukáže prázdna klauzula). Ak
predpokladáme, že $S$ nie je splniteľná, potom podľa Herbrandovej vety (1.
variant), nie je splniteľná práve vtedy, keď je možné jej priradiť konečný
uzavretý sémantický strom. 

\par Môže as stať, že strom $T$ pozostáva jedine z koreňa -- odmieta jedinú
klauzulu a v tomto prípade veta platí. Teraz predpokladajme, že je konečný a má
viac ako 1 vrchol. V tomto prípade, tak má aspoň jeden akceptujúci vrchol. Potom
$i_v$ je čiastočná interpretácia končiaca v tom vrchole. Ďalej, každý
nasledovník je odmietajúci. Vrchol je akceptujúci, ak čiastočná interpretácia v
ňom existuje, a každý nasledujúci vrchol je odmietajúci.

\par Predpokladajme, že by tento strom nema akceptujúci vrchol. Potom každý
vrchol obsahuje nasledovníka, ktorý nie je odmietajúci. Týmto pádom by sme
vytvorili nekonečne dlhú vetvu, čo je spor (strom je konečný).

\par Ideme pracovať s akceptujúcim vrhcholom. Nech $v$ je akceptujúci vrchol
stromu $T$ a $v_1$, $v_2$ sú odmietajúci nasledovníci $v$. $I(v)$ (čiastočná
interpretácia končiaca vo vrchole $v$) vyzerá nasledovne:

\begin{eqnarray*}
I(v) &=& \{ m_1, m_2, \ldots, m_n \} \\
I(v_1) &=& \{ m_1, m_2, \ldots, m_n, m_{n+1} \}  \\
I(v_2) &=& \{ m_1, m_2, \ldots, m_n, \neg m_{n+1} \} 
\end{eqnarray*}

$C_1'$ a $C_2'$ sú dve základne inštancie klauzúl $C_1$ a $C_2$ -- $C_1'$ a
$C_2'$ neplatia v $I(v_1)$ a $I(v_2)$. $C_1'$ a $C_2'$ sa neodmietajú v $I(v)$.
$C_1'$ musí obsahovať $\neg m_{n+1}$ a $C_2'$ musí obsahovať $m_{n+1}$. $L_1' =
\neg m_{n+1}$ a $L_2' = m_{n+1}$. AAk vynecháme $L_1'$ a $L_2'$, dostaneme
rezolventu. $C'$ je rezolventa $C_1$ a $C_2$. $C' = (C_1' - L_1') \cup (C_2' =
L_2')$. $C'$ -- musí byť nepravdivá v $I(v)$. Podľa predchádzajúcej lemy musí
existovať rezolventa $C$ taká, že $C'$ je základná inštancia $C$.

\par Vezmime si $T''$ -- uzavretý sémantický strom, $C \cup \{C\}$. ... (niečo
ďalej?)



\paragraph{Príklad} Majme množinu formúl $F_1: (\forall x) (C(x) \implies (W(x)
\land R(x))$, $F_2: (\exists x)(C(x) \land Q(x))$ $G: (\exists x) (Q(x) \land
R(x))$. Ukážte, že $G$ je logickým dôsledkom $F_1$ a $F_2$.

\paragraph{Riešenie} Pre $F_1$, $F_2$ a $\neg G$ vytvoríme štandardné formy.
Dostávame nasledujúcich 5 klauzúl:
\begin{enumerate}
	\item $(\forall x) (C(x) \implies (W(x) \land R(x)) \iff (\forall
	x)(\neg C(x) \lor (W(x)\land R(x)) \iff (\forall x) ((\neg C(x) \lor
	W(x)) \land (\neg C(x) \lor R(x)))$
	\par (1) $\neg C(x) \lor W(x)$ -- $F_1$
	\par (2) $\neg C(x) \lor R(x)$ -- $F_1$
	\par $C(a)$ -- $F_2$
	\par $Q(a)$ -- $F_2$
\end{enumerate}
$$\neg G \iff \neg (\exists x)(Q(x)\land R(x)) \iff (\forall x) (\neg Q(x) \lor
\neg R(x)) $$. Štandardná formula pre túto formulu je:
\par (5) $\neg Q(x) \lor \neg R(x)$ -- $G$.

\par Rezolventy: 
\par (6) $R(a)$ -- rezolventa (2), (3)
\par (7) $\neg R(a)$ ($\sigma = \{a / x \})$ -- rezolventa (5), (4)
\par (8) $\square$ -- rezolventa (6), (7)

\par Záver: $G$ je logickýkm dôsledkom $F_1$ a $F_2$

\section{Stratégia vymazávania}
Na základe vety o úplnosti, majme nejakú množinu klauzúl $S$, postupne si z nej
vytváram rezolventy. Ak nie je splniteľná, po konečnom počte krokov dostávam
prázdnu klauzulu.

\paragraph{Príklad} Majme množinu klauzúl $S=\{P\lor Q, \neg P\lor Q, P \lor
\neg Q, \neg P \lor \neg Q\}$. Metódou rezolvent ukážte, že $S$ nie je
splniteľná.

\par nejakým spôsobom zostrojujeme postupnosť klauzúl, nakoniec niekde dostanem
prázdnu klauzulu. 
$S^0 = S$, $S^n = \{ \mbox{rezolventy} C_1, C_2 | C_1 \in S^0 \cup S^1,
\ldots \cup S^{n-1} \land C_2 \in S^{n-1}\}, n=1, 2, \ldots$. Takýmto spôsobom
by sme po $39$ krokoch dostali odpoveď. Niektoré klauzuly sa vyskytnú v
popísanom prístupe dvakrát. Môžu sa tam vyskytnúť tautológie.


\paragraph{Definícia} Klauzula $C$ je podklauzulou klauzuly $D$ (alebo pohlcuje
$D$) práve vtedy, keď existuje substitúcia $\sigma$ taká, že platí $C\sigma
\subseteq D$. $D$ nazývame nadklauzulou $C$.

\paragraph{Príklad} $C = P(x)$, $D = P(A) \lor Q(a)$, $\sigma = \{a/x\}$,
$C\sigma = P(a)$, $C\sigma \subseteq D$.

\paragraph{Poznámka} $D$ je identicky rovná $C$. (??)
Ak klauzula $D$ je inštancia $C$, tak $D$ je nadklauzula $C$. (??). 

\par
Možno to sformulovať takto: Majme $(S^0 \cup S^1 \cup \ldots \cup
S^{n-1}\}$, berieme z nej klauzulu $C_1$ a $C_2$. V prípade, že ako rezolventa
$C_1$ a $C_2$ nevznikne tautológia a nebude to ani nadklauzula niektorej z tých,
ktoré sú už vypísané, vypíšem ju. Takto sa zbavujem tautológií a nadklauzúl.

\paragraph{Príklad} $S=S^0$ a :
\begin{itemize}
	\item $P\lor Q$
	\item $\neg P \lor Q$
	\item $P\lor \neg Q$
	\item $\neg P \lor \neg Q$
\end{itemize}
$$
S_1:
\begin{array}{ll}
(5)& Q (1) (2)\\
(6)& P (1) (3)\\
(7)& \neg P (2) (4)\\
(8)& \square (6) (7)\\
\end{array}
$$

$$\theta = \{ a_1 / x_1, a_2 / x_2, \ldots, a_n / x_n \}$$. $x_1, x_2, \ldots
x_n$ sú premenné, ktoré sa vyskytujú v $D$. $a_1, a_2, \ldots a_n$ sú nové
konštanty, kotére sa nevyskytujú v $C$, $D$.

$D\Theta$ je základná... .

$$D = L_1 \lor L_2 \lor \ldots \lor L_m$$
$$D \theta = L_1\Theta \lor L_2 \Theta \lor \ldots \lor L_m \Theta$$
$$\neg D\Theta = \neg L_1 \Theta \land \neg L_2 \Theta \land \ldots \land \neg
L_m \Theta$$
Algoritums preveruje, či klauzula $C$ je podklauzulou $D$.

\paragraph{Algoritmus  pohltenia}
\begin{enumerate}
	\item Nech $W = \{ \neg L_1 \Theta, \neg L_2 \Theta, \ldots, \neg L_m
	\Theta \}$

	\item Kladieme $k=0$ a $\mathcal{U}^0 = \{ C \}$
	\item Ak $\mathcal{U}^k$ obsahuje $\square$, tak koniec $C$ je pod D. V
	opačnom prípade kladieme $\mathcal{U}^{k+1} = \{ \mbox{rezolventa} C_1 a
	C_2 | C_1 \in \mathcal{U}^{k} \land C_2 \in U\}$. 
	\item Ak $\mathcal{U}^{k+1}$ je $\emptyset$, tak koniec, $C$ nie je
	podklauzula $D$. V opačnom prípade kladieme $k=k+1$ a prejdeme ku kroku

\end{enumerate}

\paragraph{Poznámka} $\mathcal{U}^k, \mathcal{U}^{k+1}$, klauzuly z
$\mathcal{U}^{k}$ sú konečné. $\mathcal{U}^0, \mathcal{U}^1, \ldots \square$.

\paragraph{Dôkaz} Predpokladajme, že $C$ je podklauzula $D$. Na základe našej
definície existuje substitúcia $\sigma$, že $C\sigma \subseteq D$. Teda
$C(\sigma \circ \Theta) \subseteq D\Theta$. Literály z $C\sigma \circ \Theta$
môžeme vynechať pomocou jednotkových klauzúl z $W$. ... Algoritmus skončí svoju
činnosť.
\par
Obrátené tvrdenie: predpokladajme, že algoritmus zakončuje prácu na treťom
kroku. Odmietnutie môžeme znázorniť nasledujúcim obrázkom:

%pic1

$$C_0, N_1 ,\ldots B_r \in W$$
$$C(\sigma_0 \circ \sigma_1 \circ \sigma \circ \sigma_r) = \{ \neg B_0, \neg
B_1, \ldots \neg B_r\} \subseteq D\Theta$$
$$\lambda = \sigma_0 \circ \sigma_1 \circ \sigma_2 \ldots \circ \sigma_r \implies
C \lambda \subseteq D\Theta$$

$\sigma$, ktorá dostaneme z $\lambda$ tak, že v každom komponente $\lambda$
nahradíme  konštantu $a_i$ premennou $x_i$, $i=1, 2, 3, \ldots$. $C\sigma
\subseteq D$. $C$ je pod $D$.

\paragraph{Príklad} $C = \neg P(x) \lor Q(f(x), a)$. $D = \neg P(h(y)) \lor
Q(f(h(y)),a) \lor P(z)$. Zistite, či klauzula $C$ je podklauzulou $D$.

\par $y$ a $z$ sú premenné v $D$. $\Theta = \{ b/y, c/z\}$. Konštanty $b$, $c$
nevystupujú v $C$, $I$. najprv vypočítame $D\Theta \neg P(h(b)) \lor
Q(f(h(b)),a) \lor \neg P(c)$

$$\neg D \Theta = P(h(b)) \land \neg Q(f(h(b)),a) \ lor P(c)$$
$$W = \{P(h(b)), \neg Q(f(h(b)),a), P(c) \}$$
$$\mathcal{U}^0 = C = \neg P(x) \lor Q(f(x),a)\}$$
$\mathcal{U}^0$ neobsahuje $\square$, musíme vytvoriť $\mathcal{U}^1$. Urobíme
príslušnú substitúciu v množine $\mathcal{U}^0$. Dostávam nasledovné rezolventy:
$$\mathcal{U}^1 = \{ Q(f(h(b)),a), \neg P(h(b)), Q(f(b),a)\}$$. 
\par
$\mathcal{U}^1$
nie je prádzna a neobsahuje prádznu klauzulu -- musím vytvoriť $\mathcal{U}^2$.
V tomto sa už vyskytne prádzna klauzula, čo znamená, že $C$ pohlcuje klauzulu
$D$.


\paragraph{Príklad 1} $C=P(x,x)$ a $D=P(f(x),y) \lor P(y,f(x))$. Zistite, či $C$
je podklauzula $D$.

\paragraph{Riešenie} (1) $x$, $y$ sú premenné v $D$. $a$ a $b$ sú konštanty,
ktoré sa nevyskytujú $C$, $D$. $\Theta = \{ a/x, b/y\}$. $D\Theta = P(f(a),b), \lor P(b,
f(a))$.

$$\neg D\Theta = \neg P(f(a),b) \lor \neg O(b,f(a))$$
$$W = \{ \neg P(f(a),b), \neg P(b,f(a))\}$$
$$\mathcal{U}^0 = P(x,x)$$


\par (2) $\mathcal{U}^0$ neobsahuje $\square$, tak sa môže zistiť
$\mathcal{U}^1$
\par (3) $\mathcal{U}^1 = \emptyset$. Záver: $C$ nie je podklauzula $D$.


\paragraph{Príklad 2} Majme formuly:

\begin{enumerate}
	\item $P\implies S$
	\item $S \implies U$
	\item $P$
	\item $U$
\end{enumerate}

Dokážte, že formula 4 vyplýva z formúl 1, 2 a 3. 

\paragraph{Riešenie} Prepíšeme si formuly do správneho tvaru, aby sme mohli
použiť pravidlo rezolventy:
\begin{enumerate}
	\item $\neg P \lor S$
	\item $\neg S\lor U$
	\item $P$
	\item $U$
\end{enumerate}
Snažíme sa nájsť negáciu -- chceme ukázať, že 
\begin{enumerate}
	\item $\neg P \lor S$
	\item $\neg S\lor U$
	\item $P$
	\item $\neg U$
\end{enumerate}

nie je splniteľná. Zoberiem si rezolventu 1 a 3, dostávam $S$ (5). Keď zoberiem 
2 a 4, dostávam $\neg S$ (6). Zoberiem 5 a 6, dostávam $\square$ (7).

\paragraph{Príklad 3} Predpoklad: Študenti sú občania. Záver: Hlasy študentov sú
hlasy občanov.

\paragraph{Riešenie} 
\begin{itemize}
	\item $S(x)$ označuje \uv{$x$ je študent}.
	\item $C(x)$ označuje \uv{$x$ je občan}.
	\item $V(x,y)$ označuje \uv{$x$ je hlas $y$}.
\end{itemize}
Predpoklad: $(\forall y) (S(y)\implies C(y))$. Študenti sú občania.
Záver: $(\forall x) ((\exists y) (S(y) \land V(x,y)) \implies (\exists z)(C(z)
\implies V(x,z)))$. Hlasy študentov sú hlasy občanov.

Aká bude štandardná forma pre vyjadrenie predpokladu?
\begin{enumerate}
	\item $\neg S(y) \lor C(y)$
		\par $\neg ((\forall x) ((\exists y)(S(y)\land V(x,y)) \implies
		(\exists x)(C(z) \land V(x,z)))) \iff 
		\neg ((\forall x)(\forall y)(\neg S(y)\lor \neg V(x,y)) \lor
		(\exists z)(C(z)\land V(x,z))) \iff 
		\neg ((\forall x)(\forall y)(\exists x)(\neg S(y) \lor
		\neg U(x,y)\lor (C(z) \land V(x,z)))) \iff
		(\exists x)(\exists y)(\forall z) (S(y) \land V(x,y)) \land
		(\neg C(z) \lor \neg V(x,z))$
		Teraz potrebujeme Skolemov normálny tvar:
		$(\forall z)(S(b) \land U(a,b)) \land(\neg C(z) \lor \neg
		V(a,z)))$
		\par Pre negáciu záver dostávane nasledujúce klauzuly:
	\item $S(b)$
	\item $V(a,b)$
	\item $\neg C(z) \lor \neg V(a,z)$.

	\item $C(b)$ z (1) a (2) (miesto $y$ dosadíme $b$)
	\item $\neg V(a,b)$ zo (4) a (5)
	\item $\square$ z (3)  a (6)
\end{enumerate}

\par Predpokladajme, že $b$ je študent, $a$ je hlas študenta $b$ a nie je hlas
žiadneho občana. Pretože $b$ je študent, $b$ je občan. Okrem toho $a$ nemôže byž
hlas $b$, pretože $b$ je občan a to nie je možné.

\section{Rozširovanie teórie}
\paragraph{Veta} Nech $A$ je formula teórie $T$ s jazykom $L$, nech všetky voľné
premenné formuly $A$ sú $x_1, x_2, \ldots x_n$, ďalej nech $L'$ je rozšírenie
$L$ pridaním nového predikátového symbolu $P$ a nech $T'$ vznikne z $T$ pridaním
axiomy $P(x_1, \ldots x_n) \iff A$ (označme ju *). 
\par Potom teória $T'$ s jazykom $L'$ je konzervatívne rozšírenie teórie $T$ s
jazykom $L$. Ďalej, pre každú formulu $B$ jazyka $L'$ existuje formula $B^*$ z
jazyka $L$, že platí $T \models  B \iff B^*$. 

\par Najprv budeme konštruovať $B^*$. Dôkaz: Nech $B$ je formula na jazyku $L'$,
nech $A'$ je variant  formuly $A$ (definujúcej predikát $P$), že žiadna
premennáz formuly $B$ nie je viazaná vo formule $A'$.

\par Nech $B^*$ vznikne z $B$ tak, že každú podformulu $P(a_1, \ldots a_n)$
nahradíme podformulou $A'(a_1, \ldots a_n)$ a podľa vety o variantoch nám platí,
že:
$$ T' \models P(a_1, \ldots, a_n) \iff A'(a_1, \ldots a_n)$$
$$ T' \models B \iff B^*$$.

\par Ukážeme, že $T'$ v jazyku $L'$ je konzervatívne rozšírenie teórie $T$ v
$L$ (teda pre ľubovoľnú formulu na pôvodnom jazyku, ktorá je dokázateľná v
teórii $L'$, je dokázateľná aj v pôvodnej teórii --  nevzniknú žiadne nové
teorémy). Nech $C$ je ľubovoľné formula na jazyku $L'$, ďalej nech platí $T'
\models C$. Nám stačí dokázať, že $T \models C^*$. Ak je potom formula $C$ na
jazyku $L$, potom $C^*$ je $C$.
\par Uvažujme $C_1, C_2, \ldots C_n$ je odvodenie formuly $C$ v teórii $T'$. Nám
stačí ukázať, že $C_i^*$ sú dokázateľné (odvoditeľné) v $T$.

Pri odvodení $C$ sa môže stať nasledovné:
\begin{enumerate}
	\item $C_i$ je axioma predikátovej logiky $L$, potom $C_i^*$ je axioma
	rovnakého druhu
	\item $C_i$ je axioma z $T$, potom $C_i^*$ je $C_i$ a teda je
	dokázateľná z $T$.
	\item Môže sa stať, že $C_i$ je axioma (x) $P(x_1,  \ldots, x_n) \iff
	A(x_1, \ldots, x_n)$. \par $C_i^*: A_i \iff A_i$, je dokázateľná pomocou
	vety o variantoch.
	\item Ak $C_i$ je odvodená z $C_j$ a $C_k$, použijeme pravidlo modus
	ponens. $C_i^*$ je odvodené z $C_j^*$ a $C_k^*$ tým istým pravidlom.
	\item $C_i$ je odvodené z $C_k$ pravidlom zovšeobecnenia ($k<i$). Potom
	$C_i^*$ je odvodené z $C_k^*$ pomocou toho istého pravidla.
\end{enumerate}



Ak máme teóriu $T$ a každá jej $A \in T$ je otvorená (každá premenná je voľná),
hovoríme o otvorenej teórii.

\par Majme danú teóriu $T$ s jazykom $L$. Máme k nej zostrojiť teóriu $T_H$
(henkinova teória) s jazykom $L(C)$, kde $C$ je zjednotenie konštánt všetkých
rádov. Dôležité pre nás budú konštanty $c_A, c_{\neg A}$.

\begin{itemize}
	\item $A_{c_{\neg A}} \implies (\forall x) A$ (*)
	\item $(\exists ) A \implies A(c_A)$ (**)
\end{itemize}

$$A \iff \neg \neg A$$
$$c_{\neg A} \iff \neg A(x)$$
$$\neg A(c_{\neg\neg A}) \implies (\forall x)\neg A$$
$$(\exists x) A \implies A(c_{\neg\neg A})$$
(Použili sme prenexnú operáciu).

\par Ak $c_{\neg A}$ je konštanta z axiomy s (*), tak hovoríme, že prislúcha
(patrí) formule A. Keď zoberieme formulu $A$.

$$(\forall x) A \implies A(t)$$ -- $t$ je term bez premenných v $L(C)$

$$A_{[c_{\neg A}]} \implies (\forall x) A \mbox{.......} c_{\neg A}$$

$\delta(T)$ množina na jazyku $L(C)$ prislchajúu konštante $C$ má ... axiomy
identity a ... a formúl z $T$.


\paragraph{Lema 1} Nech $A$ je formula na jazyku $L$ a ďalej, nech $A'$ je
uzavretá inštancia formuly $A$ na jazyku $L(C)$. Ak platí, že formula $A$ je
dokázateľná v pôvodnej teórii, potom $A'$ je tautologickým dôsledkom konečne
mnoho formúl z $\delta(T)$

\paragraph{Dôkaz} Nech $A_1, A_2, \ldots, A_n$ je dôkaz formuly $A$ v teórii
$T$. Indukciou podľa dĺžky dôkazu ukážeme, že ľubovoľná uzavretá formula v
jazyku $L(C)$ inštancia $A_j'$ formuly $A_j$ je tautologickým dôsledkom konečne
veľa formúl z $\delta(T)$.

\par Môžu nastať tieto prípady:
\begin{itemize}
	\item $A_j$ je axioma výrokovej logiky. Potom $A_j'$ je opäť tautológia.
	$A$ je tautologickým dôsledkom prázdnej množiny predpokladov.
	\item $A_j$ je axioma tvaru $(\forall x) B \implies B_x[t]$. Axioma
	špecifikácie. $A_j: (\forall x) B' \implies B_x'[t]]$ -- leží v
	$\delta T$. $t'$ je term bez premenných a je v $L(C)$. 
	\item $A_j$ je axioma $(\forall x) (C \implies D) \implies (C \implies
	(\forall x) D)$, $x$ nie je voľná v $C$.
	\par Uvažujme formulu $(\forall x) (C \implies D) \implies (C \implies
	(\forall x) D_x[c_{\neg D}])$. Táto formula patrí to $\delta(T)$, lebo
	je to inštancia axiomy špecifikácie.
	\par $ D_x[c_{\neg D}] \implies (\forall x) D$ -- opäť patrí do
	$\delta(T)$ (opäť (*))
	\par Tvrdíme, že formula $A_j$ je tautologický dôsledok horeuvedených 
	formúl. Budeme uvažovať takúto formulu:
	t -- $(A_1 \implies (B_1 \implies D_1)) \implies ((D_1 \implies C_1)
	\implies (A_1 \implies (B_1 \implies C_1)))$ Tvrdíme, že táto formula
	je teoréma.

	\par Dôkaz formuly: Za predpoklady si zoberiem formuly $A_1, (A_1
	\implies (B_1 \implies D_1))$, $B_1$, $(D_1 \implies C_1)$. Z týchto
	predpokladov dokážem odvodiť $C_1$.

	\par $A_1$ bude formula $(\forall x)(C \implies D)$, $B_1: C$, $C_1:
	(\forall x) D$, $D_1: D_x[c_{\neg D}]$

	\item $A_j$ je axioma z $T$, potom $A_j' \in \delta(T)$.
	\item $A_j$ je axioma identity alebo axioma rovnosti, potom $A_j'$ patrí
	do $\delta(T)$.

	\item $A_j$ dostaneme aplikáciou pravidla modu
	ponens z formúl $A_k$ a $A_l$, pričom $k, l < i$. $A_j'$ dostávame z
	$A_k'$ a $A_l'$ pomocou pravidla modus ponens.
	\item $A_j$ je odvodená z $A_k$ pravidlom zovšeobecnenia, pričom
	predpoklad je $k<j$. Teda $A_j: (\forall x) C$, teda $A_k: C$, x --
	premenná. Inštancia $A_j': (\forall x) C'$. Uvažujme inštanciu $A_k'$
	formuly $A_k$ v takomto tvare: $A_k:  C_k'[c_{\neg C'}]$. Podľa
	indukčného predpokladu je táto formula tautologickým dôsledkom konečne
	veľa formúl z $\delta (T)$. Uvažujme formulu $$C_x'[c_{\neg C'}]
	\implies (\forall x) C' \qquad \mbox{(***)}$$
	$A_j'$ je tautologickým dôsledkom formuly z $\delta(T)$. Dokázali sme
	formulu $A'$ tak, že sme nepoužili pravidlo zovšeobecnenia.
\end{itemize} 

\paragraph{Definícia} Hovoríme, že formula $A$ je \emph{kvázitautológia}, ak je
tautologickým dôsledkom inštancií axiom identity a rovnosti.

\paragraph{Veta (Hilbert-Ackermann)} Otvorená teória $T$ v jazyku $L$ (s
rovnosťou) je sporná práve vtedy, keď existuje (kvázi-)tautológia, ktorá je
disjunkciou negácií inštancií axiom z $T$.

\paragraph{Dôkaz} Najprv ľahšia implikácia: Predpokladajme, že sú splnené
podmienky vety, tak potom $T$ je sporná. Zoberme si nejakú formulu $A$. Berieme
$\neg A_1 \lor \neg A_2 \lor \ldots \lor \neg A_n$, kde $A_i$ sú inštancie axiom
z $T$. Pomocou de Morganovho zákona dostaneme $A_1 \land A_2 \land \ldots \land
A_n$. Ak platí $T \models A_i$, potom $T \models A_1 \land \ldots \land A_m$, a
teda $T$ je sporná, lebo je z nej dokázateľná ľubovoľná formula $B$, pretože
$\models A \implies (\neg A \implies B)$ a $T \models B$.

\par Naopak: Predpokladajme, že otvorená teória $T$ je sporná. $x \neq x$, $T
\models x \neq x$. Z toho vyplýva, že ak si zoberieme ľubovoľnú konštantu $r \in
C$, tak potom dostávame inštanciu formuly: $r \neq r$. To je uzavretá inštancia
teorémy vety z $T$, čiže podľa  Lemy 1 eistuje $A_1, A_2, \ldots A_k$  z $\delta
(T)$ také, že $r \neq r$ je tautologickým dôsledkom $A_1$ až $A_k$, teda z
predpokladov inštancií je dokázateľné $A_1, A_2, \ldots, A_k \models r \neq r$.
Platí $A_1 \implies A_2 \implies \ldots \implies A_k \implies r \neq r$ --
tautológia. $p \implies q \iff \neg p \lor q$. $\neg A_1 \lor \neg A_2 \lor
\ldots .. \lor \neg A_k \lor \neg (r=r)$ je tautológia. $A_i$ sú inštancie z
$\delta (T)$. Posledná formula je tautológia, ktorá je disjunkciou negácií
inštancií z $\delta (T)$.

\paragraph{Definícia} Postupnosť formúl $A_1, A_2, \ldots A_n$ nazveme
\emph{špeciálnou}, ak $\neg A_1\lor \neg A_2 \lor \ldots \lor \neg A_n$ je
tautológia. (Vieme) Ak $T$ je sporná teória, potom existuje špeciálna postupnosť
z $\delta (T)$



(pokračovanie)
\paragraph{Definícia} \emph{Stupeň konštanty} $c_{\neg A}$. $(\forall x) A$
uzavretá formula na jazyku $L(C)$. Hovoríme, že konštanta $c_{\neg A}$ spojená
axiomou $$A_{c_{\neg A}} \implies (\forall x) A$$ s formulou $(\forall x) A$ je
stupňa $n$, ak formula $(\forall x) A$ obsahuje $n$ výskytov kvantifikátorov
$\forall$ alebo $\exists$. 

\par Množina formúl $\delta_n(T)$ vznikne z $\delta(T)$ vynechaním všetkých
formúl stupňa $>n$ pre konštantu $c_{\neg A} \in C$. Treba si uvedomiť, že
stupeň konštanty $c_{\neg A}$ je vždy aspoň $1$. $\delta_0(T)$ -- uzavreté
inštancie, axiomy z množiny $T$ a axiom identity a rovnosti.

\par (Dokázali sme) Ak teória $T$ je sporná, existuje špeciálna postupnosť,
ktorá patrí do $\delta(T)$. Nech $n$ je najmenšie také $n$, že špeciálna
postupnosť z $\delta(T)$ je obsiahnutá v $\delta_n(T)$:

$n=0$. Musíme nájsť špeciálnu postupnosť prislúchajúcu do $T$.
Predpokladáme, že pre $n=0$ máme postupnosť $A_1, A_2, \ldots A_k$ --
špeciálna a patrí do $\delta_O(T)$. $c_A, c_{\neg A}$ nahradíme pomocou
premenných $A_1', A_2', \ldots A_k'$ (opäť špeciálna postupnosť). $A_1',
A_2', \ldots, A_n'$ sú všetky formuly, ktoré patria do $T$. $A_1', A_2',
\ldots A_k'$ budú inštancie axiom identity a rovnosti. $\neg A_1' \lor
\neg A_2' \lor \ldots \lor \neg A_n' \leftarrow A_{n+1} \leftarrow
\ldots \leftarrow A_k'$ -- kvázitautológia (vyplýva z toho, že táto
postupnosť je špeciálna). $A_1', A_2', \ldots A_n' \in T$.

\paragraph{Lema} Ak $n>0$ a existuje špeciálna postupnosť z $\delta_n(T)$, tak
potom existuje špeciálna postupnosť aj z $\delta_{n-1}(T)$.

\paragraph{Dôkaz} (pokračovanie dôkazu) Hilbert-Ackermanovej vety.
$\delta_0(T)$. Vieme vytvoriť postupnosť $B_1, B_2, \ldots, B_q \in T$.

\paragraph{Poznámka} Elementárna aritmetika je otvorená teória s konečným počtom
axiom. Hilbertova arigmetika, Presburgerova (??) aritmetika (aritmetika so
symbolmi $0, S, +$. Peanova aritmetika (1931).

\paragraph{Definícia} 
\begin{enumerate}
\item
Hovoríme, že formula $A$ je existenčná, ak $A$ je v
prenexnom tvare a všetky kvantifikátory v prefixe sú existenčné .
\item Hovoríme, že formula $A$ je univerzálna, ak $A$ je v prenexnom tvare a
všetky kvantifikátory sú univerzálne.

\end{enumerate}

\paragraph{Lema 3} Uzavretá existenčná formula $A$ je dokázateľná v predikátovej
logike (s rovnosťou) práve vtedy, keď istá disjunkcia otvoreného jadra formuly
$A$ je kvázitautológia.

\paragraph{Dôkaz} $A: (\exists x_1)(\exists x_2) \ldots (\exists x_n)B$
--formula v prenexnom tvare, $B$ je otvorené jadro. 

\begin{enumerate}
	\item Formulá $A$ je dokázateľná $\iff$ teória s jedinou špeciálnou
	axiomou $\neg A$ je sporná.
	\item Ak použijeme prenexné operácie, dostávame $\models \neg A \iff
	(\forall x_1)\ldots (\forall x_n) \neg B$. $\neg A \iff (\forall x_1)
	\ldots  (_n) \neg B$.

	\item Formula $A$ je dokázateľná $\iff$ keď teória $T$ s jedinou
	špeciálnou axiomou $\neg B$ je sporná. $\neg B$ je otvorená. Keď
	použijeme Hilbert-Ackermannovu vetu, dostávame $\neg B_1, \neg B_2
	\ldots, \neg B_m$ -- nad $\neg B$. Platí $\neg \neg B_1 \lor \neg \neg
	B_2 \lor \ldots \lor \neg \neg B_m$, čo je to isté ako $B_1 \lor B_2
	\lor \ldots \lor B_m$, čo je to isté, ako Hilbert-Ackermannova veta.
\end{enumerate}


\par Teóriu $T$ rozširujeme do $T_H$ s jazykom $L(C)$ a tú zase do $T_R$ s
$L(C)$. $A_x[C_{\neg A}] \implies (\forall x) A$ -- toto môžeme aj obrátiť:
$$(\forall x) A \implies A_x[c_{\neg A}]$$ -- vďaka axiome špecifikácie.
\par Ak teraz vezmeme formulu $B$ takú, že $A\iff B$, potom platí $B_x[c_{\neg
B}] \implies (\forall x) B$. 
$$
\begin{array}{lll}
T_H &\models& A_x[c_{\neg A}] \iff B_x[c_{\neg B}] \\
T_H &\models& A_x [c_{\neg A}] \iff (\forall x) A \\
T_H &\models& A \iff B \\
T_H &\models& (\forall x) A \iff (\forall x) B \\
T_H &\models& B_x[c_{\neg B}] \iff (\forall x) B \\
\end{array}
$$


$$(\forall x)(A\iff B) \implies c_{\neg A} = c_{\neg B}$$

\paragraph{Lema 4} Pre ľubovoľnú teóriu $T$ je $T_R$ konzervatívne rozšírenie.

\par Majme formulu $A$, ktorá je uzavretá a napísaná v prenexnom tvare.
Indukciou podľa počtu všeobecných kvantifikátorov budeme formulu transformovať
do tzv. Herbrandovho variantu. Ak formula $A$ má tvar:
$$ A: (\exists x_1) \ldots (\exists x_n) (\forall y) B, \qquad n \geq 0$$
$$ A^*: (\exists x_1) \ldots (\exists x_n) B_y[f(x\ldots x_n)]$$, $A^*$ je $A_H$

$A^{**}, L \cup \{ f,g,h, \ldots \}$.

\paragraph{Veta} (Herbrandova) Uzavretá formula v prenenom tvare je dokázateľnáv
predikátovej logike (s rovnosťou) práve vtedy, keď istá disjunkcia inštancií
otvoreného jadra je kvázitautológia.

\paragraph{Dôkaz} $A_H$ -- je existenčná formula. Je dokázateľná práve vtedy,
keď vezmeme do úvahy lemu 3 a platí formulácia našej vety. $L'$ .... ??? . Stačí
ukázať, že:
$$\models_L A \iff \models_{L'} A_H$$.
Implikácia zľava doprava je ľahšia -- vieme ukázať, že $\models_{L'} A \implies
\models_{L'} A^*$. Premennú $y$ sme nahradili funkciu $n$ premenných. Všimnime
si:

$$ \models_{L'} (\forall y) B \implies B_y[f(x_1\ldots x_n)]$$; axioma
špecifikácie. $(\exists x_1)(\exists x_2), \ldots (\exists x_n)$. Potom platí aj
$\models A \implies B$, $\models (\exists x) A \implies (\exists x) B$



\paragraph{Veta} (o zavedení funkčného symbolu) Nech formula $(\exists y) A$ je
dokázateľná v teórii $T$ s jazykom $L$. Nech $x_1, x_2, \ldots, x_n$ sú všetky
voľné premenné, ktoré sa vyskytujú vo formuli $(\exists y) A$. Nech $T'$ vznikne
z $T$ pridaním nového $n$-árneho funkčného symbolu $f$ a pridaním axiomy
$A_y[f(x_1, \ldots, x_n)]$. Potom $T'$ je konzervatívne rozšírenie teórie $T$.

\paragraph{Dôkaz} Nech $B$ je uzavretá formula teórie $T$ (s jazykom $L$).
Predpokladáme, že $B$ je teorémou (vetou teórie) $T'$. $T' \models B$. Máme
ukázať, že aj v $T \models B$. Predpokladajme, že formula $B$ má dôkaz v $T'$ a
že v tom dôkaze vystupujú $a_1, a_2, \ldots, a_n$ -- špeciálne axiomy teórie
$T$, prípadne axioma $(*)$. Na základe tohto predpokladu je v predikátovej
logike dokázateľná nasledovnaá formula.

$$(1) \models (\forall x_1) (\forall x_2) \ldots (\forall _n) a_y [f(x_1 \ldots
x_n)] \implies B_1 \implies B_2 \implies \ldots \implies B_n \implies B$$

$B_1, \ldots, B_n$ sú ...ny funkcií $A_1, A_2, \ldots, A_n$. 

$$ A_y[f(x_1, \ldots x_n)], A_1, \ldots, A_n \models B$$

Označme si $C$ -- prenexný tvar nasledovnej formuly:
$$A \implies B_1 \implies \ldots \implies B_n \implies B$$
Potom:

$$ (2) (\exists x_1) (\exists x_2) \ldots (\exists x_n) C_y[f(x_1, \ldots,
_n)]$$ je prenexný tvar $(1)$

Uvažujme $D: (\exists x_1) \ldots (\exists x_n) (\forall y) C$. Potom z tohto
vyplýva $D$ je formula jazyka $L$ a neobsahuje novo zavedený symbol $f$. Ak
konštruujeme herbrandovský variant tej podformuly ($D_H$) (meníme veľké
kvantifikátory) a na prvom kroku dostávame $D^*: (\exists x_1) \ldots (\exists
x_n) C_y [f(x_1, \ldots x_n)]$ -- prenexný tvar $(1)$. Z Herbrandovej vety
dostávame, že $D^*$ je dokázateľné, z tadiaľ $(D^*)_H$ je dokázateľné a z toho 
$(4) T \models D$

\par Prenexnými operáciami dostávame  $$(5) (\exists x_1) (\exists x_2) \ldots
(\exists x_n) (\forall y) (A \implies B_1 \implies \ldots \implies B_n \implies
B)$$

$$(6) T \models (\forall x_1 ) (\forall x_2) \ldots (\forall x_n) (\exists y)
A \implies B_1 \implies B_2 \implies \ldots \implies B_n \implies B$$

\par Každý z týchto predpokladov je dokázateľný v $T$, dostávame $T \models B$ a
teda $T'$ je konzervatívne rozšírenie $B$.

Mám formulu $A$, ktorá je uzavretá, idem priradiť skolemov tvar (skolemov
variant).
\begin{enumerate}
	\item Ak formula $A$ je univerzálna, potom $A_S$ je formula $A$
	\item Ak $A$ je tvaru $(\forall x_1) \ldots (\forall x_n) (\exists y) B,
	n\geq 0$, $f$ je funkčný symbol, kladieme $A_S: (\forall x_1) \ldots
	(\forall x_n) B_y[f(x_1,\ldots, x_n)]$, $(A')C ...$
\end{enumerate}

(1*) $A' \models A$ a ... $A_S \models A$
$ \models B_y[f(x_1, \ldots, x_n)] \implies (\exists y) B$ -- duálny tvar
axiomy špecifikácie. Teraz môžeme sformulovať Skolemovu vetu:

\paragraph{Veta} (Skolemova) V ľubovoľnej teórii $T$ môžeme zostrojiť otvorenú
teóriu $T'$, ktorá je konzervatívnym rozšírením teórie $T$.
\paragraph{Dôkaz} Mám teóriu $T$ s jazykom $L$. Zostrojím si teóriu $T_1$ s tým
istým jazykom a tá teória $T_1$ sú uzávery prenexných tvarov axiom z $T$. Z vety
o uzávere a vety o prenexnom tvare platí, že $T_1$ je konzervatívne rozšírenie
$T$ a naopak, $T$ je konzervatívne rozšírenie $T_1$. $T \equiv T_1$.

\par Teória $T_2$ vznikne z teórie $T_1$ tak, že ľubovoľnej formuli $A \in T_1$
zostrojíme skolemov variant $A_S$. V $T_2$ je konzervatívne rozšírenie $T_1$
podľa vety o zavedení funkčného symbolu. 

\par Ideme vytvárať teóriu $T_3$ -- vznikne z $T_2$ vynechaním všetkých axiom z
$T_1$. Podľa $(1^*)$ dostávame, že $T_2 \equiv T_3$.

\par Ďalej vytvárame teóriu $T_4$ -- pozostáva z otvorených jadier $T_3$, teda
$T_3 \equiv T_4$. Keď to zhrnieme, dostávame: $T \equiv T_1$, $T_2$ je
konzervatívne rozšírenie $T_1$. $T_2 \equiv T_3 \equiv T_4$, a z toho $T_4$ je
konzervatívne rozšírenie $T$.

\paragraph{Veta} (o zavedení funkcie pomocou definície) Majme teóriu $T$ s
jazykom $L$ a nech $x_1, \ldots, x_n, y$ sú navzájom rôzne premenné, ktoré sa
vyskytujú voľne vo formule $D$. Nech platí:
\begin{enumerate}
	\item $T\models  (\exists y) D$
	\item $T\models D \implies (D_y[y] \implies y = t)$
\end{enumerate}
Nech $L'$ vznikne z $L$ pridaním nového $n$-árneho funkčného symbolu $f$ a $T'$
z $T$ pridaním axiomy $$(3) y=f(x_1, \ldots x_n) \iff D$$ (definícia axiomy). Potom
$T'$ je konzervatívne rozšírenie $T$ a ku každej formule $A$ na jazyku $L'$
existuje formula $A^*$ na jazyku $L$ taká, že platí $$(4) T' \models A \iff A^*$$

\paragraph{Dôkaz} Najprv ukážeme, ako ku formuli $A$ priradiť formulu $A^*$, a
potom ukážeme, že $T'$ je konzervatívne rozšírenie $T$.

\par Zostrojíme $A^*$ z $A$ tak, aby platilo $(4)$. Formulu $A$ máme na jazyku
$L'$, problematický je symbol $f$ -- vyskytuje sa v atomických podformuliach.
Nech funkčný symbol $f$ sa vyskytuje vo formuli $A$ a nech je to ten
najvnútornejší výskyt. Je to nejaký $f(t_1, \ldots t_n)$, pričom $t_1, \ldots,
t_n$ už neobsahujú $f$. 

$$A: B_z[f(t_1, \ldots, t_b)]$$, a naše $z$ sa nevyskytuje vo formuli $A$, a ani
v definujúcej formuli $D$. Položme $A^*$  tvare:
$$(5) (\exists z) D'_{x_1, \ldots, x_n, y}[t_1, \ldots, t_n, z] \land B^*)$$,
pričom $D'$ je variant $D$, v ktorej nie je viazaná žiadna premenná, ktorá sa
vyskytuje vo formuli $A$. Z vety o variantoch a definujúcej axiomy $(3)$
dostávame, že na jazyku $L'$ je dokázané $ (6) \models_{L'} z = f(t_1, \ldots, t_n) \iff D'[t_1,
\ldots t_n, z]$. Z vety o ekvivalencii dostávame, že platí $$(7) T' \models
(\exists z) z = f(t_1, \ldots, t_n) \land B^*) \iff f^*$$. Z indukčného
predpokladu dostávame $$(8) T' \models B \iff B'$$

\par $(1) T' \models (\exists z) (z = f (t_1, \ldots, t_n) \land B) \iff A^*$ z
vedy o .... ($\models (\exists x) (x=t \land A) \iff A_x[t]$)
$$T' \models B_z[t_1, \ldots, t_n] \iff A^*$$, na ľavej strane je formula $A:
B_z[t_1, \ldots t_n]$.

\par Použijeme vetu o zavedení funkčného symbolu. Uvažujme $S$ -- teóriu z
jazyka $L'$, ktorá vznikne z $T$ pridaním axiomy:
$$(10) D_y[f(x_1, \ldots, x_n)]$$.

Ukážeme, že $T'$ a $S$ sú ekvivalentné a platí $T' \equiv S$. Potrebujeme
ukázať, že $(10)$ je teorémou $T'$ a $(3) y=f(x_1, \ldots, x_n) \iff D$ je
teorémou (vetou) $S$. V $T' \models f(x_1, \ldots x_n) = f(x_1, \ldots, _n)
\implies D_y[f(x_1, \ldots _n)]$, formula je inštancia axiomy $(3)$. Ak
použijeme inštanciu na axiomu identity $(3)$ $T' \models D_y[f(_1, \ldots,
x_n)]$. Potrebujeme ukázať, že $3$ je dokázateľná v $S$. Vyjdeme z tvaru $(11)
\models_{L'} y = f(x_1, \ldots x_n) \implies (D \iff D_y[f(x_1, \ldots, x_n)])$.

$$(12) \models_{L'} D_y[f'(x_1, \ldots, x_n)] \implies (y = f(x_1, \ldots _n
\implies D) $$

$$(13) S \models (y = f(x_1, \ldots, x_n) \implies D) \mbox{ z $(10)$}$$

$$(14) \models_{L'} D_y[f(x_1, \ldots, x_n) ] \implies (D \implies y = f(x_1,
\ldots x_n)] [T \models D \implies [D_y[t] \implies y = t) (2)$$


$$S \models D \implies y = f(x_1, \ldots x_n)$$

Skúška:
\begin{itemize}
	\item Dokázať $4$ vety
	\item je daná formula, nájsť prenexný tvar, skolemov tvar, metódou
	rezolvent ukázať, či platí alebo neplatí, napíšte teóriu s rovnosťou
	\item 3 alebo 4 definície
\end{itemize}


\end{document}


\end{document}
